% !TEX root = github.tex
% !TEX program = pdflatex
\documentclass[12pt,reqno]{amsart}

\usepackage[utf8]{inputenc}
\usepackage{amsmath,amssymb,amsthm,mathtools,mathrsfs,esint}
\usepackage{mathalpha}
\usepackage{enumitem}
%\usepackage[allowunsupported,grade=19]{cmgraded}
%\usepackage{libertinus}
\usepackage{microtype}
\usepackage[a4paper,top=25mm,bottom=26mm,left=29mm,right=29mm,headsep=7mm]{geometry}
\usepackage[dvipsnames]{xcolor}
\usepackage[
  breaklinks=true,
  pagebackref=true,
  hypertexnames=false
]{hyperref}
\hypersetup{
  colorlinks=true,
  bookmarksdepth=2,
  pdfpagemode=UseNone,
  citecolor=ForestGreen,
  linkcolor=MidnightBlue,
  urlcolor=RoyalBlue,
  pdfstartview=FitH,
  pdftitle={Monotonicity of the Speed of Biased Random Walks on Galton--Watson Trees},
  pdfauthor={GPT-6 Astra and Heng Ma},
  pdfsubject={Biased random walk, Galton--Watson tree, speed monotonicity}
}
\usepackage[nameinlink,noabbrev]{cleveref}

% Page geometry and paragraph spacing from the template.
\linespread{1.05}
\normalsize
\linespacing=\baselineskip
\setlength{\parindent}{1em}
\setlength{\parskip}{3pt}

\makeatletter

% First-page title, author block, abstract, and separate AI usage statement.
\def\@settitle{%
  \begin{center}%
    \uppercasenonmath\@title
    \Large\@title
  \end{center}}

\def\@setauthors{%
  \vskip21pt plus6pt minus3pt
  \begin{center}%
    {\tiny BY\par}\medskip
    {\normalsize\scshape\authors\par}\medskip
    \def\address##1##2{\par ##2\par}%
    \def\email##1##2{\par e-mail:~##2\par}%
    {\small\itshape\addresses\par}%
  \end{center}}

\def\enddoc@text{}

\renewenvironment{abstract}{%
  \global\setbox\abstractbox=\vtop\bgroup
  \vskip10pt
  \normalfont\tiny\centering\MakeUppercase{\abstractname}\par
  \Small
  \list{}{%
    \leftmargin3pc
    \rightmargin3pc
    \labelwidth0pt
    \listparindent\normalparindent
    \itemindent0pt
    \parsep0pt plus1pt
    \let\fullwidthdisplay\relax
  }%
  \item[]
}{%
  \vspace*{-10pt}
  \endlist\egroup
}

\newbox\aiusagebox
\newenvironment{aiusage}{%
  \global\setbox\aiusagebox=\vtop\bgroup
  \normalfont\Small
  \list{}{%
    \leftmargin3pc
    \rightmargin3pc
    \topsep0pt
    \partopsep0pt
    \parsep0pt
    \itemsep0pt
    \listparindent0pt
    \itemindent0pt
  }%
  \item[]%
  \nobreak\vspace{4pt}%
  {\normalfont\tiny\centering\MakeUppercase{AI usage}\par}%
  \nobreak\vspace{2pt}%
  \noindent\ignorespaces
}{%
  \par\endlist\egroup
}
\let\original@setabstracta\@setabstracta
\def\@setabstracta{%
  \original@setabstracta
  \ifvoid\aiusagebox\else
    \vskip10pt
    \box\aiusagebox
    \prevdepth\z@
  \fi
}

% Section headings.
\let\@secnumfont\@empty
\def\section{\@startsection{section}{1}{0pt}%
  {20pt plus4pt minus4pt}
  {.5\linespacing}
  {\normalfont\bfseries}}
\def\subsection{\@startsection{subsection}{2}{0pt}%
  {.5\linespacing plus.7\linespacing}
  {-.5em}
  {\normalfont\scshape}}
\def\subsubsection{\@startsection{subsubsection}{3}{0pt}%
  {.5\linespacing plus.7\linespacing}
  {-.5em}
  {\normalfont\itshape}}
\let\@afterindenttrue\@afterindentfalse
\@afterindentfalse

% Theorem and proof styles.
\newtheoremstyle{plain}
  {.5\linespacing}{.5\linespacing}
  {\itshape}{1em}{\scshape}{.}{.5em}
  {\thmname{#1}\thmnumber{ #2}\thmnote{\normalfont{} (#3)}}
\theoremstyle{plain}
\newtheorem{theorem}{Theorem}[section]
\newtheorem{lemma}[theorem]{Lemma}
\newtheorem{proposition}[theorem]{Proposition}
\newtheorem{corollary}[theorem]{Corollary}

\newtheoremstyle{definition}
  {.5\linespacing}{.5\linespacing}
  {\upshape}{1em}{\scshape}{.}{.5em}
  {\thmname{#1}\thmnumber{ #2}\thmnote{\normalfont{} (#3)}}
\theoremstyle{definition}
\newtheorem{definition}[theorem]{Definition}

\newtheoremstyle{remark}
  {.5\linespacing}{.5\linespacing}
  {\upshape}{1em}{\scshape}{.}{.5em}
  {\thmname{#1}\thmnumber{ #2}\thmnote{\normalfont{} (#3)}}
\theoremstyle{remark}
\newtheorem{remark}[theorem]{Remark}

\renewenvironment{proof}[1][\proofname]{%
  \par
  \pushQED{\qed}%
  \normalfont
  \topsep6pt plus6pt\relax
  \trivlist
  \item[\hskip\labelsep\slshape #1\@addpunct{.}]%
  \ignorespaces
}{%
  \popQED
  \endtrivlist
  \@endpefalse
}
\DeclareRobustCommand{\qed}{%
  \leavevmode\unskip\penalty9999\hbox{}\nobreak
  \ifmmode
    \makebox[0pt][l]{\qquad\hbox{\qedsymbol}}%
  \else
    \qquad\hbox{\qedsymbol}%
  \fi
}
\renewcommand{\qedsymbol}{\rule[.47pt]{4pt}{6pt}}

\makeatother

\setlength{\emergencystretch}{1.8em}
\numberwithin{equation}{section}
\allowdisplaybreaks[1]


\usepackage{booktabs,array}
\newcommand{\E}{\mathbb E}
\newcommand{\PP}{\mathbb P}
\newcommand{\R}{\mathbb R}
\newcommand{\Law}{\operatorname{Law}}
\newcommand{\Cov}{\operatorname{Cov}}
\newcommand{\id}{\mathrm{id}}
\newcommand{\dd}{\,\mathrm d}
\newcommand{\Dplus}{D^{+}}
\newcommand{\cJ}{\mathcal J}
\newcommand{\cL}{\mathcal L}
\newcommand{\cF}{\mathscr F}
\newcommand{\res}{\mathfrak R}
\crefname{theorem}{Theorem}{Theorems}
\crefname{lemma}{Lemma}{Lemmas}
\crefname{proposition}{Proposition}{Propositions}
\crefname{corollary}{Corollary}{Corollaries}
\crefname{section}{Section}{Sections}

\title[Monotonicity of the Speed on Galton--Watson Trees]{Monotonicity of the Speed of Biased Random Walks on Galton--Watson Trees}
\author[GPT-6 Astra]{GPT-6 Astra}
\author[Heng Ma]{Heng Ma}
\thanks{Only Heng Ma takes responsibility for this paper.}
\date{November 20, 2026}
\subjclass[2020]{60J10, 60J80, 60K37}
\keywords{Galton--Watson tree, biased random walk, speed monotonicity, escape probability, computer-assisted proof}

\begin{document}
\begin{abstract}
We study the speed of the $\lambda$-biased random walk on a Galton--Watson tree without leaves, where the parent edge has weight $\lambda$ and every child edge has weight one. Under only a finite offspring mean $m>1$, we give two analytic proofs of strict decrease throughout $[0,m)$ when every vertex has at least two children. When single-child vertices are allowed, we prove strict decrease on $(0,1)$ and on $[(3+\sqrt5)/2,m)$. The first proof for $\lambda>1$ modifies the speed formula by a term that vanishes under the escape recursion. The comparison then reduces to an inequality between the mean escape probability and a weighted mean. For $\lambda<1$, we bound the change in the escape probability by an affine map and use a resistance estimate and Jensen's inequality. The second proof uses linear response to express $v'(\lambda)$ in terms of the escape law for $1<\lambda<m$. For laws allowing single-child vertices, we use the same derivative formula. Bounds on moments and series tails then reduce the inequality $v'(\lambda)<0$ to finitely many rational checks. The supplement gives the exact coefficient tables and verification programs.
\end{abstract}

\begin{aiusage}
This manuscript was substantially generated and revised with the assistance of AI systems. The analytic arguments and exposition received AI-assisted checks. The finite arithmetic used in the computer-assisted part was reconstructed and checked with exact arithmetic; the resulting coefficient data and verification programs accompany the paper. These checks do not constitute independent human verification or formal proof verification.
\end{aiusage}

\maketitle
\begingroup
\setcounter{tocdepth}{1}
\footnotesize
\setlength{\linespacing}{8pt}
\setlength{\parskip}{0pt}
\tableofcontents
\endgroup
\clearpage

\section{Introduction}\label{sec:introduction}

Let $\mathbb T$ be a Galton--Watson tree with offspring variable $\nu$ satisfying
\begin{equation}\label{eq:assumptions}
 \PP(\nu=0)=0,\qquad 1<m:=\E\nu<\infty.
\end{equation}
At a nonroot vertex with $k$ children, the $\lambda$-biased random walk moves to the parent with probability $\lambda/(\lambda+k)$ and to each child with probability $1/(\lambda+k)$. At the root it chooses a child uniformly. Write $W_n$ for the walk started at the root and $|W_n|$ for its distance from the root. In the transient regime $0\le\lambda<m$, its speed
\begin{equation*}
 v(\lambda)=\lim_{n\to\infty}\frac{|W_n|}{n}
\end{equation*}
exists almost surely and is a deterministic positive constant; see~\cite{LPP,Aidekon}. Increasing $\lambda$ strengthens the bias towards the root. Lyons, Pemantle and Peres asked whether $v$ is nonincreasing in $\lambda$ for every offspring law satisfying~\eqref{eq:assumptions}~\cite[Question~2.1]{LPP}.

The local drift at a vertex of degree $k+1$ decreases with $\lambda$, but this observation does not settle the question: changing the bias also changes the distribution of the environments visited by the walk. A proof must control both effects. Our analytic arguments make this comparison through the recursive law of the escape probability.

\subsection{Analytic results}
For a real-valued function $f$, denote its upper right Dini derivative by
\begin{equation*}
 \Dplus f(\lambda)=\limsup_{\varepsilon\downarrow0}
 \frac{f(\lambda+\varepsilon)-f(\lambda)}{\varepsilon}.
\end{equation*}
The first result covers the entire transient regime when single-child vertices are absent.

\begin{theorem}\label{thm:twochildren}
Suppose $\PP(\nu\ge2)=1$ and $m=\E\nu<\infty$. Then $v$ is strictly decreasing on $[0,m)$. Moreover,
\begin{equation*}
 \Dplus v(\lambda)<0,\qquad 1<\lambda<m.
\end{equation*}
\end{theorem}

For laws with single-child vertices, we prove two complementary analytic results.

\begin{theorem}\label{thm:low}
Suppose~\eqref{eq:assumptions} holds and $\PP(\nu=1)>0$. Then $v$ is strictly decreasing on $(0,1)$, and $\Dplus v(\lambda)<0$ for every $0<\lambda<1$.
\end{theorem}

\begin{theorem}\label{thm:high}
Suppose~\eqref{eq:assumptions} holds and $\PP(\nu=1)>0$. If $m>(3+\sqrt5)/2$, then $v$ is strictly decreasing on $[(3+\sqrt5)/2,m)$, and
\begin{equation*}
 \Dplus v(\lambda)<0,\qquad \frac{3+\sqrt5}{2}\le\lambda<m.
\end{equation*}
\end{theorem}

All three statements require only a finite first offspring moment. Their proofs in Sections~\ref{sec:preliminaries}--\ref{sec:high} are analytic. For laws with single-child vertices, Theorem~\ref{thm:computer} gives the stronger conclusion $v'(\lambda)<0$ throughout $(1,m)$ by a computer-assisted argument. Section~\ref{sec:computer} describes the reductions to finite rational inequalities and identifies the exact computations; coefficient data and verification programs are supplied separately. Theorems~\ref{thm:twochildren}--\ref{thm:high} are proved independently of that extension. Lemma~\ref{lem:continuity-one} proves continuity at $\lambda=1$. Together with Theorems~\ref{thm:low} and~\ref{thm:computer}, this gives strict decrease throughout $(0,m)$ for laws with single-child vertices as well.

\subsection{Relation to earlier work}
A\"id\'ekon's explicit speed formula~\cite{Aidekon} provides the starting point for comparisons through the escape law. If the minimum offspring number is $m_1\ge2$, Song--Wang--Xiang proved strict decrease on $[0,m_1/(1+\sqrt{1-1/m_1})]$~\cite[Theorem~1.2]{SWX}. This interval contains $\lambda=1$ in its interior, so it overlaps the interval $(1,m)$ on which we prove strict decrease in Theorem~\ref{thm:twochildren}. Bowditch--Tokushige proved differentiability of speed under an exponential offspring moment assumption~\cite[Theorem~1.1]{BT}. The analytic arguments here establish a strictly negative upper right Dini derivative under a finite first moment.

\begin{remark}[Concurrent work]\label{rem:concurrent-work}
Schwarzfischer's recent preprint~\cite{Schwarzfischer}, submitted on November 19, 2026, proves $v'(\lambda)<0$ on $(0,m)$ under the assumptions $\nu\ge4$ almost surely and $m<\infty$. Its introduction also states that related techniques can handle $\nu\ge2$, without presenting that extension. Theorem~\ref{thm:twochildren} gives a full analytic argument under $\nu\ge2$, while Theorems~\ref{thm:low} and~\ref{thm:high} treat laws with single-child vertices. Our representation of $v'(\lambda)$ is equivalent to that in~\cite[Proposition~2.9]{Schwarzfischer}. Both build on A\"id\'ekon's speed formula~\cite{Aidekon} and the linearization of the escape-law recursion. The estimates used here to establish strict negativity of $v'(\lambda)$ differ from those in~\cite{Schwarzfischer}.

The result under $\nu\ge2$ was obtained in September 2026. An early draft containing this result and its proof is preserved in the repository's initial commit, dated September 25, 2026.\footnote{See GitHub commit \href{https://github.com/hengmamath/monotone-speed-biased-RW/commit/9374a17682392e39152f1f39f8e7193ef8365c97}{\texttt{9374a17}}.} Posting to arXiv or Hexagon was postponed while we sought a proof of the full conjecture that would not require computer verification.
\end{remark}

\subsection{Comparing the speed at two biases}
Attach an extra parent to the root of a Galton--Watson tree. For $0<\lambda<m$, $\lambda\ne1$, let $X,Y$ be independent copies of the quenched probability that the walk, started at the root, never reaches that parent. The speed representation used throughout the paper is
\begin{equation}\label{eq:intro-speed}
 v(\lambda)=
 \frac{\displaystyle\E\frac{XY}{\lambda-1+X+Y-XY}}
 {\displaystyle\E\frac{X+Y-XY}{\lambda-1+X+Y-XY}}.
\end{equation}
The law of $X$ decreases in stochastic order when $\lambda$ increases. For $\lambda>1$, however, the sign of the corresponding speed change is not immediate from~\eqref{eq:intro-speed}. Fix $v=v(\lambda)$ and subtract $v$ times the denominator expectation from the numerator expectation. To control how this difference changes with the law of $X$, we subtract a suitable multiple of
\begin{equation*}
 \E u(X)-\E u\left(\frac{\sum_{i=1}^{\nu}X_i}{\lambda+\sum_{i=1}^{\nu}X_i}\right),
\end{equation*}
where the $X_i$ are independent copies of $X$, independent of $\nu$. This expression vanishes by the escape recursion. We choose the derivative of $u$ to cancel one term in the change with respect to the law. The remaining change in the law and the change in the bias with the law fixed have the required signs whenever
\begin{equation}\label{eq:intro-weighted}
 \E X<2\frac{\E[X/(\lambda-1+X)^2]}{\E[1/(\lambda-1+X)^2]}.
\end{equation}
Proposition~\ref{prop:criterion} makes this implication precise. Weighting the joint law of the children by $(\lambda-1+\sum_iX_i)^{-2}$ gives~\eqref{eq:intro-weighted} for $1<\lambda<m$ if there are at least two children, and for $(3+\sqrt5)/2\le\lambda<m$ if there is at least one child.

For $\lambda<1$, the deterministic bound $X\ge1-\lambda$ gives an affine upper bound on the escape probability at a larger bias. The speed quotient preserves stochastic order for laws supported on $[1-\lambda,1]$, so the affine bound gives an upper bound on the new speed. We prove that its derivative is negative using the mean resistance $\E(1/X)$ and two applications of Jensen's inequality. Section~\ref{sec:low} gives the complete argument.

Section~\ref{sec:linear-response} proves local Lipschitz continuity of the escape law and a formula for the derivative of its expectations. A positive inverse of the resulting linear operator expresses $v'(\lambda)$ as the negative of a quantity to be estimated, divided by a positive denominator. This formula gives the second proof in Section~\ref{sec:two-child-response} and the starting point for the computer-assisted estimates in Section~\ref{sec:computer}. Moment inequalities and an identity involving a logarithmic function give lower bounds for that quantity. For each of four bias intervals, estimates on the moments and infinite tails reduce the required strict sign to finitely many rational inequalities. The exact computations verify those inequalities on domains that contain the parameters of every escape law under consideration.

\section{Preliminaries}\label{sec:preliminaries}

Throughout the paper the offspring variable satisfies $\PP(\nu\ge1)=1$ and $1<m=\E\nu<\infty$. The offspring law is fixed when the bias changes. We first express the speed through two independent escape probabilities. For $\lambda>1$, we establish the kernel identities and linear-response formula used in the two proofs in Section~\ref{sec:twochildren} and in the later estimates.

\subsection{The escape law and the speed formula}
Attach an artificial parent to the root of a descendant tree $T$, and let $\beta_\lambda(T)$ be the quenched probability, starting from that root, of never visiting its parent. For $0<\lambda<m$, transience gives $0<\beta_\lambda(T)<1$ almost surely. Write $\pi_\lambda=\Law(\beta_\lambda)$ and
\begin{equation}\label{eq:rde}
 F_\lambda(s)=\frac{s}{\lambda+s},\qquad
 X\overset{\mathrm d}=F_\lambda(S),\qquad S=\sum_{i=1}^{\nu}X_i,
\end{equation}
where the $X_i$ have law $\pi_\lambda$, independently of one another and of $\nu$. To obtain the recursion, condition on the child subtrees. A step to child $i$ leads to escape before returning to the present vertex with probability $X_i$; after a return the same experiment restarts. The first-step equation is
\[
 X=\frac1{\lambda+\nu}\sum_{i=1}^{\nu}
 \{X_i+(1-X_i)X\},
\]
which rearranges to~\eqref{eq:rde}. At finite depth, terminal escape value one is the boundary condition for reaching that generation before the artificial parent. Those events decrease with the terminal depth, and their probabilities tend to the probability of never hitting the parent. Since $F_\lambda$ increases in its argument and decreases in $\lambda$, induction on the finite recursion followed by this limit gives the same-tree order
\begin{equation}\label{eq:same-tree}
 \beta_\alpha(T)\ge\beta_\gamma(T),\qquad 0<\alpha<\gamma<m.
\end{equation}

Set
\begin{equation}\label{eq:kernels}
 t(x,y)=x+y-xy,\qquad A_\lambda(x,y)=\lambda-1+t(x,y).
\end{equation}
For independent $X,Y$ with common probability law $\rho$ on $[0,1]$, define
\begin{equation}\label{eq:NZ}
 N_\lambda(\rho)=\E_\rho\frac{XY}{A_\lambda(X,Y)},\qquad
 Z_\lambda(\rho)=\E_\rho\frac{t(X,Y)}{A_\lambda(X,Y)},\qquad
 \mathcal V_\lambda(\rho)=\frac{N_\lambda(\rho)}{Z_\lambda(\rho)}
\end{equation}
whenever the integrals are finite and the denominator is positive. If $\lambda>1$, then $\lambda-1\le A_\lambda\le\lambda$ on $[0,1]^2$. If $0<\lambda<1$, the finite recursion gives $X\ge1-\lambda$ under $\pi_\lambda$, and
\begin{equation}\label{eq:low-denominator}
 A_\lambda(x,y)\ge\lambda(1-\lambda)>0
 \quad\text{for }x,y\in[1-\lambda,1].
\end{equation}

\begin{proposition}[Symmetric speed formula]\label{prop:symmetric}
For $0<\lambda<m$, $\lambda\ne1$, the speed satisfies
\begin{equation}\label{eq:pair-speed}
 v(\lambda)=\mathcal V_\lambda(\pi_\lambda).
\end{equation}
\end{proposition}

\begin{proof}
We use A\"id\'ekon's speed formula~\cite[Theorem~1.1]{Aidekon}, in the leafless form stated in~\cite[(2.1)]{SWX}:
\begin{equation}\label{eq:degree-speed}
 v(\lambda)=
 \frac{\E[(\nu-\lambda)X_0/(\lambda-1+\sum_{i=0}^{\nu}X_i)]}
 {\E[(\nu+\lambda)X_0/(\lambda-1+\sum_{i=0}^{\nu}X_i)]}.
\end{equation}
The $X_i$ are independent with law $\pi_\lambda$, independently of $\nu$. The speed is unchanged by attaching an artificial parent at the root; see~\cite[Lemma~2.1]{SWX}. Conditional exchangeability yields
\begin{equation*}
 \E\frac{\nu X_0}{\lambda-1+X_0+S}
 =\E\frac{S}{\lambda-1+X_0+S}.
\end{equation*}
Put $X=X_0$ and $Y=F_\lambda(S)$. Then $X,Y$ are independent with law $\pi_\lambda$ and $S=\lambda Y/(1-Y)$. Multiplication of the denominator by $1-Y$ gives $A_\lambda(X,Y)$. The numerator of~\eqref{eq:degree-speed} becomes
\begin{equation*}
 \lambda\E\frac{Y-X+XY}{A_\lambda(X,Y)}
 =\lambda N_\lambda(\pi_\lambda),
\end{equation*}
because $A_\lambda$ is symmetric and $X,Y$ have the same law. The denominator of~\eqref{eq:degree-speed} becomes $\lambda Z_\lambda(\pi_\lambda)$. The bounds preceding the proposition and $\E\nu<\infty$ justify these operations.
\end{proof}

The following lemma will join the strict comparisons on either side of bias one.

\begin{lemma}[Continuity at bias one]\label{lem:continuity-one}
Suppose that $\PP(\nu\ge1)=1$ and $1<m=\E\nu<\infty$. Then
\begin{equation*}
 \lim_{\lambda\to1}v(\lambda)
 =v(1)=\E\left[\frac{\nu-1}{\nu+1}\right].
\end{equation*}
\end{lemma}

\begin{proof}
For each bias $\lambda$ near one, let $X_0,X_1,\ldots$ be independent with law $\pi_\lambda$, independently of $\nu$. Couple these variables for different biases by using the same independent trees, with $X_0=\beta_\lambda(T)$.

A\"id\'ekon's formula~\eqref{eq:degree-speed} holds under the stated assumptions for all $\lambda$ in a neighborhood of one, including $\lambda=1$~\cite[Theorem~1.1]{Aidekon}. Conditional on $\nu$, the common denominator inside the expectations is symmetric in $X_0,\ldots,X_\nu$. Averaging the distinguished variable $X_0$ over these $\nu+1$ variables therefore gives
\begin{equation}\label{eq:speed-exchangeable}
 v(\lambda)=
 \frac{
  \displaystyle\E\left[
   \frac{\nu-\lambda}{\nu+1}
   \frac{\sum_{i=0}^{\nu}X_i}
        {\lambda-1+\sum_{i=0}^{\nu}X_i}
  \right]
 }{
  \displaystyle\E\left[
   \frac{\nu+\lambda}{\nu+1}
   \frac{\sum_{i=0}^{\nu}X_i}
        {\lambda-1+\sum_{i=0}^{\nu}X_i}
  \right]
 }.
\end{equation}

Fix $\Lambda\in(1,m)$. By~\eqref{eq:same-tree} and transience at bias $\Lambda$, almost surely,
\begin{equation*}
 \sum_{i=0}^{\nu}X_i
 \ge X_0=\beta_\lambda(T)
 \ge\beta_\Lambda(T)>0,
 \qquad 0<\lambda<\Lambda.
\end{equation*}
Consequently, whenever $|\lambda-1|<\beta_\Lambda(T)/2$,
\begin{equation*}
 \left|
 \frac{\sum_{i=0}^{\nu}X_i}
      {\lambda-1+\sum_{i=0}^{\nu}X_i}-1
 \right|
 \le\frac{2|\lambda-1|}{\beta_\Lambda(T)}
 \longrightarrow0.
\end{equation*}

We next bound the integrands in~\eqref{eq:speed-exchangeable} uniformly near one. For $\lambda\ge1$,
\begin{equation*}
 0<\frac{\sum_{i=0}^{\nu}X_i}
        {\lambda-1+\sum_{i=0}^{\nu}X_i}\le1.
\end{equation*}
For $0<\lambda<1$, the escape recursion with terminal value one gives $X_i\ge1-\lambda$. Indeed, if every child value is at least $1-\lambda$, then their sum $S$ is at least $1-\lambda$, and
\begin{equation*}
 \frac{S}{\lambda+S}
 \ge\frac{1-\lambda}{\lambda+1-\lambda}=1-\lambda.
\end{equation*}
Induction followed by the limit in depth proves the bound. Since $\nu\ge1$, it follows that
\begin{equation*}
 0<\frac{\sum_{i=0}^{\nu}X_i}
        {\lambda-1+\sum_{i=0}^{\nu}X_i}
 =1+\frac{1-\lambda}{\sum_{i=0}^{\nu}X_i-(1-\lambda)}
 \le1+\frac1\nu\le2.
\end{equation*}
Also, for $0<\lambda<\Lambda$,
\begin{equation*}
 \frac{|\nu-\lambda|}{\nu+1}\le1+\Lambda,
 \qquad
 \frac{\nu+\lambda}{\nu+1}\le1+\Lambda.
\end{equation*}
Thus both integrands in~\eqref{eq:speed-exchangeable} are bounded in absolute value by $2(1+\Lambda)$. Dominated convergence shows that the numerator tends to $\E[(\nu-1)/(\nu+1)]$ and the denominator tends to one. Evaluating~\eqref{eq:speed-exchangeable} at $\lambda=1$ gives
\begin{equation*}
 \lim_{\lambda\to1}v(\lambda)
 =\E\left[\frac{\nu-1}{\nu+1}\right]
 =v(1).\qedhere
\end{equation*}
\end{proof}

\subsection{Continuity of the escape law and speed for \texorpdfstring{$1<\lambda<m$}{1 < lambda < m}}
The comparison of nearby biases requires continuity of $\pi_\lambda$. We prove it from the recursion using only $\E\nu<\infty$.

\begin{lemma}\label{lem:continuity}
The escape law $\pi_\lambda$ is continuous in the $1$-Wasserstein metric on $(1,m)$. Consequently, $v$ is continuous on $(1,m)$.
\end{lemma}

\begin{proof}
By~\eqref{eq:same-tree}, the one-sided limits of the escape probabilities exist. We will identify their laws by proving uniqueness of the nonzero fixed law of~\eqref{eq:rde}. To distinguish two ordered fixed laws, we choose a test function whose expected change under one recursion step is strictly monotone along their linear interpolation if the laws differ. That change vanishes at both fixed laws, forcing the laws to agree.

Fix $1<\lambda<m$ and write $c=\lambda-1$. For a nonzero fixed law $\rho$, define
\begin{equation*}
 r_\rho(x)=\E_\rho\frac{Y}{A_\lambda(x,Y)^2},\qquad
 u_\rho(x)=\int_0^x r_\rho(s)\dd s.
\end{equation*}
To differentiate expectations after one step of the recursion, we use the following operator. If the children have law $\sigma$, let
\begin{equation}\label{eq:linear-operator}
 (\cL_{\lambda,\sigma}r)(x)=
 \E_{\nu,\sigma}\left[
 \frac{\nu\lambda}{(\lambda+x+\sum_{i=2}^{\nu}X_i)^2}
 r\left(F_\lambda\left(x+\sum_{i=2}^{\nu}X_i\right)\right)\right].
\end{equation}
The denominator identity
\begin{equation}\label{eq:denominator-cancellation}
 A_\lambda(F_\lambda(s),y)=\frac{\lambda(c+s+y)}{\lambda+s}
\end{equation}
gives
\begin{equation*}
 \frac{\lambda}{(\lambda+s)^2}r_\rho(F_\lambda(s))
 =\frac1\lambda\E_\rho\frac{Y}{(c+s+Y)^2}.
\end{equation*}
The right-hand side decreases in $s$. If $\sigma$ stochastically dominates $\rho$, take independent copies of an ordered coupling, with the $\sigma$ variable above the $\rho$ variable. This gives $\cL_{\lambda,\sigma}r_\rho\le\cL_{\lambda,\rho}r_\rho$. At $\sigma=\rho$, the variable $Y$ and the remaining children form $\nu$ independent variables with law $\rho$. Conditional exchangeability and then the fixed-point equation give
\[
 (\cL_{\lambda,\rho}r_\rho)(x)
 =\frac1\lambda\E_\rho\frac{\nu Y}{(c+x+\sum_{i=2}^{\nu}X_i+Y)^2}
 =\frac1\lambda\E_\rho\frac{S}{(c+x+S)^2}
 =\E_\rho\frac{Y(1-Y)}{A_\lambda(x,Y)^2}.
\]
Subtracting this identity from the definition of $r_\rho$ proves
\begin{equation}\label{eq:positive-gap}
 r_\rho(x)-(\cL_{\lambda,\rho}r_\rho)(x)
 =\E_\rho\frac{Y^2}{A_\lambda(x,Y)^2}
 \ge\frac{\E_\rho Y^2}{\lambda^2}>0.
\end{equation}

To compare two ordered fixed laws, couple $X_0\le X_1$ with laws $\rho,\sigma$ and interpolate $X_s=(1-s)X_0+sX_1$. For the recursion residual
\begin{equation*}
 \res_{u_\rho}(\lambda,\sigma)
 =\E_\sigma u_\rho(X)
 -\E_{\nu,\sigma}u_\rho\left(F_\lambda\left(\sum_{i=1}^{\nu}X_i\right)\right),
\end{equation*}
the chain rule along independent copies of this coupling yields
\begin{align*}
 \frac{\dd}{\dd s}\res_{u_\rho}(\lambda,\Law(X_s))
 &=\E\left[(X_1-X_0)
 \{r_\rho-\cL_{\lambda,\Law(X_s)}r_\rho\}(X_s)\right]\\
 &\ge\frac{\E_\rho Y^2}{\lambda^2}\E(X_1-X_0).
\end{align*}
For a fixed number of children, this is the ordinary chain rule: the derivative of the sum is the sum of the coupled increments, and their contributions agree by exchangeability. The recursion term is bounded by $\nu\|r_\rho\|_\infty/\lambda$, since each increment is at most one; the direct term is bounded by $\|r_\rho\|_\infty$. These bounds justify differentiation under the offspring expectation. The residual vanishes at both $s=0$ and $s=1$, because both endpoint laws satisfy the recursion. Integration in $s$ therefore forces $\E(X_1-X_0)=0$ and hence $\rho=\sigma$. Iteration from terminal value one dominates every fixed law: couple its terminal ones above independent samples of that fixed law and iterate. Its limit has law $\pi_\lambda$, so the preceding ordered comparison proves uniqueness of the nonzero fixed law.

At any $\lambda\in(1,m)$, each one-sided same-tree limit of $\beta_\alpha$ as $\alpha\to\lambda$ satisfies~\eqref{eq:rde}. This follows by bounded convergence, since the number of children is finite almost surely. Choose $\Lambda\in(\lambda,m)$. The one-sided limits dominate $\beta_\Lambda>0$, and hence their laws are nonzero. Uniqueness identifies both with $\pi_\lambda$. The order~\eqref{eq:same-tree} then gives
\begin{equation*}
 W_1(\pi_\alpha,\pi_\lambda)
 =\E|\beta_\alpha-\beta_\lambda|\longrightarrow0.
\end{equation*}
The kernels in~\eqref{eq:NZ} are bounded and continuous for biases near $\lambda>1$. Since $Z_\lambda(\pi_\lambda)>0$, the speed formula gives continuity of $v$.
\end{proof}

\subsection{Kernel identities and the weighted mean}\label{sec:kernel-identities}
Fix $1<\lambda<m$, put $c=\lambda-1$, and let $X,Y$ be independent with law $\pi_\lambda$. In the rest of this subsection, $v=v(\lambda)$, $A=A_\lambda$, $N=N_\lambda(\pi_\lambda)$ and $Z=Z_\lambda(\pi_\lambda)$ are fixed. Define the mean with weight $(c+X)^{-2}$ by
\begin{equation}\label{eq:weighted-mean}
 k=\frac{\E[X/(c+X)^2]}{\E[1/(c+X)^2]}.
\end{equation}
Define
\begin{equation}\label{eq:J-h}
 h_v(y)=(1+v)y-v,\qquad
 (\cJ f)(x)=\E\frac{f(Y)}{A(x,Y)^2}.
\end{equation}

The derivative of the integrand of $N-vZ$ in its first coordinate is
\begin{equation}\label{eq:variation-kernel}
 \partial_x\frac{xy-vt(x,y)}{A(x,y)}
 =\frac{y^2+c h_v(y)}{A(x,y)^2}.
\end{equation}
With $\cL=\cL_{\lambda,\pi_\lambda}$ from~\eqref{eq:linear-operator}, set
\begin{equation}\label{eq:correction}
 r(x)=\E\frac{Y}{A(x,Y)^2},\qquad u(x)=\int_0^x r(s)\dd s.
\end{equation}
Identity~\eqref{eq:denominator-cancellation}, exchangeability of the children, and~\eqref{eq:rde} yield
\begin{align}\label{eq:Lr}
 (\cL r)(x)
 &=\frac1\lambda\E\frac{\nu Y}{(c+x+\sum_{i=2}^{\nu}X_i+Y)^2}
 =\frac1\lambda\E\frac{S}{(c+x+S)^2}\notag\\
 &=\E\frac{Y(1-Y)}{A(x,Y)^2}.
\end{align}
Thus $(I-\cL)r=\cJ(y^2)$. This identity will both cancel a term in the direct comparison and give a positive inverse of $I-\cL$.

The explicit bias dependence of the speed formula and of the recursion will be expressed through
\begin{equation}\label{eq:parameter-margin}
 b(x)=\frac{x(1-x)}{\lambda},\qquad
 C=\E\frac{v\,t(X,Y)-XY}{A(X,Y)^2},\qquad
 K=2\E[b(X)r(X)]-C.
\end{equation}
By symmetry,
\begin{equation*}
 \lambda K=\E\frac{XY(2-X-Y)+\lambda\{XY-vt(X,Y)\}}{A(X,Y)^2}.
\end{equation*}
Use $\lambda-A(x,y)=(1-x)(1-y)$ and $N-vZ=0$ to remove the term with a single power of $A$ in the denominator. The remaining numerator satisfies
\begin{align*}
 &(1-x)(1-y)\{xy-vt(x,y)\}+xy(2-x-y)\\
 &\qquad=x(1-x)h_v(y)+y(1-y)h_v(x)
 +(1+v)xy(1-x)(1-y).
\end{align*}
After division by $A(x,y)^2$ and averaging, this gives
\begin{equation}\label{eq:positive-K}
 \lambda K=2\E[X(1-X)(\cJ h_v)(X)]
 +(1+v)\E\frac{XY(1-X)(1-Y)}{A(X,Y)^2}.
\end{equation}

\begin{lemma}[Positivity from the weighted mean]\label{lem:weighted-positivity}
If $\E X<2k$, then
\begin{equation*}
 \inf_{0\le x\le1}(\cJ h_v)(x)>0,\qquad K>0.
\end{equation*}
\end{lemma}

\begin{proof}
We first show that $\cJ h_v$ is uniformly positive. Put $W=\E[X/A(X,Y)]$. Symmetry gives $N+Z=2W$. For fixed $x$, the function $y\mapsto A(x,y)^{-1}$ is decreasing; its covariance with $Y$ is nonpositive. Hence $N\le(\E X)W$ and
\begin{equation}\label{eq:weighted-speed}
 \frac{v}{1+v}=\frac{N}{2W}\le\frac{\E X}{2}<k.
\end{equation}
Relative to the law weighted by $(c+y)^{-2}$, the law weighted by $A(x,y)^{-2}$ has density proportional to $\{(c+y)/A(x,y)\}^2$. This density increases in $y$, because
\begin{equation*}
 \frac{\partial}{\partial y}\frac{c+y}{A(x,y)}
 =\frac{\lambda x}{A(x,y)^2}\ge0.
\end{equation*}
The mean of $Y$ under this weight is therefore at least $k$. Since $A\le\lambda$,
\begin{equation}\label{eq:J-positive}
 (\cJ h_v)(x)\ge\frac{(1+v)k-v}{\lambda^2}>0,
 \qquad 0\le x\le1.
\end{equation}

By~\eqref{eq:positive-K}, this positivity and $0<X,Y<1$ give $K>0$.
\end{proof}

\subsection{Linear response of the escape law}\label{sec:linear-response}
Fix $1<\lambda<m$. Throughout this subsection, $c=\lambda-1$, $X,Y$ are independent with law $\pi_\lambda$, and $v=v(\lambda)$. We retain $A$, $\cJ$, $\cL$, $r$, $b$ and $C$ from Section~\ref{sec:kernel-identities}, evaluated at this bias and law. To differentiate the speed formula we use
\begin{equation}\label{eq:full-response-defs}
 \Psi=\cJ\bigl(y^2+c h_v(y)\bigr),\qquad
 \mathcal R=(I-\cL)^{-1}.
\end{equation}
We first construct the inverse $\mathcal R$, then use it to differentiate expectations under $\pi_\lambda$ using only the finite mean of $\nu$.

\begin{lemma}[Positive inverse]\label{lem:positive-inverse}
The operator $\cL$ is bounded and positive on $C([0,1])$. Its Neumann series $\mathcal R=\sum_{j\ge0}\cL^j$ converges in operator norm and defines a positive inverse of $I-\cL$. In a neighborhood of any fixed $\lambda\in(1,m)$, the tails of the series admit a common geometric bound in an equivalent norm.
\end{lemma}
\begin{proof}
For $f\in C([0,1])$, the integrand defining $\cL f$ is bounded in absolute value by $\nu\|f\|_\infty/\lambda$. Thus finite mean gives boundedness, and dominated convergence gives continuity in $x$. Positivity follows directly from the formula for $\cL$. Identity~\eqref{eq:Lr} gives
\[
 (I-\cL)r=\cJ(y^2),\qquad
 r(x)\ge\frac{\E Y}{\lambda^2}>0,\qquad
 \cJ(y^2)(x)\ge\frac{\E Y^2}{\lambda^2}>0.
\]
The norm $\|f\|_r=\sup_{0\le x\le1}|f(x)|/r(x)$ is equivalent to the uniform norm. Since $\cL$ is positive,
\[
 |\cL f|\le\|f\|_r\cL r,
 \qquad \|\cL\|_r\le
 1-\min_{0\le x\le1}\frac{\cJ(y^2)(x)}{r(x)}<1.
\]
The geometric series therefore converges and gives the asserted positive inverse.

For the local statement, keep this base function $r$ fixed. If $\alpha\to\lambda$ and $\rho\to\pi_\lambda$ in $W_1$, then
\begin{equation}\label{eq:fixed-test-continuity}
 \|\cL_{\alpha,\rho}f-\cL f\|_\infty\longrightarrow0
 \quad\text{for each fixed }f\in C([0,1]).
\end{equation}
Indeed, for $\nu\le M$ there are only finitely many continuous kernels on compact cubes, and their integrals converge uniformly in $x$. For biases $\alpha\ge a>1$, the sum of the two discarded tails is at most $2\|f\|_\infty\E[\nu;\nu>M]/a$. First choose $M$ and then take the limit. Apply~\eqref{eq:fixed-test-continuity} with $f=r$. The strictly positive minimum of $r-\cL r$ persists for nearby operators $\cL_{\alpha,\pi_\alpha}$, so their norms in $\|\cdot\|_r$ have a common bound strictly below one. This gives the local geometric bound using convergence on each fixed test function.
\end{proof}


\begin{lemma}[Local Lipschitz continuity]\label{lem:law-lipschitz}
The map $\lambda\mapsto\pi_\lambda$ is locally Lipschitz in $W_1$ on $(1,m)$.
\end{lemma}
\begin{proof}
Continuity was proved in Lemma~\ref{lem:continuity} without differentiation. Fix the base parameter, put $u(x)=\int_0^x r(s)\dd s$, and retain the recursion residual
\[
 \res_u(\lambda,\rho)=\E_\rho u(X)
 -\E_{\nu,\rho}u\left(F_\lambda\left(\sum_{i=1}^{\nu}X_i\right)\right).
\]
By~\eqref{eq:fixed-test-continuity}, every law on the ordered interpolation between $\pi_\lambda$ and $\pi_{\lambda+\varepsilon}$ satisfies
\[
 r-\cL_{\lambda,\rho}r\ge\tfrac12\min_{[0,1]}\cJ(y^2)
\]
when $|\varepsilon|$ is small. Integrating the residual's law derivative along that interpolation, in its increasing direction, gives
\[
 |\res_u(\lambda,\pi_{\lambda+\varepsilon})|
 \ge\tfrac12\min_{[0,1]}\cJ(y^2)\,
 W_1(\pi_{\lambda+\varepsilon},\pi_\lambda).
\]
The residual at $\pi_\lambda$ is zero. The fixed-point equation at $\lambda+\varepsilon$ also gives
\[
 \res_u(\lambda,\pi_{\lambda+\varepsilon})
 =\E_{\nu,\pi_{\lambda+\varepsilon}}
 \left[u(F_{\lambda+\varepsilon}(S))-u(F_\lambda(S))\right].
\]
For $\alpha\ge a>1$ and $s\ge0$,
\[
 |\partial_\alpha F_\alpha(s)|=\frac{s}{(\alpha+s)^2}
 \le\frac1{4a}.
\]
The residual is consequently at most $|\varepsilon|\|r\|_\infty/(4a)$ in absolute value. Combining the two inequalities proves the bound centered at the fixed parameter. To obtain a common Lipschitz constant on a neighborhood, keep the same base function $r$ and use $\res_u(\alpha,\pi_\beta)$ for any two nearby biases $\alpha,\beta$. The continuity in~\eqref{eq:fixed-test-continuity} preserves the same positive lower bound for $r-\cL_{\alpha,\rho}r$ for all their interpolating laws. The residual at $\pi_\alpha$ is zero, and its parameter bound is at most $|\beta-\alpha|\|r\|_\infty/(4a)$. Thus the same constant bounds $W_1(\pi_\alpha,\pi_\beta)/|\beta-\alpha|$ throughout a sufficiently small neighborhood. Neither estimate requires a second offspring moment.
\end{proof}

\begin{lemma}[Derivatives of expectations and of the speed]\label{lem:full-response}
For every $f\in C^1([0,1])$,
\begin{equation}\label{eq:weak-response}
 \frac{\dd}{\dd\lambda}\E f(\beta_\lambda)
 =-\E[b(X)(\mathcal R f')(X)].
\end{equation}
The derivative is continuous on $(1,m)$. In particular, $v\in C^1((1,m))$ and
\begin{equation}\label{eq:full-sign}
 v'(\lambda)=-\frac{D_\lambda}{Z},\qquad
 D_\lambda=2\E[b(X)(\mathcal R\Psi)(X)]-C,
 \qquad Z=Z_\lambda(\pi_\lambda)>0.
\end{equation}
\end{lemma}
\begin{proof}
Couple $X=\beta_\lambda$ and $X_\varepsilon=\beta_{\lambda+\varepsilon}$ on the same tree. For $\varepsilon\ne0$ small, define a positive finite measure by
\[
 \eta_\varepsilon(a)=\E\left[\frac1\varepsilon
 \int_{X_\varepsilon}^{X}a(s)\dd s\right],
 \qquad a\in C([0,1]).
\]
Positivity holds for both signs of $\varepsilon$, by the same-tree order. Lemma~\ref{lem:law-lipschitz} bounds its mass:
\[
 \eta_\varepsilon(1)=
 \frac{W_1(\pi_{\lambda+\varepsilon},\pi_\lambda)}{|\varepsilon|}.
\]
Let $u$ be a primitive of a fixed continuous $a$. Then
$(\E u(X_\varepsilon)-\E u(X))/\varepsilon=-\eta_\varepsilon(a)$.
Use the fixed-point equations to compute the same difference by first changing the bias at fixed child law $\pi_{\lambda+\varepsilon}$, and then changing the child law at fixed bias $\lambda$. The parameter contribution tends to $-\E[b(X)a(X)]$. Its integrand is bounded by $\|a\|_\infty/(4a_0)$ for a positive lower bound $a_0$ on the nearby biases, because $|\partial_\alpha F_\alpha(s)|\le1/(4a_0)$.

For the child-law contribution, let $X_t=(1-t)X+tX_\varepsilon$ and $\rho_t=\Law(X_t)$. Differentiating the finite-child recursion and using exchangeability gives the exact quotient
\[
 -\int_0^1\E\left[
 \frac{X-X_\varepsilon}{\varepsilon}
 (\cL_{\lambda,\rho_t}a)(X_t)\right]\dd t.
\]
For each fixed $a$, finite-mean truncation as in~\eqref{eq:fixed-test-continuity} yields
\[
 \sup_{0\le t\le1}
 \|\cL_{\lambda,\rho_t}a-\cL a\|_\infty\longrightarrow0.
\]
Multiplying this uniform error by the bounded mass $\eta_\varepsilon(1)$ shows that the child-law quotient is $-\eta_\varepsilon(\cL a)+o(1)$. Hence
\begin{equation}\label{eq:response-difference-measures}
 \eta_\varepsilon((I-\cL)a)=\E[ba]+o(1).
\end{equation}
Take $a=\mathcal R f'$, which is continuous by Lemma~\ref{lem:positive-inverse}. Equation~\eqref{eq:response-difference-measures} gives~\eqref{eq:weak-response} for two-sided increments. This passage uses bounded expectations of $(X-X_\varepsilon)/\varepsilon$ and a uniform test-function error; it does not require uniform integrability of the random difference quotients.

For continuity of the derivative, nearby operators converge on each fixed continuous test by~\eqref{eq:fixed-test-continuity}. Each finite Neumann sum therefore converges strongly, by induction on its length and uniform boundedness of the operators. The common geometric tail bound in Lemma~\ref{lem:positive-inverse} then proves strong convergence of the resolvents. Together with $W_1$ continuity of $\pi_\lambda$ and continuity of $b$, this proves continuity in~\eqref{eq:weak-response}.

Finally, freeze $v=v(\lambda)$ in
\[
 H_\alpha(x,y)=\frac{xy-vt(x,y)}{\alpha-1+t(x,y)}.
\]
On any compact bias interval in $(1,m)$, this kernel and its first derivatives are bounded and uniformly continuous on $[0,1]^2$. Its averaged coordinate derivative at the base parameter is $\Psi$, by~\eqref{eq:variation-kernel}; its averaged explicit parameter derivative is $C$. Applying~\eqref{eq:weak-response} to the two marginals gives
\[
 \left.\frac{\dd}{\dd\alpha}
 \E H_\alpha(\beta_\alpha,\widetilde\beta_\alpha)
 \right|_{\alpha=\lambda}
 =C-2\E[b\mathcal R\Psi].
\]
Here the two escapes are independent. For precision, the averaged test in one marginal can depend on the increment through the other marginal. Its derivative converges uniformly to the base averaged derivative, by compact uniform continuity and $W_1$ convergence. The resulting difference-quotient error is bounded by that uniform difference times $\eta_\varepsilon(1)$, so it tends to zero. The same argument applied separately to $N_\lambda$ and $Z_\lambda$ first establishes differentiability of their quotient. Since the last display equals $Z v'(\lambda)$, it proves~\eqref{eq:full-sign}. Strong resolvent continuity proves that $v'$ is continuous.
\end{proof}


\section{Strict decrease on \texorpdfstring{$[0,m)$}{[0,m)} when every vertex has at least two children}\label{sec:twochildren}

We give two proofs of Theorem~\ref{thm:twochildren}: one compares nearby biases directly, and the other uses the linear-response formula. Both use the inequality $\E X<2k$ for $1<\lambda<m$, which we prove first. We weight the joint law of the children by $(\lambda-1+\sum_iX_i)^{-2}$. Under this measure, the mean of the first child is at least $k$. Conditional exchangeability then turns $\nu\ge2$ into a lower bound on their sum.

\begin{lemma}\label{lem:two-child}
Let $\lambda>1$, $c=\lambda-1$, and let $X\in(0,1)$ satisfy~\eqref{eq:rde} with $\nu\ge2$ almost surely. For $k$ defined in~\eqref{eq:weighted-mean}, one has $\E X<2k$.
\end{lemma}

\begin{proof}
Let $M$ be the probability law obtained by weighting the law of $X$ with $(c+X)^{-2}$, and write
\begin{equation}\label{eq:tilted-moments}
 k=\E_M X,\qquad j=\E_M X^2,\qquad \ell=\E_M X^3.
\end{equation}
Undoing this change of measure expresses the original mean as
\begin{equation}\label{eq:untilted-mean}
 \E X=\frac{c^2k+2cj+\ell}{c^2+2ck+j}.
\end{equation}
We will use $\ell\le j$ and an upper bound on $j$ to prove $\E X<2k$. To obtain the bound on $j$, weight the joint law of $(\nu,X_1,X_2,\ldots)$ by $(c+S)^{-2}$ and denote the resulting probability measure by $Q$. Under the original law, $X_1$ is independent of $H=\sum_{i=2}^{\nu}X_i$. Its marginal under $Q$, relative to $M$, has density proportional to
\begin{equation*}
 R(x)=\E\left[\left(\frac{c+x}{c+x+H}\right)^2\right].
\end{equation*}
The function $R$ is increasing. The covariance inequality for two increasing functions therefore gives
\begin{equation*}
 \E_QX_1=\frac{\E_M[XR(X)]}{\E_MR(X)}\ge k.
\end{equation*}
For example, the covariance sign follows by taking independent $U,V$ with law $M$ and writing $2\Cov_M(X,R(X))=\E[(U-V)(R(U)-R(V))]\ge0$. Conditional exchangeability under the symmetric weight and $\nu\ge2$ now imply
\begin{equation}\label{eq:two-child-constraint}
 \E_QS=\E_Q[\nu X_1]\ge2\E_QX_1\ge2k.
\end{equation}
By~\eqref{eq:rde}, $S$ has the same law as $\lambda X/(1-X)$ and $c+S$ has the same law as $(c+X)/(1-X)$. Thus
\begin{equation*}
 \E_QS=\lambda\frac{k-j}{1-2k+j}.
\end{equation*}
The denominator is $1-2k+j=\E_M(1-X)^2>0$. Combining this formula with~\eqref{eq:two-child-constraint} gives
\begin{equation}\label{eq:two-j-bound}
 j\le\frac{k(c-1+4k)}{c+1+2k}.
\end{equation}
Since $j\ge k^2$ and
\begin{equation*}
 \frac{k(c-1+4k)}{c+1+2k}-k^2
 =\frac{k(1-k)(c+2k-1)}{c+1+2k},
\end{equation*}
we also have $c+2k-1\ge0$.

If $k\ge1/2$, then $\E X<1\le2k$. Suppose $k<1/2$. First use $\ell\le j$ in~\eqref{eq:untilted-mean}. Since $2k-2c-1<0$, the upper bound for $j$ in~\eqref{eq:two-j-bound} then gives a lower bound for $2k-\E X$:
\begin{align*}
 &(2k-\E X)(c^2+2ck+j)\\
 &\quad\ge c^2k+4ck^2+(2k-2c-1)\frac{k(c-1+4k)}{c+1+2k}\\
 &\quad=\frac{k}{c+1+2k}\Bigl[
 (c+2k-1)^3+2(c+2k-1)^2\\
 &\hspace{7em}+\{2+2k(1-2k)\}(c+2k-1)
 +2(1-k)(1-2k)\Bigr]>0.
\end{align*}
Each term in brackets is nonnegative, and the last is strictly positive. This proves $\E X<2k$.
\end{proof}

For $1<\lambda<m$, the assumption $\nu\ge2$ enters through~\eqref{eq:two-child-constraint}. Independence is used under the original law and conditional exchangeability under $Q$. If one child is allowed, the same calculation gives only $\E_QS\ge k$. Section~\ref{sec:high} uses this weaker inequality for $(3+\sqrt5)/2\le\lambda<m$.

\subsection{Proof by direct comparison}\label{sec:direct-comparison}
We first prove a criterion that also applies when single-child vertices are allowed. Fix $1<\lambda<m$ and use the notation of Section~\ref{sec:kernel-identities}, evaluated at this bias and escape law.

\begin{proposition}\label{prop:criterion}
If $\E X<2k$, then $\Dplus v(\lambda)<0$.
\end{proposition}

\begin{proof}
We compare a larger bias with $\lambda$ by changing the escape law and the bias separately. The escape law decreases in stochastic order. We modify $N_\alpha(\rho)-vZ_\alpha(\rho)$ by a term that vanishes when $\rho=\pi_\alpha$, so that the modified expression increases with $\rho$ near $\pi_\lambda$ and has a negative derivative in $\alpha$ at $(\lambda,\pi_\lambda)$. Lemma~\ref{lem:weighted-positivity} gives the two signs needed below.

For a nearby parameter $\alpha>1$ and a child law $\rho$ on $[0,1]$, keep $u$ and $v$ fixed and set
\begin{align}\label{eq:corrected-functional}
 \cF(\alpha,\rho)
 =N_\alpha(\rho)-vZ_\alpha(\rho)
 -2\left\{\E_\rho u(X)
 -\E_{\nu,\rho}u\left(F_\alpha\left(\sum_{i=1}^{\nu}X_i\right)\right)\right\}.
\end{align}
The expression in braces vanishes at $\rho=\pi_\alpha$ by the recursion. Consequently,
\begin{equation}\label{eq:on-curve}
 \cF(\alpha,\pi_\alpha)
 =Z_\alpha(\pi_\alpha)\{v(\alpha)-v(\lambda)\}.
\end{equation}
For an ordered coupling $X_0\le X_1$, write $X_s=(1-s)X_0+sX_1$ and $\rho_s=\Law(X_s)$. Differentiating~\eqref{eq:corrected-functional} gives
\begin{align}\label{eq:transport}
 \frac{\dd}{\dd s}\cF(\alpha,\rho_s)
 &=2\E[(X_1-X_0)G_{\alpha,\rho_s}(X_s)],\\
 G_{\alpha,\rho}(x)
 &=\int\frac{y^2+(\alpha-1)h_v(y)}{A_\alpha(x,y)^2}\rho(\dd y)
 -r(x)+(\cL_{\alpha,\rho}r)(x).\notag
\end{align}
At $(\alpha,\rho)=(\lambda,\pi_\lambda)$, the choice of $r$ gives $G_{\lambda,\pi_\lambda}=c\cJ h_v$, which is uniformly positive by~\eqref{eq:J-positive}. The map $(\alpha,\rho)\mapsto G_{\alpha,\rho}$ is continuous in the uniform norm when $\rho$ varies in $W_1$. For the pair-kernel term this follows from uniform continuity on the compact parameter and coordinate domain. For $\cL_{\alpha,\rho}r$, truncate $\nu$ at $M$. For each of the finitely many possible offspring numbers, couple the child variables so their mean absolute differences tend to zero. Uniform continuity of the finite-child integrand then gives convergence uniformly in $x\in[0,1]$. The tail of each operator, on a compact interval with $\alpha\ge a>1$, is bounded by
\begin{equation*}
 \frac{\|r\|_\infty}{a}\E[\nu;\nu>M]\longrightarrow0.
\end{equation*}
The same bound justifies the differentiation in~\eqref{eq:transport}. Hence $G_{\alpha,\rho}>0$ throughout a neighborhood of $(\lambda,\pi_\lambda)$.

We next vary $\alpha$ with the law fixed at $\pi_\lambda$. Since $\partial_\alpha F_\alpha(s)|_{\alpha=\lambda}=-b(F_\lambda(s))$, the derivative of $\cF(\alpha,\pi_\lambda)$ at $\alpha=\lambda$ equals $C-2\E[br]=-K<0$, by Lemma~\ref{lem:weighted-positivity}.

We now combine the two comparisons. Couple $\beta_{\lambda+\varepsilon}\le\beta_\lambda$ on the same tree. By Lemma~\ref{lem:continuity}, the entire linear interpolation of these variables has law converging to $\pi_\lambda$ in $W_1$, uniformly along the interpolation. Equation~\eqref{eq:transport} therefore gives, for small $\varepsilon>0$,
\begin{align*}
 Z_{\lambda+\varepsilon}(\pi_{\lambda+\varepsilon})
 \{v(\lambda+\varepsilon)-v(\lambda)\}
 &=\cF(\lambda+\varepsilon,\pi_{\lambda+\varepsilon})\\
 &\le\cF(\lambda+\varepsilon,\pi_\lambda)
 =-K\varepsilon+o(\varepsilon).
\end{align*}
Since $Z_{\lambda+\varepsilon}(\pi_{\lambda+\varepsilon})\to Z>0$, we obtain $\Dplus v(\lambda)\le-K/Z<0$.
\end{proof}

We will also use the elementary fact that a continuous function with strictly negative upper right Dini derivative at every point of an interval is strictly decreasing. Indeed, if $f(b)\ge f(a)$ for some $a<b$, a minimum on $[a,b]$ is attained at some $x<b$. All sufficiently small right difference quotients at that $x$ are nonnegative, a contradiction.

\begin{proof}[First proof of Theorem~\ref{thm:twochildren}]
For each $1<\lambda<m$, Lemma~\ref{lem:two-child} and Proposition~\ref{prop:criterion} give $\Dplus v(\lambda)<0$. Lemma~\ref{lem:continuity} and the Dini criterion above imply strict decrease on $(1,m)$.

To include $\lambda\le1$, we use a known interval of strict decrease that overlaps $(1,m)$. Let $m_1=\min\{k:\PP(\nu=k)>0\}\ge2$. Song--Wang--Xiang~\cite[Theorem~1.2 and Lemma~2.5]{SWX} prove strict decrease on
\begin{equation}\label{eq:SWX-range}
 \left[0,\frac{m_1}{1+\sqrt{1-1/m_1}}\right].
\end{equation}
The right endpoint is larger than one and smaller than $m_1\le m$. Therefore~\eqref{eq:SWX-range} overlaps $(1,m)$, and the two strict comparisons yield strict decrease on $[0,m)$.
\end{proof}

\subsection{Proof via linear response}\label{sec:two-child-response}
The second proof uses the same weighted-mean inequality to determine the sign of the derivative formula from Section~\ref{sec:linear-response}.

\begin{proof}[Second proof of Theorem~\ref{thm:twochildren}]
Fix $1<\lambda<m$. Lemma~\ref{lem:two-child} gives $\E X<2k$, so Lemma~\ref{lem:weighted-positivity} gives $\cJ h_v>0$ and $K>0$. Recall $\Psi=\cJ(y^2+c h_v)$ and the positive inverse $\mathcal R=(I-\cL)^{-1}$ from~\eqref{eq:full-response-defs} and Lemma~\ref{lem:positive-inverse}. Applying $\mathcal R$ to $(I-\cL)r=\cJ(y^2)$ yields
\begin{equation}\label{eq:two-child-response}
 \mathcal R\Psi=r+c\mathcal R\cJ h_v.
\end{equation}
Substituting this identity into~\eqref{eq:full-sign} gives
\begin{align*}
 D_\lambda
 &=2\E[b(X)r(X)]-C
   +2c\E[b(X)(\mathcal R\cJ h_v)(X)]\\
 &=K+2c\E[b(X)(\mathcal R\cJ h_v)(X)]\ge K>0.
\end{align*}
Here $c>0$, $b\ge0$ and positivity of $\mathcal R$ give the inequality. Since $Z>0$, Lemma~\ref{lem:full-response} gives $v'(\lambda)=-D_\lambda/Z<0$. The mean value theorem now gives strict decrease on $(1,m)$. Combining this interval with the overlapping Song--Wang--Xiang interval~\eqref{eq:SWX-range} proves strict decrease on $[0,m)$.
\end{proof}

The second proof gives $v'(\lambda)<0$ on $(1,m)$ under the finite first-moment assumption. If the offspring law has an exponential moment, Bowditch--Tokushige~\cite[Theorem~1.1]{BT} give differentiability of the speed throughout $(0,m)$ in the leafless case. The bound $\Dplus v(\lambda)<0$ for $1<\lambda<m$, together with~\cite[Lemma~2.5]{SWX}, then gives $v'(\lambda)<0$ for every $0<\lambda<m$. The strict monotonicity in Theorem~\ref{thm:twochildren} itself uses only the finite first moment.

\section{Strict decrease for \texorpdfstring{$0<\lambda<1$}{0 < lambda < 1}}\label{sec:low}

We now allow vertices with a single child. For $0<\lambda<1$, the escape probability is at least $1-\lambda$ on every tree. This bound allows us to compare the escape laws at two biases. We combine the comparison with a resistance estimate to show that the decrease in escape probabilities outweighs the explicit dependence on $\lambda$ in the speed formula.

\subsection{An affine bound for the escape probabilities}
\begin{lemma}\label{lem:low-transport}
For $0<\lambda<\alpha<1$, the same-tree escape probabilities satisfy
\begin{equation}\label{eq:low-transport}
 \beta_\alpha\le T_{\lambda,\alpha}(\beta_\lambda),\qquad
 T_{\lambda,\alpha}(x)
 =1-\frac{\alpha(1+\lambda)}{\lambda(1+\alpha)}(1-x).
\end{equation}
The affine map $T_{\lambda,\alpha}$ is increasing and sends $[1-\lambda,1]$ into $[1-\alpha,1]$. Moreover, $\lambda\mapsto\beta_\lambda(T)$ is locally Lipschitz on $(0,1)$, uniformly in the tree.
\end{lemma}

\begin{proof}
We first bound the derivative of the finite-depth recursion, integrate the bound, and then let the depth tend to infinity. Take terminal value one; the resulting recursion is differentiable for positive bias. Put $d=1-\lambda$ and $a_\lambda=1/[\lambda(1+\lambda)]$. At a vertex with $n\ge1$ children, write $x=S/(\lambda+S)$, where $S=\sum_{i=1}^n x_i$. If every child value is at least $d$, then $S\ge nd\ge d$ and $S/(\lambda+S)\ge d/(\lambda+d)=d$. Thus induction gives $d\le x\le1$ at every depth. With $q=-\partial_\lambda x$ and $q_i=-\partial_\lambda x_i$,
\begin{equation}\label{eq:finite-sensitivity}
 q=\frac{S+\lambda\sum_{i=1}^n q_i}{(\lambda+S)^2},
 \qquad 0\le q\le\frac{x}{d}\le\frac1d.
\end{equation}
For the upper bound, the induction hypothesis $q_i\le x_i/d$ gives
\[
 q\le\frac{S+\lambda S/d}{(\lambda+S)^2}
 =\frac{S}{d(\lambda+S)^2}
 \le\frac{x}{d},
\]
because $\lambda+S\ge\lambda+d=1$. At the terminal generation the derivative is zero, so the induction starts.

We also claim $q\ge a_\lambda(1-x)$. The claim holds at the terminal vertices. If it holds for all children, subtracting the required numerator from the numerator in~\eqref{eq:finite-sensitivity} gives
\begin{align*}
 S+\lambda a_\lambda(n-S)-a_\lambda\lambda(\lambda+S)
 &=\frac{n-\lambda-dS}{1+\lambda}\\
 &\ge\frac{\lambda(n-1)}{1+\lambda}\ge0,
\end{align*}
where $S\le n$ was used. Integration of $(1-x)'\ge a_\lambda(1-x)$ between $\lambda$ and $\alpha$ yields~\eqref{eq:low-transport} at finite depth, since
\begin{equation*}
 \exp\left(\int_\lambda^\alpha\frac{\dd s}{s(1+s)}\right)
 =\frac{\alpha(1+\lambda)}{\lambda(1+\alpha)}.
\end{equation*}
The affine map $T_{\lambda,\alpha}$ is continuous, so letting the depth tend to infinity proves the comparison for $\beta_\lambda$ and $\beta_\alpha$. On every compact interval $[a,b]\subset(0,1)$, the finite-depth derivative bound is at most $1/(1-b)$. Integrating that bound between any two biases and then passing to the depth limit gives
\[
 |\beta_{\lambda_1}(T)-\beta_{\lambda_2}(T)|
 \le\frac{|\lambda_1-\lambda_2|}{1-b},
 \qquad \lambda_1,\lambda_2\in[a,b],
\]
uniformly in the tree. Finally,
\begin{equation*}
 T_{\lambda,\alpha}(1-\lambda)
 =1-\frac{\alpha(1+\lambda)}{1+\alpha}\ge1-\alpha,
 \qquad T_{\lambda,\alpha}(1)=1,
\end{equation*}
which verifies the support assertion.
\end{proof}

\subsection{Stochastic order and the speed quotient}
Fix $0<\lambda<1$ and write $d=1-\lambda$. For a probability law $\sigma$ supported on $[d,1]$, put $N=N_\lambda(\sigma)$ and $Z=Z_\lambda(\sigma)$. Its speed quotient $v=N/Z=\mathcal V_\lambda(\sigma)$ satisfies
\begin{equation}\label{eq:low-speed-range}
 \frac{d}{1+\lambda}\le v\le1.
\end{equation}
Indeed $\mathcal V_\lambda(\sigma)$ is a weighted average of $xy/t(x,y)$, which lies between $d/(2-d)$ and one on $[d,1]^2$. With $v$ fixed at this quotient, let
\begin{equation}\label{eq:low-Psi}
 \begin{gathered}
 P_v(y)=y^2-d\{(1+v)y-v\},\\
 \Psi(x)=\E_\sigma\frac{P_v(Y)}{A_\lambda(x,Y)^2},\qquad
 C=\E_\sigma\frac{v\,t(X,Y)-XY}{A_\lambda(X,Y)^2}.
 \end{gathered}
\end{equation}
The coordinate derivative of $(xy-vt(x,y))/A_\lambda(x,y)$ is $P_v(y)/A_\lambda(x,y)^2$. Since
\begin{equation*}
 P_v(y)=y(y-d)+dv(1-y)>0,\qquad d\le y\le1,
\end{equation*}
increasing the escape variables increases the quotient. To see this, couple $X_0\le X_1$ with laws supported on $[d,1]$ and put $X_s=(1-s)X_0+sX_1$. Write $v_s=\mathcal V_\lambda(\Law(X_s))$, and let $Y_s$ be an independent copy. Quotient differentiation gives
\[
 \frac{\dd}{\dd s}\mathcal V_\lambda(\Law(X_s))
 =\frac{2}{Z_\lambda(\Law(X_s))}
 \E\left[(X_1-X_0)
 \frac{P_{v_s}(Y_s)}{A_\lambda(X_s,Y_s)^2}\right]\ge0.
\]
For each $s$, the quotient obeys~\eqref{eq:low-speed-range}, so $P_{v_s}>0$ by the displayed factorization. The positive denominator bound~\eqref{eq:low-denominator} bounds every kernel derivative on this compact square and justifies differentiation under the expectation. Thus $\mathcal V_\lambda$ preserves stochastic order among laws supported on $[1-\lambda,1]$. More generally, if $X$ varies differentiably within this interval with uniformly bounded derivative $\dot X$, the same quotient rule gives the derivative $2\E[\dot X\Psi(X)]/Z$. Changing $\lambda$ with the law fixed gives the derivative $C/Z$.

From now on, $X,Y$ are independent with law $\pi_\lambda$ at the fixed parameter $\lambda$, $N=N_\lambda(\pi_\lambda)$, $Z=Z_\lambda(\pi_\lambda)$ and $v=v(\lambda)$. All expectations in~\eqref{eq:low-Psi} refer to this law. The affine comparison in Lemma~\ref{lem:low-transport} bounds the speed at a larger bias by $\mathcal V_\alpha(\Law(T_{\lambda,\alpha}(X)))$, with equality at $\alpha=\lambda$. Its derivative there is
\[
 \frac1Z\left(C-\frac{2}{\lambda(1+\lambda)}
 \E[(1-X)\Psi(X)]\right).
\]
We will prove that this derivative is negative. Identity~\eqref{eq:low-identity} below expresses the negative of its numerator with a positive coefficient of $C$ and a nonpositive coefficient of $\E[1/A_\lambda(X,Y)]$. The next three lemmas provide the corresponding bounds, together with $v\E(1/X)\le1$.

\begin{lemma}[Speed and mean resistance]\label{lem:resistance}
Write $\rho=\E(1/X)$ and $e=\E(1/\nu)$. Then
\begin{equation}\label{eq:resistance}
 v\le1-\lambda e,\qquad
 1<\rho\le\frac1{1-\lambda e},\qquad v\rho\le1.
\end{equation}
\end{lemma}

\begin{proof}
To bound $v$, we write~\eqref{eq:degree-speed} as a weighted average of local drifts. Let $T_j=\sum_{i=1}^jX_i$ for independent $X_i$ with law $\pi_\lambda$. Conditional on $\nu=n$, the weight in the denominator is
\begin{equation*}
 W_{n+1},\qquad
 W_j=\frac{j-d}{j}\E\frac{T_j}{T_j-d}
 =1+d\E\frac{1-X_1}{T_j-d},\qquad j\ge2.
\end{equation*}
The identity for $W_j$ follows by exchangeability: $\E[X_1/(T_j-d)]=j^{-1}\E[T_j/(T_j-d)]$. Coupling all $T_j$ by one independent sequence shows that $W_j$ decreases with $j$, since the nonnegative numerator $1-X_1$ stays fixed while the denominator increases. The degree formula is exactly
\[
 v=\frac{\E\left[W_{\nu+1}(\nu-\lambda)/(\nu+\lambda)\right]}
 {\E W_{\nu+1}}.
\]
The local drift $(n-\lambda)/(n+\lambda)$ increases with $n$, whereas $W_{n+1}$ decreases. Their covariance is nonpositive. All terms are integrable: $T_j-d\ge(j-1)d$ and the local drift is bounded. Consequently,
\begin{equation*}
 v\le\E\frac{\nu-\lambda}{\nu+\lambda}
 =1-2\lambda\E\frac1{\nu+\lambda}
 \le1-\lambda\E\frac1\nu.
\end{equation*}
For the last inequality, $\lambda<1\le\nu$ implies $2/(\nu+\lambda)\ge1/\nu$.

The resistance $R=1/X$ is bounded by $1/d$ and satisfies
\begin{equation*}
 R\overset{\mathrm d}=1+\frac{\lambda}{\sum_{i=1}^{\nu}R_i^{-1}}.
\end{equation*}
The harmonic-mean inequality gives $(\sum_{i=1}^{\nu}R_i^{-1})^{-1}\le\nu^{-2}\sum_{i=1}^{\nu}R_i$. Conditional on $\nu=n$, independence gives the expectation of the right-hand side as $\rho/n$. Thus $\rho\le1+\lambda e\rho$. Since $\lambda e<1$, this gives $\rho\le(1-\lambda e)^{-1}$. Both $v$ and $\rho$ are positive, and their respective upper bounds give $v\rho\le(1-\lambda e)/(1-\lambda e)=1$. Strictness of $\rho>1$ follows from $X<1$ almost surely.
\end{proof}

\begin{lemma}[Sign of the explicit parameter term]\label{lem:Cpositive}
The quantity $C$ in~\eqref{eq:low-Psi}, evaluated under $\pi_\lambda$, is nonnegative.
\end{lemma}

\begin{proof}
We expand $C$ and the identity $vZ-N=0$ in the same powers of $\lambda^{-1}$. Their coefficients will express $C$ as a covariance of two increasing sequences. Put $a_j=\E(1-X)^j$ and $b_j=a_j-a_{j+1}$ for $j\ge0$. Since $1-X\le\lambda$, the geometric expansion of $A_\lambda^{-1}=[\lambda-(1-X)(1-Y)]^{-1}$ gives absolutely convergent series
\begin{align*}
 N&=\lambda^{-1}\sum_{j\ge0}b_j^2\lambda^{-j},\\
 Z&=\lambda^{-1}\sum_{j\ge0}b_j(2a_j-b_j)\lambda^{-j},\\
 C&=\lambda^{-2}\sum_{j\ge0}(j+1)
 \{2va_jb_j-(1+v)b_j^2\}\lambda^{-j}.
\end{align*}
For example $a_j,b_j\le\lambda^j$, so the series for $C$ is dominated by a constant times $\sum(j+1)\lambda^j$. The expansion follows from
\[
 \frac1{A_\lambda(X,Y)}=\frac1\lambda
 \sum_{j\ge0}\left(\frac{(1-X)(1-Y)}\lambda\right)^j,
 \quad b_j=\E[X(1-X)^j].
\]
The expansion of $A_\lambda^{-2}$ has the additional factors $j+1$. Since $0<X<1$, all $a_j,b_j$ are positive. Cauchy--Schwarz gives $a_{j+1}^2\le a_ja_{j+2}$, so $a_{j+1}/a_j$ increases and $b_j/a_j=1-a_{j+1}/a_j$ decreases. Therefore $2v-(1+v)b_j/a_j$ increases with $j$. Its weighted sum is
\[
 \sum_{j\ge0}a_jb_j\lambda^{-j}
 \left(2v-(1+v)\frac{b_j}{a_j}\right)
 =\lambda(vZ-N)=0.
\]
Normalize the summable positive weights $a_jb_j\lambda^{-j}$ to a probability distribution on the nonnegative integers. The covariance of the increasing sequences $j+1$ and $2v-(1+v)b_j/a_j$ is nonnegative. The zero weighted mean shows that the series for $C$ equals a positive constant times this covariance. Its first moment is finite by the geometric bound already proved, so $C\ge0$.
\end{proof}

\begin{lemma}[Reciprocal kernel estimate]\label{lem:reciprocal}
With $\rho=\E(1/X)$, set $B_\rho=2\rho-1-d\rho^2$. Then $B_\rho>0$ and
\begin{equation}\label{eq:reciprocal}
 K:=\E\frac1{A_\lambda(X,Y)}\le\frac{\rho^2}{B_\rho}.
\end{equation}
\end{lemma}

\begin{proof}
For $R,S\in[1,1/d]$, the reciprocal kernel is
\begin{equation*}
 \frac1{A_\lambda(1/R,1/S)}=\frac{RS}{R+S-1-dRS}.
\end{equation*}
Its second derivative in $R$ equals
\begin{equation*}
 -\frac{2S(S-1)(1-dS)}{(R+S-1-dRS)^3}\le0,
\end{equation*}
and the same holds in $S$. All denominators are positive on $[1,1/d]^2$, since $R+S-1-dRS=RS A_\lambda(1/R,1/S)$. Apply Jensen's inequality first in $R=1/X$, conditional on $S=1/Y$, and then in $S$:
\[
 K\le\E\frac{\rho S}{\rho+S-1-d\rho S}
 \le\frac{\rho^2}{2\rho-1-d\rho^2}.
\]
Independence makes the conditional mean of $R$ equal to $\rho$. Since $\rho\in[1,1/d]$, the same denominator identity gives $B_\rho=\rho^2A_\lambda(1/\rho,1/\rho)>0$.
\end{proof}

\subsection{A negative upper bound for \texorpdfstring{$\Dplus v(\lambda)$}{the upper right Dini derivative}}
\begin{proof}[Proof of Theorem~\ref{thm:low}]
For $\alpha>\lambda$ close to $\lambda$, Lemma~\ref{lem:low-transport} and stochastic monotonicity of $\mathcal V_\alpha$ on $[1-\alpha,1]$ give
\begin{equation}\label{eq:low-upper-comparison}
 v(\alpha)\le
 \mathcal V_\alpha\bigl(\Law(T_{\lambda,\alpha}(\beta_\lambda))\bigr).
\end{equation}
Equality holds at $\alpha=\lambda$. We will show that the right-hand side has a strictly negative derivative there.


At $\alpha=\lambda$, the derivative of $T_{\lambda,\alpha}(X)$ with respect to $\alpha$ is $-a_\lambda(1-X)$ by Lemma~\ref{lem:low-transport}, where $a_\lambda=1/[\lambda(1+\lambda)]$. By the derivative formulas following~\eqref{eq:low-Psi}, the derivative of the upper bound in~\eqref{eq:low-upper-comparison} is $-D_a/Z$, where
\begin{equation}\label{eq:low-margin}
 D_a=2a_\lambda\E[(1-X)\Psi(X)]-C.
\end{equation}
We prove $D_a>0$ using the three preceding lemmas. Write $\Delta=(1+\lambda)v-d\ge0$, where nonnegativity follows from~\eqref{eq:low-speed-range}. The exact identity is
\begin{equation}\label{eq:low-identity}
 \lambda(1+\lambda)D_a=(1+v)-\Delta K+d\lambda C.
\end{equation}
To verify~\eqref{eq:low-identity}, expand
\begin{align*}
 &(1-x)P_v(y)+(1-y)P_v(x)-\lambda(1+\lambda)\{vt(x,y)-xy\}\\
 &\qquad=A_\lambda(x,y)\{x+y-(2-d+d^2)v\}-d\lambda(xy-dv).
\end{align*}
Divide by $A_\lambda(x,y)^2$ and average. Using $xy-dv=xy-vt(x,y)+vA_\lambda(x,y)$ and $(2-d+d^2)+d\lambda=2$, we obtain
\begin{align*}
 \lambda(1+\lambda)D_a
 &=\E\frac{X+Y-2v}{A_\lambda(X,Y)}+d\lambda C\\
 &=(1+v)(1+dK)-2vK+d\lambda C\\
 &=(1+v)-\Delta K+d\lambda C.
\end{align*}
The second equality uses $x+y=t(x,y)+xy$, $N=vZ$ and $Z=1+dK$. This proves~\eqref{eq:low-identity}.

Lemmas~\ref{lem:Cpositive} and~\ref{lem:reciprocal} imply
\begin{equation*}
 \lambda(1+\lambda)D_a\ge1+v-\Delta\frac{\rho^2}{B_\rho}.
\end{equation*}
The right-hand side decreases as a function of $v$, since
\begin{equation*}
 B_\rho-(1+\lambda)\rho^2=-(\rho-1)^2-\rho^2<0.
\end{equation*}
Using $v\le1/\rho$ from Lemma~\ref{lem:resistance} gives
\begin{equation}\label{eq:low-strict}
 D_a\ge\frac{1-1/\rho}{\lambda(1+\lambda)B_\rho}>0.
\end{equation}

It remains to justify differentiating the right-hand side of~\eqref{eq:low-upper-comparison}. On a sufficiently small right interval $[\lambda,\lambda+\delta]\subset(0,1)$, the variables $T_{\lambda,\alpha}(\beta_\lambda)$ lie in $[1-\alpha,1]$, so~\eqref{eq:low-denominator} bounds the denominators away from zero. Both $T_{\lambda,\alpha}(\beta_\lambda)$ and its derivative in $\alpha$ are uniformly bounded there. Differentiation under the expectation is therefore legitimate, and the derivative of the upper bound at $\alpha=\lambda$ is $-D_a/Z$, by~\eqref{eq:low-Psi} and~\eqref{eq:low-margin}. Subtracting the common value at $\alpha=\lambda$ in~\eqref{eq:low-upper-comparison}, dividing by $\alpha-\lambda$, and taking the upper limit gives $\Dplus v(\lambda)\le-D_a/Z<0$. The local Lipschitz bound in Lemma~\ref{lem:low-transport} is uniform over trees. Together with the same denominator bound and the speed formula, it gives continuity of $v$ on $(0,1)$. The Dini criterion in Section~\ref{sec:direct-comparison} completes the proof.
\end{proof}

\section{Strict decrease for \texorpdfstring{$(3+\sqrt5)/2\le\lambda<m$}{(3+sqrt(5))/2 <= lambda < m}}\label{sec:high}

When $\PP(\nu=1)>0$, the change of measure from Section~\ref{sec:twochildren} gives $\E_QS\ge k$ in place of $\E_QS\ge2k$. We show that this still implies $\E X<2k$ for $\lambda\ge(3+\sqrt5)/2$, so Proposition~\ref{prop:criterion} applies.

\begin{lemma}\label{lem:one-child}
Let $\lambda\ge(3+\sqrt5)/2$, $c=\lambda-1$, and let $X\in(0,1)$ satisfy~\eqref{eq:rde} with $\nu\ge1$. Then the weighted mean~\eqref{eq:weighted-mean} satisfies $\E X<2k$.
\end{lemma}

\begin{proof}
Use the probability measure $M$ obtained by weighting the law of $X$ with $(c+X)^{-2}$, and let $k=\E_M X$, $j=\E_M X^2$, $\ell=\E_M X^3$. As in Lemma~\ref{lem:two-child}, weight the joint law of $\nu,X_1,X_2,\ldots$ by $(c+S)^{-2}$ to obtain $Q$. These normalizing constants are finite and positive because $c>0$ and $X,S>0$.

The comparison of marginal densities in Lemma~\ref{lem:two-child} still applies when $\nu=1$: the remainder $H=\sum_{i=2}^{\nu}X_i$, which is independent of $X_1$ under the original law, is then zero, and $((c+x)/(c+x+H))^2$ is the constant one. For all $H\ge0$ this ratio is nondecreasing in $x$. Thus $\E_QX_1\ge k$. Conditional exchangeability gives $\E_QS=\E_Q[\nu X_1]\ge\E_QX_1\ge k$. The recursion transforms this expectation into
\[
 \E_QS=\frac{\E[\lambda X(1-X)/(c+X)^2]}
 {\E[(1-X)^2/(c+X)^2]}
 =\lambda\frac{k-j}{1-2k+j}.
\]
Here $1-2k+j=\E_M(1-X)^2>0$. Rearranging $\E_QS\ge k$, with $\lambda=c+1$, gives
\begin{equation}\label{eq:one-j-bound}
 \lambda\frac{k-j}{1-2k+j}\ge k,
 \qquad j\le\frac{k(c+2k)}{c+1+k}.
\end{equation}
Undoing the change of measure gives exactly~\eqref{eq:untilted-mean}, with positive denominator $c^2+2ck+j=\E_M(c+X)^2$. If $k\ge1/2$, then $\E X<1\le2k$. Suppose $k<1/2$. Since $0<X<1$, we have $\ell\le j$, and therefore
\[
 (2k-\E X)(c^2+2ck+j)
 =c^2k+4ck^2+(2k-2c)j-\ell
 \ge c^2k+4ck^2+(2k-2c-1)j.
\]
The coefficient $2k-2c-1$ is negative. Substituting the upper bound for $j$ from~\eqref{eq:one-j-bound} consequently gives the lower bound
\begin{align}\label{eq:one-child-polynomial}
 &(2k-\E X)(c^2+2ck+j)\notag\\
 &\quad\ge c^2k+4ck^2+(2k-2c-1)\frac{k(c+2k)}{c+1+k}\notag\\
 &\quad=\frac{k}{c+1+k}
 \left[c(c^2-c-1)+(5c^2+2c-2)k+4(c+1)k^2\right]>0.
\end{align}
Indeed $c\ge(1+\sqrt5)/2$, so $c^2-c-1\ge0$ and $5c^2+2c-2>0$. Since $k>0$, strictness holds even at the threshold.
\end{proof}

\begin{proof}[Proof of Theorem~\ref{thm:high}]
For each $(3+\sqrt5)/2\le\lambda<m$, Lemma~\ref{lem:one-child} verifies the hypothesis of Proposition~\ref{prop:criterion}. Thus $\Dplus v(\lambda)<0$, including at the left endpoint. Continuity and the Dini criterion in Section~\ref{sec:direct-comparison} give strict decrease on $[(3+\sqrt5)/2,m)$.
\end{proof}

For $1<\lambda<(3+\sqrt5)/2$, the constant term in the brackets in~\eqref{eq:one-child-polynomial} is negative, so this estimate no longer proves $\E X<2k$. The computer-assisted argument in Section~\ref{sec:computer} treats the remaining interval.

\clearpage
\section{Computer-assisted proof for \texorpdfstring{$1<\lambda<m$}{1 < lambda < m}}\label{sec:computer}

We prove that $v'(\lambda)<0$ for $1<\lambda<m$ when the offspring law allows a single child. All bias intervals below are restricted to $\lambda<m$. We derive bounds for the infinite series and their tails, then reduce the sign of the derivative to finitely many rational inequalities. The coefficient tables and verification programs accompany the paper~\cite{Supplement}.

\begin{theorem}[Computer-assisted monotonicity]\label{thm:computer}
Fix an offspring law $\mu$ with $\mu(0)=0$, $\mu(1)>0$ and $1<m=\E_\mu\nu<\infty$. With the root and transition conventions of Section~\ref{sec:introduction},
\begin{equation}\label{eq:computer-result}
 v_\mu\in C^1((1,m)),\qquad v_\mu'(\lambda)<0\quad(1<\lambda<m).
\end{equation}
\end{theorem}

Fix $1<\lambda<m$, put $c=\lambda-1$, and let $X,Y$ be independent with law $\pi_\lambda$. We retain $v=v(\lambda)$ and the notation of Sections~\ref{sec:kernel-identities}--\ref{sec:linear-response}. Lemma~\ref{lem:full-response} gives $v'(\lambda)=-D_\lambda/Z$ with $Z>0$. The escape recursion gives moment inequalities and an identity for a logarithmic function. Together they yield a lower bound for $D_\lambda$ and specify when further estimates are needed. We prove $D_\lambda>0$ on four intervals covering $1<\lambda<(3+\sqrt5)/2$. Section~\ref{sec:high} supplies the sign of the derivative for $(3+\sqrt5)/2\le\lambda<m$.

\subsection{Moment inequalities from the escape recursion}
To estimate $D_\lambda$, we compare the moments of $X/(c+X)$ with an explicit sequence. This comparison gives a positive lower bound for $\Psi$ and bounds on the tails of the series used below.

Put
\begin{equation}\label{eq:archive-parameters}
 \begin{gathered}
 W=\frac1{c+X},\quad w=\E W,\quad h=w^{-1}-c,\quad d=1-h,\\
 U=\frac{X}{c+X},\qquad u=\E U=\frac{h}{c+h}.
 \end{gathered}
\end{equation}
Here $0<h<1$ and $0<U<1/\lambda$ almost surely. In these moment estimates, $d$ denotes $1-h$ and $u$ denotes $\E U$. In Section~\ref{sec:low}, $d$ denoted $1-\lambda$, and in Section~\ref{sec:preliminaries}, $u$ denoted a primitive. Define
\begin{equation}\label{eq:canonical-moment-sequences}
 a_n=\E U^n,\qquad \beta_n=\frac{h}{\lambda^n-d},\qquad
 \rho_n=\frac{a_n}{\beta_n},\qquad n\ge0.
\end{equation}
The sequence $\beta_n$ in~\eqref{eq:canonical-moment-sequences} consists of deterministic comparison moments. We continue to write $\beta_\lambda$ for the random escape probability.

\begin{lemma}[Monotonicity of the moment ratios]\label{lem:canonical-moments}
For the fixed offspring law satisfying~\eqref{eq:assumptions},
\begin{equation}\label{eq:rho-monotone}
 \rho_0=\rho_1=1,\qquad
 \rho_n>\rho_{n+1}>0\quad(n\ge1),\qquad \rho_n\longrightarrow0.
\end{equation}
In particular, $a_n/a_{n-1}<\beta_n/\beta_{n-1}$ for $n\ge2$.
\end{lemma}
\begin{proof}
Let $S=\sum_{i=1}^{\nu}X_i$ and $T=\sum_{i=2}^{\nu}X_i$. The first child is independent of $T$, including the random number of summands. Write
\[
 H=\frac1{c+S}=\frac{W_1}{1+W_1T}.
\]
The recursion gives $\lambda W\overset{\mathrm d}=1+H$. Consequently,
$\E H/w=\lambda-1/w=d$. Since $0<X<1$, we have $1/\lambda<W<1/c$, and thus $0<h<1$.

For a bounded increasing function $f$ on $[0,1]$, $H\le W_1$ gives
\[
 \E[Hf(1-cH)]\ge\E[Hf(1-cW_1)].
\]
Condition on $W_1$. The multiplier $M(a)=\E(1+aT)^{-1}$ decreases in $a$. Under the probability measure that weights the law of $W_1$ by $W_1/w$, both $M(W_1)$ and $f(1-cW_1)$ are decreasing functions of $W_1$. Their covariance is nonnegative. Restoring the original measure gives
\begin{equation}\label{eq:chosen-child-comparison}
 \E[Hf(1-cH)]\ge
 \frac{\E H}{w}\E[Wf(1-cW)]
 =d\E[Wf(U)].
\end{equation}
No change of law for $\nu$ is made here. Since $m>1$ and $\nu\ge1$, the event $\nu\ge2$ has positive probability. On that event $T>0$, so the first inequality is strict for $f(t)=t^n$, $n\ge1$.

Using $H\overset{\mathrm d}=\lambda W-1$ and $1-cH\overset{\mathrm d}=\lambda U$ in~\eqref{eq:chosen-child-comparison}, and multiplying by $c$, yields
\[
 (\lambda^n-d)a_n-(\lambda^{n+1}-d)a_{n+1}>0,
 \qquad n\ge1.
\]
Division by $h$ proves the strict decrease in~\eqref{eq:rho-monotone}; $a_0=\beta_0=1$ and $a_1=\beta_1=h/(c+h)$ prove its initial values. Finally,
\[
 \rho_n=\frac{\E[(\lambda U)^n]-d\E U^n}{h}\longrightarrow0,
\]
by bounded convergence, since $0<\lambda U<1$ almost surely. This proves the limit without a deterministic escape bound strictly smaller than one.
\end{proof}

\begin{lemma}[An upper bound for the speed and a lower bound for the kernel]\label{lem:harmonic-speed}
For $X$ with law $\pi_\lambda$, one has $v<v_p<h$, where $v_p$ is the symmetric speed quotient of the comparison law whose $U$-moments are $\beta_n$. Moreover, $\inf_{0\le x\le1}\Psi(x)>0$.
\end{lemma}
\begin{proof}
Let $V$ be an independent copy of $U$. Substituting $X=cU/(1-U)$ and $Y=cV/(1-V)$ into the pair kernels gives
\[
 \frac{XY}{A(X,Y)}=\frac{cUV}{1-\lambda UV},\qquad
 \frac{t(X,Y)}{A(X,Y)}=
 \frac{U+V-(\lambda+1)UV}{1-\lambda UV}.
\]
Since $\lambda UV\le1/\lambda<1$, geometric expansion and absolute convergence give
\begin{equation}\label{eq:moment-speed-expansion}
 N=cT,\qquad Z=2Q-(\lambda+1)T,\qquad
 T=\sum_{n\ge1}\lambda^{n-1}a_n^2,\quad
 Q=\sum_{n\ge1}\lambda^{n-1}a_na_{n-1}.
\end{equation}
For example, the summands in $T$ are at most $\lambda^{-n-1}$ and those in $Q$ at most $\lambda^{-n}$. Thus
\[
 v=\frac{c}{2Q/T-(\lambda+1)},\qquad
 \frac QT=
 \frac{\sum_{n\ge1}\lambda^{n-1}a_n^2(a_{n-1}/a_n)}
 {\sum_{n\ge1}\lambda^{n-1}a_n^2}.
\]

The sequence $\beta_n$ is realized by $U_p=\lambda^{-L}$, where
$\PP(L=j)=h d^{j-1}$ for $j\ge1$. Its corresponding escape coordinate is $cU_p/(1-U_p)\in(0,1]$, so its symmetric denominator is positive. Write $T_p,Q_p$ for the sums in~\eqref{eq:moment-speed-expansion} with $a_n$ replaced by $\beta_n$, and $v_p=cT_p/[2Q_p-(\lambda+1)T_p]$. The ratios
\[
 \frac{\beta_{n-1}}{\beta_n}
 =\lambda+\frac{dc}{\lambda^{n-1}-d}
\]
decrease with $n$. Lemma~\ref{lem:canonical-moments} implies $a_{n-1}/a_n\ge\beta_{n-1}/\beta_n$, strictly for $n\ge2$. The weights $\lambda^{n-1}a_n^2$ are the comparison weights $\lambda^{n-1}\beta_n^2$ multiplied by the decreasing sequence $\rho_n^2$. The covariance inequality for two decreasing sequences therefore gives
\[
 \frac QT>
 \frac{\sum_{n\ge1}\lambda^{n-1}\beta_n^2\rho_n^2
       (\beta_{n-1}/\beta_n)}
 {\sum_{n\ge1}\lambda^{n-1}\beta_n^2\rho_n^2}
 \ge\frac{Q_p}{T_p}.
\]
All sums are absolutely convergent; the covariance inequality follows first for finite truncations and then by passage to the limit. Positivity of both quotient denominators gives $v<v_p$.

To prove $v_p<h$, put
\[
 Q_* =\sum_{n\ge1}\frac{\lambda^{n-1}}{(\lambda^n-d)^2},\qquad
 R_* =\sum_{n\ge1}
 \frac{\lambda^{n-1}}{(\lambda^n-d)^2(\lambda^{n-1}-d)}.
\]
Then $Q_p/T_p=\lambda+dcR_*/Q_*$ and $v_p=1/(1+2dR_*/Q_*)$. Set $s_n=(\lambda^{n-1}-d)^{-1}$. The identity
$s_n-s_{n+1}=c\lambda^{n-1}s_ns_{n+1}$ yields
\begin{align*}
 2hR_*-Q_*
 &=\frac1c\sum_{n\ge1}(s_n-s_{n+1})s_{n+1}(2h-s_n^{-1})\\
 &=\frac1c\sum_{n\ge1}(s_n-s_{n+1})^2(s_n^{-1}-h)>0.
\end{align*}
To verify the second equality, subtract its summands: the difference is
$h(s_n^2-s_{n+1}^2)-(s_n-s_{n+1})$. The series telescopes to $hs_1^2-s_1=0$, because $s_1=1/h$ and $s_n\to0$. Each term in the final series is nonnegative and every term with $n\ge2$ is positive. Since $d=1-h>0$, the inequality $2hR_*>Q_*$ is equivalent to $v_p<h$.

It remains to deduce positivity of $\Psi$. With $z=\E(c+X)^{-2}$, expansion at $x=0$ gives
\begin{equation}\label{eq:Psi-zero-positive}
 \Psi(0)=1-c(1-v)w-c\lambda vz
 >1-(c+v)w=(h-v)w>0.
\end{equation}
Here $cz<w$, since $X>0$, and $v>0$. The polynomial $y^2+c((1+v)y-v)$ is increasing on $[0,1]$. Relative to the law weighted by $(c+y)^{-2}$, the law weighted by $A(x,y)^{-2}$ has increasing density proportional to $((c+y)/A(x,y))^2$. Its expectation of this polynomial is at least $\Psi(0)/z$. Since $\cJ1(x)\ge1/\lambda^2$,
\[
 \Psi(x)\ge\frac{\Psi(0)}{\lambda^2z}>0\qquad(0\le x\le1).
\]
\end{proof}

\subsection{An identity for \texorpdfstring{$g$}{g} and a lower bound for \texorpdfstring{$D_\lambda$}{D(lambda)}}
To bound $D_\lambda$ while retaining contributions from all powers of $\cL$ in the series for $\mathcal R$, we use the function
\begin{equation}\label{eq:entropy-function}
 g(x)=\frac{x(c+x)\log(1+c/x)}{\lambda c\log\lambda}-\frac{x}{c}
 \quad(x>0),\qquad g(0)=0.
\end{equation}
The formula defines $g$ for all positive $x$, since it will also be applied to a sum of child escapes. On $[0,1]$ it is continuous, nonnegative and bounded.

\begin{lemma}[An identity for $g$ under the escape recursion]\label{lem:entropy-response}
For a finite positive child list $x_1,\ldots,x_n$, with $S=\sum_i x_i$,
\begin{equation}\label{eq:entropy-pair}
 \sum_i g(x_i)-g(S)
 =\frac1{\lambda c\log\lambda}\int_0^c(c-a)
 \sum_{i<j}\frac{x_ix_j(2a+x_i+x_j)}
 {(a+S)(a+x_i)(a+x_j)}\dd a.
\end{equation}
For the random child escapes $X_1,\ldots,X_\nu$, define
\[
 \delta_g=F_\lambda'(S)\left(\sum_{i=1}^{\nu}g(X_i)-g(S)\right).
\]
This random variable is integrable, nonnegative, and positive on $\{\nu\ge2\}$. For every $f\in C([0,1])$,
\begin{equation}\label{eq:entropy-stationarity}
 \E[g(X)(I-\cL)f(X)]
 =\E[b(X)f(X)]-\E[\delta_g f(F_\lambda(S))].
\end{equation}
Consequently the positive measure $\eta(f)=\E[b(X)(\mathcal Rf)(X)]$ satisfies
\begin{equation}\label{eq:response-lower-measures}
 \eta\ge g\pi_\lambda,\qquad \eta\ge b\pi_\lambda,
 \qquad \eta\ge\max(g,b)\pi_\lambda.
\end{equation}
\end{lemma}
\begin{proof}
Polynomial division under the integral gives
\[
 x(c+x)\log(1+c/x)
 =cx+\frac{c^2}{2}-\int_0^c\frac{a(c-a)}{a+x}\dd a.
\]
Subtract the value at $S$ from the sum over the children. The linear terms cancel. The remaining integrand is simplified by
\begin{align*}
 \frac{n-1}{a}-\sum_i\frac1{a+x_i}+\frac1{a+S}
 &=\frac1a\left(\sum_i\frac{x_i}{a+x_i}-\frac{S}{a+S}\right)\\
 &=\frac1a\sum_{i<j}
 \frac{x_ix_j(2a+x_i+x_j)}{(a+S)(a+x_i)(a+x_j)}.
\end{align*}
Multiplication by $a(c-a)$ proves~\eqref{eq:entropy-pair}; the simplified integrand is integrable at $a=0$. Differentiating the one-variable integral identity gives
\[
 g''(x)=-\frac2{\lambda c\log\lambda}
 \int_0^c\frac{a(c-a)}{(a+x)^3}\dd a<0.
\]
Together with $g(0)=g(1)=0$, this proves the claimed properties on $[0,1]$. Strict positivity of the summands in~\eqref{eq:entropy-pair} gives the assertions about the sign of $\delta_g$.

Direct substitution in~\eqref{eq:entropy-function} gives the identity for a single child,
\[
 g(F_\lambda(s))=b(F_\lambda(s))+F_\lambda'(s)g(s),\qquad s>0.
\]
For example, $c+F_\lambda(s)=\lambda(c+s)/(\lambda+s)$ and
$1+c/F_\lambda(s)=\lambda(1+c/s)$ separate the logarithm into the two required terms. Thus
\[
 g(F_\lambda(S))
 =b(F_\lambda(S))+F_\lambda'(S)\sum_i g(X_i)-\delta_g.
\]
Multiply by $f(F_\lambda(S))$, take expectations, and use the fixed-point equation and conditional exchangeability. The middle term becomes $\E[g(X)(\cL f)(X)]$, proving~\eqref{eq:entropy-stationarity}.

These expectations are justified for unbounded $\nu$. If $M=\max_{[0,1]}g$, then $g(S)\ge-S/c$ and $S\le\nu$, so
\begin{equation}\label{eq:entropy-integrability}
 0\le\delta_g\le\frac{M+1/c}{\lambda}\nu.
\end{equation}
The exchangeability term is at most $M\|f\|_\infty\nu/\lambda$ in absolute value. Both bounds are integrable. No derivative of $g$ at zero is used as a test in Lemma~\ref{lem:full-response}.

Apply~\eqref{eq:entropy-stationarity} with $\mathcal R f$ in place of $f$. For nonnegative continuous $f$, positivity of the resolvent gives
\[
 \eta(f)=\E[g(X)f(X)]+\E[\delta_g(\mathcal Rf)(F_\lambda(S))]
 \ge\E[g(X)f(X)].
\]
The Neumann series separately gives $\mathcal R f\ge f$, hence $\eta\ge b\pi_\lambda$. These inequalities hold as inequalities of finite Borel measures, by continuous approximation on the compact interval. Splitting a Borel set into its intersections with $\{g\ge b\}$ and $\{g<b\}$ proves the last assertion in~\eqref{eq:response-lower-measures}.
\end{proof}

\begin{lemma}[A lower bound for $D_\lambda$]\label{lem:calibration}
Set
\begin{equation}\label{eq:calibrated-floor}
 \begin{gathered}
 \alpha=\min\left\{1,\frac{\Psi(0)}{\cJ(y^2)(0)}\right\}>0,
 \qquad H_* =\Psi-\alpha\cJ(y^2),\\
 F_* =2\alpha\E[br]-C+2\E[gH_*].
 \end{gathered}
\end{equation}
Then $H_*\ge0$ on $[0,1]$ and
\begin{equation}\label{eq:floor-remainder}
 D_\lambda=F_*+2\E[\delta_g(\mathcal RH_*)(F_\lambda(S))]\ge F_*.
\end{equation}
If $\alpha=1$, then $D_\lambda>0$. In the remaining case $0<\alpha<1$,
\begin{equation}\label{eq:hard-calibration}
 \begin{gathered}
 A_0=1-(\lambda+1)u+\lambda a_2>0,\qquad
 \alpha=\frac{u-(v/c)A_0}{a_2}>\lambda^2u,\\
 u<\lambda^{-2},\qquad h<\frac1{\lambda+1}.
 \end{gathered}
\end{equation}
In particular, $(\lambda+1)h\ge1$ implies $\alpha=1$ and hence $D_\lambda>0$.
\end{lemma}
\begin{proof}
The polynomial $(1-\alpha)y^2+c((1+v)y-v)$ is increasing on $[0,1]$. Its inverse-square weighted expectation at zero is nonnegative by the choice of $\alpha$. The increasing-density comparison used in Lemma~\ref{lem:harmonic-speed} therefore gives $H_*\ge0$ everywhere. Since $(I-\cL)r=\cJ(y^2)$, we have
\[
 \mathcal R\Psi=\alpha r+\mathcal RH_*.
\]
Apply~\eqref{eq:entropy-stationarity} with $f=\mathcal RH_*$ and insert this decomposition into~\eqref{eq:full-sign}. This proves~\eqref{eq:floor-remainder}.

If $\alpha=1$, then $\cJ h_v(0)\ge0$, and the same increasing-density comparison gives $\cJ h_v(x)\ge0$ for all $x$. Identity~\eqref{eq:positive-K} remains valid without assuming the strict weighted-mean inequality. Its second term is strictly positive because $0<X,Y<1$, so $2\E[br]-C>0$. Formula~\eqref{eq:floor-remainder} and $gH_*\ge0$ prove $D_\lambda>0$.

Suppose $\alpha<1$. The identities
\[
 \cJ(y^2)(0)=a_2,\qquad
 \Psi(0)=u-\frac vc\{1-(\lambda+1)u+\lambda a_2\}
\]
give the formula for $\alpha$. Also $A_0=\E[(1-U)(1-\lambda U)]>0$. Since $v<h=cu/(1-u)$,
\[
 \alpha>
 \frac{u-uA_0/(1-u)}{a_2}
 =\frac{\lambda u(u-a_2)}{(1-u)a_2}.
\]
Lemma~\ref{lem:canonical-moments} gives
$a_2<\beta_2=u/(\lambda+1-\lambda u)$. Hence
$(u-a_2)/[(1-u)a_2]>\lambda$, proving $\alpha>\lambda^2u$. Because $\alpha<1$, this implies $u<\lambda^{-2}$, and $h=cu/(1-u)<c/(\lambda^2-1)=1/(\lambda+1)$.
\end{proof}

At $v=N/Z$ we also have the identity
\begin{equation}\label{eq:C-kernel-identity}
 C=\E\left[b(X)\cJ\bigl((1-y)(2v-(1+v)y)\bigr)(X)\right].
\end{equation}
To check it without a sign assumption, put $q=xy-vt(x,y)$. Direct polynomial expansion gives
\[
 x(1-x)(1-y)(2v-(1+v)y)
 +y(1-y)(1-x)(2v-(1+v)x)
 =-2(1-x)(1-y)q.
\]
After averaging and using symmetry, the right-hand side of~\eqref{eq:C-kernel-identity} is
$-\lambda^{-1}\E[(1-X)(1-Y)q/A^2]$.
Since $(1-x)(1-y)=\lambda-A(x,y)$ and $\E[q/A]=N-vZ=0$, this is $-\E[q/A^2]=C$. If $v$ is replaced by an upper bound with the escape law fixed, $N-vZ$ need no longer vanish. We retain this term in the next subsection.

\subsection{Integral formulas for moment ratios and a series lower bound}
\label{subsec:marker-floor}
The lower bound $F_*$ in Lemma~\ref{lem:calibration} contains the
logarithmic function $g$. We bound its contribution using the moments
of $U$ and the successive differences of the ratios $a_s/\beta_s$.
An integral representation relates these differences at all orders.

For real $s\ge0$ retain
\[
 a_s=\E U^s,\qquad \beta_s=\frac{h}{\lambda^s-d},\qquad
 \rho(s)=\frac{a_s}{\beta_s}.
\]
For $s\ge1$ put
\begin{equation}\label{eq:marker-deficits}
 \begin{gathered}
 r_s=\frac{\beta_{s+1}}{\beta_s},\qquad
 \Delta_s=\rho(s)-\rho(s+1),\qquad
 \delta_s=a_s r_s-a_{s+1}=\beta_{s+1}\Delta_s,\\
 r_0=u,\qquad \delta_0=0,\qquad
 k=\frac{c}{\lambda\log\lambda}.
 \end{gathered}
\end{equation}
In particular, $1/\lambda<k<1$.

\begin{lemma}[Integral representation of the moment ratios]\label{lem:canonical-marker}
There is a probability measure $\sigma$ on $(0,1)$, constructed from the
escape law $\pi_\lambda$, such that, for every real $s\ge1$,
\begin{equation}\label{eq:marker-representation}
 \rho(s)=\int_0^1 t^{s-1}\{s-(s-1)t\}\,\sigma(\dd t),
 \qquad
 \Delta_s=s\int_0^1 t^{s-1}(1-t)^2\,\sigma(\dd t).
\end{equation}
The derivative $\rho'(s)$ exists, with the right derivative taken at $s=1$,
and
\begin{equation}\label{eq:strict-marker-derivative}
 -\rho'(s)>\left(1-\frac1{2s}\right)\Delta_s.
\end{equation}
Thus, with $e=\Delta_1=1-\rho_2$, one has $0<e<1$ and
$0<\Delta_i\le i e$ for every integer $i\ge1$.
\end{lemma}
\begin{proof}
Apply the chosen-child inequality~\eqref{eq:chosen-child-comparison} to
increasing indicator functions.  After the same change of variables as
in the proof of Lemma~\ref{lem:canonical-moments}, it says that
\[
 A_m(t)=\E[(1-\lambda U)\mathbf1_{\{\lambda U>t\}}]
       -d\E[(1-U)\mathbf1_{\{U>t\}}]\ge0,
 \qquad 0<t<1.
\]
Layer-cake integration gives
\[
 h\{\rho(s)-\rho(s+1)\}=s\int_0^1t^{s-1}A_m(t)\dd t,
 \qquad s\ge1.
\]
Indeed, the right side equals
$(\lambda^s-d)a_s-(\lambda^{s+1}-d)a_{s+1}$.
The bounded-convergence argument in Lemma~\ref{lem:canonical-moments}
also gives $\rho(s+n)\to0$ for every real $s\ge1$.  Sum the last
display over $s,s+1,\ldots$, using Tonelli and
\[
 \sum_{j\ge0}(s+j)t^{s+j-1}
 =\frac{t^{s-1}\{s-(s-1)t\}}{(1-t)^2}.
\]
At $s=1$ the resulting left side is $\rho(1)=1$. Consequently
\[
 \sigma(\dd t)=\frac{A_m(t)}{h(1-t)^2}\dd t
\]
is a probability measure and gives~\eqref{eq:marker-representation}.
The measure gives full mass to $(0,1)$.

For $s>1$, differentiation under this probability integral is justified
locally by boundedness of $t^{s-1}|\log t|$.  At $s=1$, the derivative
of $a_s/\beta_s$ is finite because $U^s|\log U|$ is bounded when $s$
is bounded away from zero.  Fatou's lemma applied to
\[
 \frac{1-t^{\varepsilon}\{1+\varepsilon(1-t)\}}{\varepsilon}
\]
then shows that $-\log t-(1-t)$ is $\sigma$-integrable.  Hence
$|\log t|$ is integrable and dominated convergence also permits
differentiation from the right at $s=1$. We obtain
\[
 -\rho'(s)=\int_0^1 t^{s-1}
 \bigl[\{1+(s-1)(1-t)\}(-\log t)-(1-t)\bigr]\sigma(\dd t).
\]
Write $z=1-t$.  Since $-\log(1-z)>z+z^2/2$ for $0<z<1$, the
bracket is strictly larger than
$(s-\tfrac12)z^2+(s-1)z^3/2$.
Integration proves~\eqref{eq:strict-marker-derivative}.  Finally,
the formula for $\Delta_i$ and $t^{i-1}\le1$ give $\Delta_i\le i e$.
Strict positivity follows from the support of $\sigma$.
\end{proof}

We now replace the speed $v$ by a real parameter $w$ while keeping the
law of $X$ fixed. Since $wZ-N$ vanishes only at $w=v$, we keep this
term when substituting an upper bound for $v$. In the next two lemmas,
$w$ denotes this parameter. It previously denoted $\E W$
in~\eqref{eq:archive-parameters}.

\begin{lemma}[An affine comparison at fixed escape law]\label{lem:assigned-floor}
Keep the escape law $\pi_\lambda$ and its moments fixed, and put
\[
 A_0=1-(\lambda+1)u+\lambda a_2,\qquad
 a(w)=\frac{u-(w/c)A_0}{a_2},
\]
\[
 P_c(w;y)=(1-y)\{2w-(1+w)y\}.
\]
Define the affine function
\begin{equation}\label{eq:assigned-floor-definition}
 \widetilde F(w)=2a(w)\E[br]
 -\E[b\cJ P_c(w;\cdot)]
 +2\E\!\left[g\cJ\{(1-a(w))y^2+c((1+w)y-w)\}\right].
\end{equation}
If $0<\alpha<1$ in Lemma~\ref{lem:calibration}, then
$\widetilde F(v)=F_*$ and $\widetilde F'(w)<0$ for every real $w$.
In particular, $F_*>\widetilde F(w)$ when $w>v$.
\end{lemma}
\begin{proof}
Set $\Gamma=\E[br]-\E[g\cJ(y^2)]$.  Apply
\eqref{eq:entropy-stationarity} with the strictly positive function $r$,
using $(I-\cL)r=\cJ(y^2)$.  It gives
\[
 \Gamma=\E[\delta_g r(F_\lambda(S))]>0.
\]
All these expectations are finite by~\eqref{eq:entropy-integrability}.
To find the dependence on $w$, symmetrize the bounded
pair kernels.  For $t=x+y-xy$,
\[
 \frac{b(x)P_c(w;y)+b(y)P_c(w;x)}2
 =\frac{(1-t)(wt-xy)}\lambda.
\]
Since $1-t=\lambda-A$, this implies the exact correction
\begin{equation}\label{eq:assigned-C-correction}
 \E\frac{wt-XY}{A^2}-\E[b\cJ P_c(w;\cdot)]
 =\frac{wZ-N}\lambda.
\end{equation}
At $w=v=N/Z$ the correction vanishes, giving
$\widetilde F(v)=F_*$.  Differentiating the affine expression with the
law fixed gives
\[
 \widetilde F'(w)=2c\E[g\cJ(y-1)]
 -\E[b\cJ((1-y)(2-y))]-\frac{2A_0\Gamma}{c a_2}<0.
\]
Here $A_0>0$, $a_2>0$, $\Gamma>0$, $g\ge0$, and
$0<X,Y<1$ almost surely. Only the scalar $w$ varies in this derivative.
\end{proof}

\begin{lemma}[A series lower bound in terms of the moments]\label{lem:mellin-floor}
Fix $w$ satisfying $v\le w<h$ and $0<a(w)<1$.
Set
\begin{align}
 Q_{a(w)}(z)&=-\frac wc+\left(1+(\lambda+1)\frac wc\right)z
                   -\left(a(w)+\lambda\frac wc\right)z^2,
 \nonumber\\
 t_n(w)&=\E[U^nQ_{a(w)}(U)],\label{eq:assigned-t}\\
 J_{s-1}(w)&=-2w a_{s-1}
  +\{c(1+2a(w))+(3\lambda+1)w\}a_s
  \nonumber\\
 &\quad-\{c(\lambda+2a(w))+\lambda(\lambda+1)w\}a_{s+1},
 \label{eq:assigned-J}\\
 \overline\gamma_{s-1}
 &=a_s(r_s-r_{s-1})-\delta_s
   +k\left(1-\frac1{2s}\right)\frac{\delta_s}{r_s}.
 \label{eq:gamma-floor}
\end{align}
Then $t_0(w)=0$, $t_n(w)>0$ for every $n\ge1$, and the series
\begin{equation}\label{eq:complete-mellin-floor}
 F_{\rm floor}(w)=\sum_{s\ge1}s\lambda^{s-1}
 \left\{\frac{(a_s-\lambda a_{s+1})J_{s-1}(w)}{\lambda c}
             +2\overline\gamma_{s-1}t_{s-1}(w)\right\}
\end{equation}
converges absolutely.  Moreover,
\begin{equation}\label{eq:strict-Mellin-gap}
 \widetilde F(w)-F_{\rm floor}(w)
 =2\sum_{s\ge2}s\lambda^{s-1}
 \left\{k\beta_s[-\rho'(s)]
 -k\left(1-\frac1{2s}\right)\frac{\delta_s}{r_s}\right\}
 t_{s-1}(w)>0.
\end{equation}
In particular, $D_\lambda\ge F_*>F_{\rm floor}(v)$.
\end{lemma}
\begin{proof}
The choice of $a(w)$ gives $t_0(w)=0$. The quadratic $Q_{a(w)}$ is negative
at zero, strictly concave, and has value $(\lambda-a(w))/\lambda^2>0$
at $1/\lambda$.  It therefore has exactly one zero $z_0$ in
$(0,1/\lambda)$, with the sign of $z-z_0$ on that interval.  For
$n\ge1$,
\[
 t_n(w)=\E[(U^n-z_0^n)Q_{a(w)}(U)]>0.
\]
The integrand is nonnegative and vanishes only at $U=z_0$.
The tree starts with a chain of single-child vertices of any given
length with positive probability. Along such a chain of length $j$,
the transform $U$ is multiplied by $\lambda^{-j}$. Thus the
essential infimum of $U$ is zero, although $U>0$ almost surely. In particular, its law cannot be concentrated at $z_0>0$.

Write $V=Y/(c+Y)$ for an independent copy of $U$.  Direct substitution
gives, with $a=a(w)$,
\begin{align*}
 \frac{b(X)Y}{A^2}
 &=\frac{UV(1-\lambda U)(1-V)}{\lambda(1-\lambda UV)^2},\\
 \frac{b(X)P_c(w;Y)}{A^2}
 &=\frac{U(1-\lambda U)(1-\lambda V)
            \{2w-(c+(\lambda+1)w)V\}}
          {\lambda c(1-\lambda UV)^2},\\
 \frac{(1-a)Y^2+c((1+w)Y-w)}{A^2}
 &=\frac{(1-U)^2Q_a(V)}{(1-\lambda UV)^2}.
\end{align*}
Since $\lambda UV\le1/\lambda<1$, expand the common denominator in
its uniformly convergent series.  With
\[
 \gamma_{s-1}=\E[g(X)(1-U)^2U^{s-1}]
 =k\E[U^s\log(1/U)]-a_s+a_{s+1},
\]
the first two kernels produce
$(a_s-\lambda a_{s+1})J_{s-1}(w)/(\lambda c)$, and the third produces
$2\gamma_{s-1}t_{s-1}(w)$.  Their sum is exactly
$\widetilde F(w)$.  In particular, the first factor is
$a_s-\lambda a_{s+1}$ for every index $s\ge1$, including the first.

Differentiate $a_s=\beta_s\rho(s)$ with respect to the moment order $s$.
Since
$-k\beta_s'/\beta_s=1-r_{s-1}$, one obtains
\[
 \gamma_{s-1}=a_s(r_s-r_{s-1})-\delta_s
                    +k\beta_s[-\rho'(s)].
\]
Use Lemma~\ref{lem:canonical-marker},
$\beta_s\Delta_s=\delta_s/r_s$, and $t_{s-1}(w)>0$ for $s\ge2$.
The term $s=1$ has $t_0(w)=0$.  This proves the equality and strict
inequality in~\eqref{eq:strict-Mellin-gap}.

For completeness, $a_s\le\lambda^{-s}$ and
$\delta_s\le\lambda^{-s-1}$.  The ratios $r_s$ are bounded away from
zero at fixed $\lambda,h$, and $r_s-r_{s-1}$ is bounded.  The integral
defining $\gamma_{s-1}$ is bounded by
$\|g\|_\infty\lambda^{-s+1}$; its identity in the preceding display
therefore bounds $k\beta_s[-\rho'(s)]$ by a constant times
$\lambda^{-s}$.  Also $t_{s-1}(w)$ and $J_{s-1}(w)$ have bounds of
the same geometric order.  Every expansion and regrouping above is
dominated by a constant times $\sum_{s\ge1}s\lambda^{-s}$.
No unweighted logarithmic moment or inverse escape moment is used.
\end{proof}

\subsection{Series estimates for \texorpdfstring{$1<\lambda<2/\sqrt3$}{1 < lambda < 2/sqrt(3)}}
\label{subsec:low-computer-region}
Fix $1<\lambda<\min(m,2/\sqrt3)$. Lemma~\ref{lem:calibration}
gives $D_\lambda>0$ when $\alpha=1$. For $0<\alpha<1$, put
\[\begin{gathered}
 q=\frac{u}{1-u}=\frac hc,\qquad R_* =\frac vc,\qquad
 w_i=\lambda^{i-1}\beta_i^2,\\
 T_p=\sum_{i\ge1}w_i,\qquad Z_p=2q-(\lambda+1)T_p,\qquad
 R=\frac{T_p}{Z_p},\quad \theta=\frac Rq.
\end{gathered}\]
The denominator $Z_p$ is positive by
Lemma~\ref{lem:harmonic-speed}, which also gives $R_*<R<q$.
Write
\begin{equation}\label{eq:low-reference-calibration}
 \alpha_p=\frac{u-R\{1-(\lambda+1)u+\lambda\beta_2\}}{\beta_2}
 =\frac{\lambda^2q+(1-\theta)d(\lambda+1)}{q+1},
 \qquad C_\alpha=\alpha_p+\lambda R.
\end{equation}
Thus $\lambda^2u<\alpha_p<a(cR)<\alpha<1$.  In particular,
$0<h<1/(\lambda+1)<1/2$. Throughout the comparison, the offspring
law is fixed.

\subsubsection{A lower bound for \texorpdfstring{$F_{\rm floor}(v)$}{the moment series at v}}
We bound the series for $F_{\rm floor}(v)$ in terms of
$e=1-\rho_2$. We keep the term at $i=1$, the boundary terms from
summation by parts, and the contribution of $R-R_*$.
For $i\ge1$ define
\begin{align}
 L_i&=\frac1\lambda-r_{i-1}
       =\frac{cd}{\lambda(\lambda^i-d)},\nonumber\\
 W_i^0&=\frac{ci}{\lambda}-1+2iL_i,\nonumber\\
 W_i^p&=cW_i^0+4iL_i
       -2r_i\{iL_i+(i+1)\lambda L_{i+1}\},
 \nonumber\\
 W_i(R)&=W_i^0+RW_i^p,\label{eq:low-W}\\
 V_i&=\sum_{s>i}w_sW_s(R),
 \nonumber\\
 F_{\rm pipe}&=-2RuL_1+u^2W_1(R)+V_1,\nonumber\\
 \tau_i&=1-C_\alpha r_i+(\lambda+1)R-R/r_{i-1},\nonumber\\
 B_i&=w_i\left\{\frac{2i(\alpha_p+cR)}\lambda
       +2iRr_iL_i+2(i+1)R\lambda r_i^2L_i
       +k(2i-1)\tau_i\right\},
 \nonumber\\
 S_i&=B_i-2V_i.\label{eq:low-B-S}
\end{align}
Set $W_1^Y=0$, $W_i^Y=k w_i(2i-1)R/r_{i-1}$ for $i\ge2$, and
\begin{equation}\label{eq:low-J}
 J_i=V_i+k w_i(2i-1)C_\alpha r_i
                         -\frac{W_i^Y+W_{i+1}^Y}{2},\qquad i\ge2.
\end{equation}
Here $J_i$ is the coefficient of $\Delta_i^2$ in the bound below.
The polynomial $J_{s-1}(w)$ in~\eqref{eq:assigned-J} carries the
argument $w$.

\begin{lemma}[A bound in terms of the differences $\Delta_i$]\label{lem:calibrated-head}
Suppose $0<\alpha<1$. Then $V_i>0$ for every
$i\ge1$, $\tau_1=0$, and $0<\tau_i<1$ for $i\ge2$.  Put
\[
 \kappa=\frac{(1-u)T_p}{u^2},
\]
\[\begin{split}
 C_{\rm head}={}&2A_0u\{-k\log(a_2/u)-1+a_2/u\}\\
 &+\frac{(u-\lambda a_2)
       \{2-(3\lambda+1)u+\lambda(\lambda+1)a_2\}}\lambda.
\end{split}\]
Then $C_{\rm head}>0$, $\kappa>1$, and
\begin{align}
 F_{\rm floor}(v)>{}&F_{\rm pipe}
 +\sum_{i\ge1}\Delta_i\rho_i S_i
 +e^2\left\{V_1+\frac{2u^2 C_\alpha}{\lambda(1-e)}
                         -\frac{W_2^Y}{2}\right\}
 +\sum_{i\ge2}\Delta_i^2J_i+C_{\rm cal}e,
 \label{eq:low-master}\\
 C_{\rm cal}&=\frac{C_{\rm head}u^2(\kappa-1)(2-e)}{Z_p}>0.
 \nonumber
\end{align}
The coefficient in braces in~\eqref{eq:low-master} is positive.  Also
\begin{equation}\label{eq:low-e-caps}
 \begin{gathered}
 0<e< E_{\rm hard}=\frac{1-\alpha_p}{1+\lambda R},\qquad
 e\le E_{\rm var}=\frac{d}{(q+1)^2},\\
 e q^2<\frac{\{1-(\lambda+1)h\}u^2}{1+cu},\qquad
 \sqrt e<\frac{d}{q+1}.
 \end{gathered}
\end{equation}
Every series in~\eqref{eq:low-master} converges absolutely.
\end{lemma}
\begin{proof}
First compare $F_{\rm floor}(v)$ with $F_{\rm floor}(cR)$.
With the moments fixed, differentiate $t_n(cR)$ in~\eqref{eq:assigned-t}
with respect to $R$. If the law of $U$ is weighted by $U^2/a_2$, then
\[\begin{aligned}
 \frac{\dd}{\dd R}t_n(cR)
 &=\frac{A_0}{a_2}a_{n+2}
       -\E[U^n(1-U)(1-\lambda U)]\\
 &=-a_2\operatorname{Cov}_{U^2/a_2}
       \left(U^n,\frac{(1-U)(1-\lambda U)}{U^2}\right)>0
\end{aligned}\]
for $n\ge1$.  The second function is strictly decreasing on
$(0,1/\lambda)$; all its weighted expectations are finite.  The
derivative is also bounded above by $(A_0/a_2)a_{n+2}$.
Thus the difference $\widetilde F(cR)-F_{\rm floor}(cR)$ in~\eqref{eq:strict-Mellin-gap} increases with
$R$, and the geometric bounds in Lemma~\ref{lem:mellin-floor}
justify termwise differentiation.

The first two terms in $-c\widetilde F'(w)$ have the lower bounds
\begin{align*}
 2c^2\E[g\cJ(1-y)]&\ge2A_0\gamma_0,\\
 c\E[b\cJ((1-y)(2-y))]&\ge
 \frac{(u-\lambda a_2)
       \{2-(3\lambda+1)u+\lambda(\lambda+1)a_2\}}\lambda.
\end{align*}
They follow by the $U,V$ substitution and
$(1-\lambda UV)^{-2}\ge1$; the remaining factors are positive.
Jensen's inequality under the probability weighted by $U/u$ gives
$\gamma_0\ge u\{-k\log(a_2/u)-1+a_2/u\}$.
The expression in braces is positive: the function
$-k\log z-1+z$ vanishes at $z=1/\lambda$ and has negative derivative
for $0<z<1/\lambda<k$. The term involving $\Gamma$ in
$-c\widetilde F'(w)$ makes the bound strict. Consequently
\begin{equation}\label{eq:low-slope-credit}
 F_{\rm floor}(v)>F_{\rm floor}(cR)+C_{\rm head}(R-R_*).
\end{equation}

We next bound $R-R_*$ for the same offspring law. The identity
$\lambda^{s-1}\beta_s\beta_{s-1}=q(\beta_{s-1}-\beta_s)$,
finite summation by parts, and the vanishing last endpoint give
\begin{align*}
 q-Q&=q\sum_{i\ge1}\Delta_i(a_i+a_{i+1}),\\
 T_p-T&=\sum_{i\ge1}\Delta_i(\rho_i+\rho_{i+1})
                              \sum_{s>i}w_s,\\
 R-R_*&=\frac{2q}{Z_p Z}\sum_{i\ge1}\Delta_i
 \left\{(\rho_i+\rho_{i+1})\sum_{s>i}w_s-T_p(a_i+a_{i+1})\right\}.
\end{align*}
All tail terms in the first two formulas are nonnegative.  The ratios
$r_i$ increase strictly.  With $f_i=(\sum_{s>i}w_s)/\beta_i$, the same
telescoping sum gives
\[
 f_i-f_{i+1}=\lambda^i
 \left\{\frac{\beta_{i+1}^2}{\beta_i}
               -\frac1q\sum_{s>i+1}w_s\right\}<0.
\]
Indeed, each term of the last tail is strictly larger than $r_i$
times the corresponding term $\lambda^{s-1}\beta_s\beta_{s-1}$.
Summing the latter over $s>i+1$ gives $q\beta_{i+1}$. Therefore
$f_i\ge f_1=T_p/u-u=T_p+u(\kappa-1)$.
Solving $\alpha_p<1$ in terms of $\kappa$ gives
$\kappa>2\lambda/(\lambda+1)>1$.  It follows that every expression
in braces in the speed-difference formula is strictly larger than
$u(\kappa-1)\beta_i(\rho_i+\rho_{i+1})$.
Since $0<Z<2q$ and
\[
 \sum_{i\ge1}\beta_i(\rho_i^2-\rho_{i+1}^2)
 \ge u(1-\rho_2^2)=ue(2-e),
\]
we conclude that
\begin{equation}\label{eq:low-speed-deficit}
 R-R_*>\frac{u^2(\kappa-1)e(2-e)}{Z_p}.
\end{equation}

To identify the coefficients of $\Delta_i^2$, put
$\bar\alpha=a(cR)$. The definitions of $a(cR)$ and $\alpha_p$ give
\[
 \bar\alpha-\alpha_p=\frac{C_\alpha e}{1-e}.
\]
For $n_i=a_{i+1}/a_i$ define
$c_i=2i/\lambda-k(2i-1)n_i>1/\lambda$.
Substitute $a_i=\beta_i\rho_i$ and
$\delta_i=\beta_{i+1}\Delta_i$ in the series
\eqref{eq:complete-mellin-floor}.  Use
$(1+L_i)r_i=1/\lambda$ and sum the diagonal terms by parts.  The
result, with every end term displayed, is
\begin{align}
 F_{\rm floor}(cR)={}&F_{\rm pipe}
 +\sum_{i\ge1}\Delta_i\rho_i S_i+\sum_{i\ge1}\Delta_i^2V_i
 +(\bar\alpha-\alpha_p)\sum_{i\ge1}w_i\rho_i\Delta_i c_i
 \nonumber\\[-2pt]
 &+k\sum_{i\ge1}w_i(2i-1)
 \left\{C_\alpha r_i\Delta_i^2
                  -\frac{R\Delta_i\Delta_{i-1}}{r_{i-1}}\right\},
 \qquad\Delta_0=0.\label{eq:low-square-identity}
\end{align}
One may check the local expansion before summation by noting that its
coefficient of $\Delta_i$ differs from $\rho_i B_i$ by
\[
 w_i\rho_i(\bar\alpha-\alpha_p)c_i
 +k w_i(2i-1)\left\{C_\alpha r_i\Delta_i
                         -R\Delta_{i-1}/r_{i-1}\right\}.
\]
More explicitly, before replacing the diagonal terms by the $V_i$
tails, the difference between the $i$th summand of $F_{\rm floor}(cR)$ and these local
terms is
\begin{align*}
 &-(i-1)(1+cR)w_i\rho_i^2
       +i(1+cR)w_{i+1}\rho_{i+1}^2\\
 &\hspace{5mm}
 -\frac{2iRw_iL_i}{r_{i-1}}\rho_{i-1}\rho_i
       +2(i+1)Rw_{i+1}L_i\rho_i\rho_{i+1}.
\end{align*}
The first pair telescopes to zero.  The second pair telescopes because
$L_{i+1}/r_i=L_i$, leaving exactly
$-2Rw_1L_1/r_0=-2RuL_1$.  Its upper endpoint tends to zero
geometrically.  Finally the diagonal telescope uses
$1-\rho_s^2=\sum_{i<s}\Delta_i(\rho_i+\rho_{i+1})$.
This verifies~\eqref{eq:low-square-identity}, including the boundary
term at $i=1$ and the limit at infinity.

At $i=1$ the two $k$ contributions cancel exactly, leaving
$e^2\{V_1+2u^2C_\alpha/[\lambda(1-e)]\}$.
For each $i\ge2$ apply
$2\Delta_i\Delta_{i-1}\le\Delta_i^2+\Delta_{i-1}^2$ to the
remaining adjacent product. The term at $i=2$ contributes
$-W_2^Y e^2/2$, and the other squares have
the coefficients~\eqref{eq:low-J}.  The other terms involving
$\bar\alpha-\alpha_p$ are positive and may be dropped.
Equations~\eqref{eq:low-slope-credit} and
\eqref{eq:low-speed-deficit} now prove~\eqref{eq:low-master}.
The tail positivity $V_i>0$ is proved in the next lemma.  Also
$C_\alpha/(\lambda R)>1$ and $k\lambda r_1<1$, so
\[
 \frac{2u^2C_\alpha}{\lambda(1-e)}>2u^2R,
 \qquad \frac{W_2^Y}{2}<\frac32u^2R.
\]
This proves positivity of the coefficient in braces in~\eqref{eq:low-master}.
Under the distribution $\PP(U_p=\lambda^{-j})=hd^{j-1}$, the quadratic
$Q_{\alpha_p}$ has zero mean and changes sign once from negative
to positive. The proof of Lemma~\ref{lem:mellin-floor} applied to this
scalar distribution gives $\tau_1=0$ and $\tau_i>0$ for $i\ge2$.
Furthermore
\[
 (\lambda+1)-\lambda r_i-1/r_{i-1}
 =-\frac{(\lambda+1)L_i(c+\lambda L_i)}
          {(1+L_i)(1-\lambda L_i)}<0,
\]
so $\tau_i<1$.

Substituting $\bar\alpha-\alpha_p=C_\alpha e/(1-e)$ into $\bar\alpha<1$
gives $e<E_{\rm hard}$. Jensen's inequality $a_2\ge u^2$ gives
$e\le1-u^2/\beta_2=(1-u)(1-\lambda u)=d/(q+1)^2$.
Writing
\[
 E_{\rm hard}=\frac{\theta d(\lambda+1)-\lambda}
                   {(q+1)(1+\lambda\theta q)}
\]
shows that it increases with $\theta$.  Use $\theta<1$ to obtain
the third bound in~\eqref{eq:low-e-caps}.
Finally, $d^2-\{1-(\lambda+1)h\}=h^2+ch>0$ gives the stated
bound on $\sqrt e$.  Absolute convergence follows from
$\beta_i,a_i\le\lambda^{-i}$ and the polynomial factors in $i$;
finite summation followed by geometric domination justifies all
signed rearrangements in this proof.
\end{proof}

\subsubsection{\texorpdfstring{Bounds for the tails $V_i$ and $\sum_{s>i}w_s$}{Bounds for V(i) and the tail sum of w(s)}}
We record the estimates that turn the infinite sums into finite rational
expressions.  Let $\delta=\log\lambda$ and $a=c/\delta$.  Taylor's
formula with integral remainder applied to
$a=\int_0^1(1+c)^s\dd s$ gives
\begin{equation}\label{eq:low-log-enclosure}
 a_L=1+c/2-c^2/12\le a\le
 a_U=1+c/2-c^2/12+c^3/24\le1+c/2.
\end{equation}
For every $z>0$,
\begin{align}
 1+z/2+z^2/12-z^4/720
 &\le\frac{z}{1-e^{-z}}\le1+z/2+z^2/12,\label{eq:low-global-F}\\
 1-z^2/12&\le\frac{z^2}{2(\cosh z-1)}
                       \le1-z^2/12+z^4/240.\label{eq:low-global-C}
\end{align}
For example, multiply the two inequalities in~\eqref{eq:low-global-F}
by the positive exponential denominator.  Beyond the vanishing low
powers, their coefficients, with the indicated sign, are respectively
$(n-3)(n-5)(n-6)(n+8)/(720n!)$ for $n\ge7$ and
$(n-3)(n-4)/(12n!)$ for $n\ge5$.  For~\eqref{eq:low-global-C},
the corresponding even coefficients are
$\{n(n-1)-12\}/(6n!)$ for $n\ge6$ and
$(n-6)(n^3-9n-40)/(120n!)$ for $n\ge8$.
All omitted lower coefficients vanish.  Absolute convergence of the
exponential series proves the inequalities on the entire positive axis.

\begin{lemma}[Bounds for the tails of the Lambert series]\label{lem:low-Lambert-tails}
For every $i\ge1$, $V_i>0$.  If $b=\lambda L_i$, put
\[
 r=\frac1{\lambda+b}=r_i,\qquad
 A=c+br,\quad B=br,\quad t=\frac r{b+c},
\]
and define the five polynomials
\begin{align*}
 P_1&=A,\qquad P_2=A(A+B),\qquad P_3=A(A^2+4AB+B^2),\\
 P_4&=A(A^3+11A^2B+11AB^2+B^3),\\
 P_5&=A(A^4+26A^3B+66A^2B^2+26AB^3+B^4).
\end{align*}
Define the following explicit rational functions of an argument $a>0$:
\begin{align}
 T_L(a)&=t\{a+A/2+P_2/(12a)-P_4/(720a^3)\},\nonumber\\
 T_U(a)&=t\{a+A/2+P_2/(12a)\},\nonumber\\
 P_U(a)&=\frac t\lambda
 \{(i+1)(aA+P_2/2+P_3/(12a))+a^2-P_2/12+P_4/(240a^2)\},\nonumber\\
 P_L(a)&=\frac t\lambda
 \{(i+1)(aA+P_2/2+P_3/(12a)-P_5/(720a^3))
                          +a^2-P_2/12\}.\label{eq:low-tail-rationals}
\end{align}
At $a=c/\log\lambda$,
\begin{align}
 V_i/w_i&\le V_U(a)
 :=-(1+cR)T_L(a)+\{1+(\lambda+1)R\}P_U(a)
                   -2R(i+1)\lambda r^2-2Rt,\nonumber\\
 V_i/w_i&\ge-(1+cR)T_U(a)+\{1+(\lambda+1)R\}P_L(a)
                   -2R(i+1)\lambda r^2-2Rt.\label{eq:low-V-bound}
\end{align}
For $q\le10$, $c\le1/6$, and every integer $i\ge1$, these imply the rational
bounds
\begin{equation}\label{eq:low-V-rational-bound}
 V_L\le V_i/w_i\le V_U,
\end{equation}
where
\[\begin{aligned}
 V_U&=V_U(a_U),\\
 V_L&=-(1+cR)T_U(a_U)+\{1+(\lambda+1)R\}P_L(a_L)\\
 &\quad-2R(i+1)\lambda r^2-2Rt.
\end{aligned}\]
The tail $\sum_{s>i}w_s$ lies between $w_iT_L(a_L)$ and
$w_iT_U(a_U)$ on this domain.
\end{lemma}
\begin{proof}
Geometric expansion of $\beta_s^2$ and summation over every $s>i$ give
\[
 \sum_{s>i}w_s=\frac{h^2}{\lambda\delta}
 \sum_{N\ge1}d^{N-1}\lambda^{-N(i+1)}
                         \frac{N\delta}{1-e^{-N\delta}}.
\]
Hold $h,d$ fixed while differentiating this deterministic series with
respect to $\lambda$.  A direct derivative of its $N$th summand gives
the identity
\begin{align}
 V_i={}&-(1+cR)\sum_{s>i}w_s-2R(i+1)w_{i+1}-\frac{2Rq\beta_{i+1}}\lambda
 \nonumber\\
 &+\{1+(\lambda+1)R\}\frac{ch^2}{\lambda^2}
 \Biggl\{\frac{i+1}{\delta}\sum_{N\ge1}N d^{N-1}\lambda^{-N(i+1)}
                                      \frac{N\delta}{1-e^{-N\delta}}
 \nonumber\\[-2pt]
 &\hspace{39mm}+\frac1{\delta^2}
 \sum_{N\ge1}d^{N-1}\lambda^{-N(i+1)}
                          \frac{(N\delta)^2}{2(\cosh(N\delta)-1)}\Biggr\}.
 \label{eq:low-full-tail-identity}
\end{align}
The terms $-2R(i+1)w_{i+1}$ and $-2Rq\beta_{i+1}/\lambda$ come from
shifting the summation indices in the terms involving $W_i^p$.
For a direct check of the differentiation, write
$y=\lambda^{-N}$ and
$T_N=h^2N d^{N-1}y^{i+1}/[\lambda(1-y)]$.  Then
\[
 T_N'=-T_N/\lambda-
 \frac{h^2N^2d^{N-1}y^{i+1}}{\lambda^2}
 \left\{\frac{i+1}{1-y}+\frac{y}{(1-y)^2}\right\}.
\]
Substitution proves~\eqref{eq:low-full-tail-identity} term by term.
On a compact neighborhood of a fixed $\lambda>1$, the summands and
their derivatives are bounded by a geometric sequence times a fixed
polynomial in $N$.  This justifies the derivative and both summations.

For $p=\lambda^{i+1}$, the geometric moments are
$M_j=\sum_{N\ge1}N^j d^{N-1}p^{-N}$, where
$M_0=(p-d)^{-1}$ and $M_{j+1}=-p\partial_p M_j$.
Applying this recurrence up to $j=5$ gives
\[
 \frac{hM_0}{w_i}=\frac{\lambda r}{q(b+c)},\qquad
 \frac{h^{j+1}M_j}{w_i}
    =\frac{q^{j-1}\lambda rP_j}{b+c}\quad(1\le j\le5).
\]
Use~\eqref{eq:low-global-F}--\eqref{eq:low-global-C} in every
$N$ summand of~\eqref{eq:low-full-tail-identity}.  The sign of the
tail coefficient is negative and those of the two derivative kernels
are positive.  The resulting inequalities are precisely
\eqref{eq:low-tail-rationals}--\eqref{eq:low-V-bound}.

For $q\le10$ we have $c\le1/6$ and
$b\le d/(q+i)$.  Hence $A\le2/3$, $0\le B\le A$,
$P_2\le8/9$, $P_3/A\le8/3$, and $P_4\le128/27$.
Differentiation of the explicit rational functions gives
\[
 0<T_L'(a)\le\frac{11t}{10},\qquad
 P_U'(a)\ge\frac{9t}{5\lambda},\qquad
 T_U'(a)>0,\quad P_L'(a)>0\qquad(a\ge1).
\]
For instance, $T_L'/t=1-P_2/(12a^2)+P_4/(240a^4)$,
and $\lambda P_U'/t=(i+1)(A-P_3/(12a^2))+2a-P_4/(120a^3)$;
the displayed polynomial bounds prove these estimates.
The derivative of $V_U(a)$ is positive since
\[
 \frac95\{1+(\lambda+1)R\}-\frac{11}{10}\lambda(1+cR)>0
 \qquad(1<\lambda\le7/6, R\ge0).
\]
Now~\eqref{eq:low-log-enclosure} proves
\eqref{eq:low-V-rational-bound} and the tail enclosure.

To prove positivity, put $z=d\lambda^{-i}$ and
$\Pi(N)=N(1-z)^2z^{N-1}$ for $N\ge1$.  This is a probability
distribution.  The same geometric summation, including the two
endpoints, gives
\begin{equation}\label{eq:low-positive-mixture}
 \frac{V_i}{w_i}=\frac1\lambda\E_\Pi
       [i(A_N+RP_N)+B_N+RC_N],
\end{equation}
where, writing $Q_N=1+\lambda+\cdots+\lambda^{N-1}$,
\[\begin{gathered}
 A_N=\frac{N}{Q_N},\quad B_N=\frac{\lambda Q_N'}{Q_N^2},\quad
 P_N=(\lambda+1)A_N-2\lambda^{1-N},\\
 C_N=(\lambda+1)B_N-
          \frac{2\lambda(1-A_N)}{cN\lambda^N}.
\end{gathered}\]
Here $A_N>0$, $B_N\ge0$, and
$P_N\ge cN\lambda^{-N}>0$: the latter is the convex-chord
inequality $N(\lambda^{N+1}+1)\ge2\sum_{j=1}^N\lambda^j$.
Also $Q_N-N\le cQ_N'$ term by term, and hence
\[
 C_N\ge B_N\left\{(\lambda+1)-\frac{2Q_N}{N\lambda^N}\right\}\ge0.
\]
This proves $V_i>0$ and completes the proof.
\end{proof}

The tail estimates also provide power-sum enclosures used below.  For
$A=q+2$ and every integer $j\ge2$, put
\begin{equation}\label{eq:low-power-enclosures}
 U_j=\frac{A^{1-j}}{j-1}+\frac1{2A^j}+\frac{j}{12A^{j+1}},\qquad
 L_j=U_j-\frac{j(j+1)(j+2)}{720A^{j+3}}.
\end{equation}
Then $L_j\le\sum_{i\ge2}(q+i)^{-j}\le U_j$.
Indeed, write the sum as
\[
 \frac1{(j-1)!}\int_0^\infty
       \frac{z^{j-1}e^{-Az}}{1-e^{-z}}\dd z
\]
and insert the global bounds~\eqref{eq:low-global-F}; Tonelli and
the gamma integral evaluate every term.  Thus~\eqref{eq:low-power-enclosures}
includes the entire infinite tail.

\subsubsection{The range \texorpdfstring{$0<q\le1/15$}{0 < q <= 1/15}}
Here $u\le1/16$. For the logarithmic term we retain the first two
terms in the definition of $\overline\gamma_{s-1}$. Define
$F_{\rm cap}$ by the right side of~\eqref{eq:complete-mellin-floor},
with $w=v$ and with $\overline\gamma_{s-1}$ replaced by
$a_s(r_s-r_{s-1})-\delta_s$.
The exact logarithmic identity in Lemma~\ref{lem:mellin-floor} gives
$D_\lambda\ge F_*\ge F_{\rm cap}$.
For this paragraph only put
\begin{align*}
 B_s&=\frac{c}{\lambda^s-1},& f_s&=\lambda^{s-1}B_s^2,&
 R_s&=1-u+uB_s,& \frac{\beta_s}{u}&=\frac{B_s}{R_s},\\
 t&=T/u^2,&q_a&=Q/u,&t_\pi&=\sum_{s\ge1}f_s/R_s^2,\\
 H&=\sum_{s\ge1}s f_s^2\rho_s^2/R_s^3,&
 H_\pi&=\sum_{s\ge1}s f_s^2/R_s^3,\\
 p_s&=\frac{cs}{\lambda}\frac{\lambda^s+1}{\lambda^s-1}
                 -\frac{\lambda+1}\lambda,&
 P&=\sum_{s\ge1}f_sp_s\rho_s^2/R_s^2,&
 M&=\frac{2q_a}{2q_a-(\lambda+1)ut}.
\end{align*}
The function $x\coth x$ increases on $(0,\infty)$, because its
derivative has numerator $\sinh(2x)-2x>0$.  Hence $p_s\ge0$.
Finally set
\[
 G_s=\frac{s\lambda^{s-1}(1+L_s)a_s\delta_s}{u^2},\quad
 G=\sum_{s\ge1}G_s,\quad
 G_+=\sum_{s\ge1}G_s\left\{\lambda+
 \frac{(s+1)\lambda a_{s+1}}{s a_s(1+L_s)}\right\}\ge\lambda G.
\]

Separating the term at $s=1$ gives
\begin{equation}\label{eq:low-small-raw}
 \frac{F_{\rm cap}}{u^2}
 =M\left\{P+u\left(\frac{t^2}{q_a}-2H\right)
                 +\frac{ut}{q_a}G_+\right\}+2\alpha G.
\end{equation}
Here is a verification that also specifies the endpoint in this
identity.  Put $H_s^{\rm loc}=s-(s+1)\lambda L_{s+1}$ and
$D_s^{\rm loc}=2sL_s-s/\lambda+H_s^{\rm loc}r_s$. Expanding
$F_{\rm cap}$ gives
\begin{align*}
 F_{\rm cap}={}&-2R_*uL_1
 +\sum_{s\ge1}\lambda^{s-1}a_s^2
             \{(1+v)W_s^0+2R_*D_s^{\rm loc}\}\\
 &+2\sum_{s\ge1}s\lambda^{s-1}a_s\delta_s
 \left\{\alpha(1+L_s)+R_*\left(c+\lambda L_s
                         +\frac{(s+1)\lambda L_{s+1}}s\right)\right\}.
\end{align*}
The elementary identities
\[
 W_s^0-1/\lambda=p_s-2us f_s/R_s
\]
and, writing $t_N=u^{-2}\sum_{s=1}^N\lambda^{s-1}a_s^2$ and
$Q_N=\sum_{s=1}^N\lambda^{s-1}a_sa_{s-1}$,
\begin{align*}
 \sum_{s=1}^N\frac{\lambda^{s-1}a_s^2}{u^2}
                  (D_s^{\rm loc}-W_s^0)
 ={}&t_N-1-\frac{Q_{N+1}-u}{\lambda u^2}\\
 &+\frac1{u^2}\sum_{s=1}^N\lambda^{s-1}a_s\delta_s
           \left\{H_s^{\rm loc}+(s+1)\lambda\frac{a_{s+1}}{a_s}\right\}\\
 &+\frac{(N+1)\lambda^Na_{N+1}^2}{u^2}
\end{align*}
give~\eqref{eq:low-small-raw} after substitution of
$R_*=ut/[2q_a-(\lambda+1)ut]$.  The last endpoint tends to zero
geometrically. The middle sum involving $\delta_s$ combines with
the other terms proportional to $R_*$ to give the term containing $G_+$.

We next bound the loss $t_\pi-t$ by $G$.  With
$g_i=i\lambda^{i-2}\beta_i^2$,
$u^2G=\sum_i g_i\rho_i\Delta_i$. For the reference moments $\beta_i$,
\[
 \sum_{n>i}\lambda^{n-1}\beta_n\beta_{n-1}=q\beta_i,
 \qquad \sum_{n>i}w_n<q\beta_i/\lambda,
 \qquad \frac{q\beta_i}{g_i}
 =\frac{\lambda^2(1-d\lambda^{-i})}{ci}
 \le\frac{\lambda^2}{1-u}.
\]
The last bound follows by Bernoulli's inequality and
$h=cu/(1-u)$. Summing $1-\rho_s^2=\sum_{i<s}\Delta_i(\rho_i+\rho_{i+1})$
against $w_s/u^2$ and using the tail bounds gives
\begin{equation}\label{eq:low-small-relative-loss}
 t_\pi-t\le\frac{2\lambda G}{1-u}.
\end{equation}
Also $1\le q_a\le(1-u)^{-1}$ and $t\ge1$.
Writing $f_s=[\sinh(\log\lambda/2)/\sinh(s\log\lambda/2)]^2$
shows that $f_s\le s^{-2}$.  Five terms and an integral give
$\sum s^{-2}<5989/3600<5/3$.  Since $f_1=R_1=1$,
\[
 t_\pi\le1+\frac{2}{3(1-u)^2}<\frac{16}9,\qquad
 \kappa=(1-u)t_\pi\le\frac53,\qquad
 M\le\frac{90}{79},\qquad t/q_a\ge\frac{15}{16}.
\]
The comparison $R_*<R$ and $a_2<\beta_2$ gives
\[
 \alpha>\alpha_p
 =\frac{(\lambda+1)(2-\kappa)-2\lambda u}
              {2-(\lambda+1)u\kappa}
 \ge\frac{2/3-(12/5)u}{2-2u(1-u)}\ge\frac{992}{3615}
 \quad(\lambda\le6/5, u\le1/16).
\]
For the middle inequality the numerator is positive, $\kappa\le5/3$,
and $\kappa\ge1-u$; the last rational function decreases on the
stated interval, since its derivative numerator is a positive constant
times $18u^2-10u-13<0$.

Put $B=(1-u)t_\pi^2-2H_\pi$.
Equation~\eqref{eq:low-small-relative-loss} gives
$t^2/q_a-2H\ge B-4\lambda t_\pi G$.
Retaining $G_+$ in~\eqref{eq:low-small-raw} now yields
\begin{equation}\label{eq:low-small-budget}
 \frac{D_\lambda}{u^2}\ge M(P+uB)+C_*G,\qquad
 C_* =\frac{1984}{3615}
 -\frac{90}{79}\frac65\frac1{16}
       \left(\frac{64}{9}-\frac{15}{16}\right)
 =\frac{389899}{18277440}>0.
\end{equation}
The $s=2$ term in $P$ and the $s=1$ term in $G$ give
\[
 P\ge\frac{c^2\rho_2^2}{(\lambda+1)^3},\qquad
 G\ge\frac{1-\rho_2}\lambda.
\]
Thus~\eqref{eq:low-small-budget} also holds in the form
\begin{equation}\label{eq:low-small-root-budget}
 \frac{D_\lambda}{u^2}\ge
 M\{P+uB+K_0(1-\rho_2)\},\qquad
 K_0=\frac{389899}{24986880}.
\end{equation}

For clarity, the finite rational input to the remaining scalar estimate
is stated separately.  For $\ell>1$ let $c_\ell=\ell-1$,
$f_s(\ell)=\ell^{s-1}c_\ell^2/(\ell^s-1)^2$,
$B_s(\ell)=c_\ell/(\ell^s-1)$, and put
\begin{align*}
 \tau(\ell)&=\frac{c_\ell^2}
 {\ell(c_\ell-c_\ell^2/2+c_\ell^3/3)(\ell^9-1)}
                         +\frac{f_9(\ell)}2,\\
 E(\ell)&=\frac{c_\ell^4\ell^{-16}}{6\ell^2}
 \left\{\frac{8}{2c_\ell/(\ell+1)}
                     \frac{3-\ell^{-8}}{(1-\ell^{-8})^3}
 +\frac1{(2c_\ell/(\ell+1))^2}
       \left(\frac{2-\ell^{-8}}{(1-\ell^{-8})^2}-\frac12\right)\right\}.
\end{align*}
\begin{lemma}[Six rational bounds used for $0<u\le1/16$]\label{lem:low-Q0-arithmetic}
The following rational inequalities hold:
\begin{align*}
 \sum_{s=1}^8 f_s(6/5)+\tau(6/5)&>1559/1000,\\
 \sum_{s=1}^8 f_s(23/20)+\tau(23/20)&>63/40,\\
 \sum_{s=1}^8 f_s(11/10)+\tau(11/10)&>19/12,\\
 \sum_{s=1}^8 f_s(6/5)(1-1/s)+\tfrac89\tau(6/5)&>9/25,\\
 \sum_{s=1}^8s^{-3}\{1-B_s(6/5)\}+1/128&<2/15,\\
 \sum_{s=1}^8s f_s(23/20)^2+E(23/20)&<239/200.
\end{align*}
All denominators in these six finite expressions are positive.
\end{lemma}

We now estimate the tails starting at $s=9$. Both $f_s$ and
$B_s$ decrease with $\lambda$; the assertion for $f_s$ follows from
the increasing function $x\coth x$, and that for $B_s$ from its finite
geometric denominator.  As a function of real $x>0$, $f(x)$ is
decreasing and convex, with
\[\begin{aligned}
 f'&=-\log\lambda\, f\coth(x\log\lambda/2),\\
 f''&=(\log\lambda)^2f
       \{\coth^2(x\log\lambda/2)+\tfrac12\operatorname{csch}^2(x\log\lambda/2)\}.
\end{aligned}\]
Convex trapezoid comparison and
$\log\lambda\le c-c^2/2+c^3/3$ imply
$\sum_{s\ge9}f_s\ge\tau(\lambda)$.
The function $xf(x)^2$ decreases; its logarithmic derivative is less
than $-3/x$.  Expand its integral tail using the positive series for
$(1-e^{-x\log\lambda})^{-4}$.  The identities
\[
 \frac{\binom{j+3}{3}}{j+2}=\frac{(j+1)(j+3)}6,\qquad
 \frac{\binom{j+3}{3}}{(j+2)^2}
       =\frac{j+2-1/(j+2)}6
\]
sum the two geometric parts.  Keeping the $j=0$ negative subtraction
and using $\log\lambda\ge2c/(\lambda+1)$ gives
$\sum_{s\ge9}s f_s^2\le E(\lambda)$.
Every remaining coefficient is positive, so this substitution has the
stated direction.  Consequently, with
\[
 T_0=\sum f_s,\quad H_0=\sum s f_s^2,\quad
 A=\sum f_s(1-B_s),\quad C_{\rm tail}=\sum s f_s^2(1-B_s),
\]
Lemma~\ref{lem:low-Q0-arithmetic}, $B_s\le1/s$, and
$\sum_{s\ge9}s^{-3}\le1/128$ prove
\[
 A>9/25,\quad C_{\rm tail}<2/15,\quad H_0<241/200,
 \quad H_0<239/200\ \hbox{if }\lambda\ge23/20,
\]
together with the three lower bounds for $T_0$ on the respective
intervals of $\lambda$ specified below. The general bound for $H_0$ is
$\sum_{s=1}^6s^{-3}+\int_6^\infty x^{-3}\dd x=86701/72000<241/200$.

The positive geometric expansion and the convex chord for
$(1-uz)^{-3}$ on $0\le z\le1$ give
\[
 B>(1-u)(L+18u/25)^2
       -2\{\overline H+(2/15)((1-u)^{-3}-1)\}.
\]
Use $(L,\overline H)=(19/12,241/200)$,
$(63/40,241/200)$, and $(1559/1000,239/200)$ on
$(1,11/10]$, $[11/10,23/20]$, and $[23/20,6/5]$, respectively.
The right side decreases in $u\le1/16$; its derivative is
\[
 (L+18u/25)(36/25-L-54u/25)-(4/5)(1-u)^{-4}<0.
\]
At $u=1/16$ the three values are
$973799/51840000$, $-21361/3240000$, and
$-566237/16200000$.
The first is positive.  On the other intervals put
$A_1=10/9261$, $A_2=180/79507$, and let $d_1,d_2$ be the absolute
values of the last two displayed fractions.  Since
$c^2/(\lambda+1)^3$ increases for $1<\lambda\le6/5$,
\[\begin{gathered}
 P\ge A_j\rho_2^2,\qquad K_0>2A_2>2A_1,\\
 A_1-d_1/16=\frac{11873177}{17781120000}>0,\qquad
 A_2-d_2/16=\frac{1636194841}{20608214400000}>0.
\end{gathered}\]
For $0<\rho<1$,
$A_j\rho^2+K_0(1-\rho)\ge A_j$.
Equation~\eqref{eq:low-small-root-budget} proves $D_\lambda>0$.
This covers $q\le1/15$, since $2/\sqrt3<6/5$.

\subsubsection{Rational bounds for \texorpdfstring{$1/15<q\le10$}{1/15 < q <= 10}}
We use the tail bounds in Lemma~\ref{lem:low-Lambert-tails} to
form a finite set of rational sign tests.  Their definitions are given
here so that the arithmetic supplement contains only polynomial
expansion and exact coefficient checks.

For $q,c>0$ set $\lambda=1+c$, $d=1-cq$, $u=q/(q+1)$ and
$\theta_g=1-1/[6(q+1)]$. At $i=1$, use
$b=d/(q+1)$ in~\eqref{eq:low-tail-rationals}, and define
\begin{equation}\label{eq:low-theta-enclosures}
 T_{\rm lo}=u^2\{1+T_L(a_L)\},\quad
 T_{\rm up}=u^2\{1+T_U(a_U)\},\qquad
 \theta_{\rm lo,up}=
 \frac{T_{\rm lo,up}}{q\{2q-(\lambda+1)T_{\rm lo,up}\}}.
\end{equation}
In the rational expressions that follow, $R=q\theta$ and
$\alpha_p$ is given by~\eqref{eq:low-reference-calibration}, with
$\theta$ treated as a variable. For variables $b$ and $i$, put
\begin{align}
 r&=(\lambda+b)^{-1},\qquad r_{\rm prev}=(1-b)/\lambda,\nonumber\\
 \tau&=1-(\alpha_p+\lambda R)r+(\lambda+1)R-R/r_{\rm prev},
 \nonumber\\
 \mathcal B(i,b,\theta)&=\frac{2i(\alpha_p+cR)}\lambda
       +\frac{2iRrb}\lambda+2(i+1)Rr^2b
                              +\frac{a_L}{\lambda}(2i-1)\tau,\nonumber\\
 K_H(i,b,\theta)&=\mathcal B(i,b,\theta)-2V_U,\nonumber\\
 K_J(i,b,\theta)&=V_L+\frac{a_L}{\lambda}(2i-1)\alpha_pr
 -\frac{a_U R}{2\lambda}
       \{(2i-1)(1/r_{\rm prev}-\lambda r)+2\lambda r\},
 \label{eq:low-age-rational-tests}\\
 i_{\rm lo}(b)&=\frac{2(d-qb)}{2b+c(d-qb)}.\nonumber
\end{align}
Here $V_U,V_L$ are exactly~\eqref{eq:low-V-rational-bound}, evaluated
at the free $i,b,\theta$.  Both $K_H$ and $K_J$ are affine in $i$
and in $\theta$ separately.

We also need a lower bound for the terms at $i=1$ and the term
$C_{\rm cal}e$ in~\eqref{eq:low-master}. In the next display all tail
functions are at $i=1$, $b=d/(q+1)$, and
$r=(q+1)/(q+2+c)$.  Put
\begin{align*}
 V_0&=-T_U(a_U)+P_L(a_L),\\
 V_R&=-cT_U(a_U)+(\lambda+1)P_L(a_L)-4\lambda r^2-2r/(b+c),\\
 W_0&=c/\lambda-1+2b/\lambda,\qquad
 W_p=cW_0+4b/\lambda
          -2r\{b/\lambda+2d/(q+2+c)\},\\
 \alpha_0&=\{\lambda^2q+d(\lambda+1)\}/(q+1),\qquad
 \alpha_1=-d(\lambda+1)/(q+1),\\
 B_{10}&=2\alpha_0/\lambda,\qquad
 B_{11}=2(\alpha_1+cq)/\lambda+2qrb/\lambda+4qr^2b,\\
 E_v&=d/(q+1)^2,\qquad
 \varepsilon=u-\lambda\beta_2,\qquad
 \phi=\frac{d-c^2(q+2+c)/6}{(q+2+c)(2q+3+c)},\\
 C_0&=\frac{E_v(2-E_v)}{2q}
                 \{2u\phi+(2-h)\varepsilon/\lambda\}.
\end{align*}
Define, for $0\le e<1$,
\begin{align}
 \mathcal A(e,\theta)={}&(1-e)^2(V_0+q\theta V_R)
 +W_0+q\theta W_p-2\theta d/\lambda+e(B_{10}+\theta B_{11})
 \nonumber\\
 &+e^2\left\{\frac{2\{\alpha_0+\theta(\alpha_1+\lambda q)\}}
                     {\lambda(1-e)}
                         -\frac32a_Uq\theta r\right\}
       +C_0\{\theta(2-h)-1\}e.\label{eq:low-root-rational}
\end{align}
Use $U_j,L_j$ from~\eqref{eq:low-power-enclosures} to define
\begin{align}
 K_4={}&\frac{q^2(q+1)^2}\lambda
       \{U_3-qL_4+c(U_2-2qL_3+q^2U_4)\},\nonumber\\
 K_6={}&\frac{q^4(q+1)^2}\lambda
       \{U_5-qL_6+c(U_4-2qL_5+q^2U_6)\},\nonumber\\
 K_7={}&\frac{q^5(q+1)^2}\lambda
       \{U_6-qL_7+c(U_5-2qL_6+q^2U_7)\},\nonumber\\
 K_{\rm sq}={}&\frac{q^4(q+1)^2}\lambda
       \{U_4-2qL_5+q^2U_6+c(U_3-3qL_4+3q^2U_5-q^3L_6)\}.
 \label{eq:low-rational-costs}
\end{align}
For $s=dz/(q+1)$, $e=s^2$, $0\le z\le1$, define
\begin{align}
 \mathcal G_1(\theta)&=\mathcal A(e,\theta)
                     -\tfrac32e(1-e)(1-s)K_4,\nonumber\\
 \mathcal G_2(\theta)&=\mathcal A(e,\theta)
       -4e(1+s)\left\{\left(1-\frac{qs}{1-s}\right)K_6
                            +\frac{qs}{1-s}K_7\right\}
       -8e^2K_{\rm sq}.\label{eq:low-finite-budgets}
\end{align}
The factor $1-s$ in the denominator of $\mathcal G_2$ is strictly
positive throughout the specified domain.

\begin{lemma}[Rational sign tests for $0<q\le3$ and $3\le q\le10$]
\label{lem:low-Q12-arithmetic}
The following finite assertions hold.
\begin{enumerate}
\item On $0<q\le3$, $0\le c\le1/6$, and also on
$q=1/x$, $c=hx$, $1/10\le x\le1/3$, $0\le h\le1/2$,
\[
 0<T_{\rm lo}\le T_{\rm up}
 <\frac{2q^2(6q+5)}{6(q+1)+(\lambda+1)q(6q+5)}.
\]
The tail expressions at $c=0$ mean their rational continuous
extensions.
\item For $0<q\le3$, $0\le c\le1/6$, $0\le v_0\le1$, and
$i=2+j$, $j\ge0$, put $b=dv_0/(q+i)$, $r=(\lambda+b)^{-1}$ and
\[
 K_R=icr^2+2ibr^3-\frac{(2i-1)\lambda}{2(1-b)}
                         +\frac{(2i-3)\lambda r}{2}.
\]
Then $K_R<0$ and
\[
 i\lambda r^2+(2i-1)\lambda u r+qK_R>0.
\]
For $b=dv_0/(q+2)$ and $\theta\in\{0,\theta_g\}$ one also has
\[
 \partial_iK_H>0,\qquad
 K_H(i_{\rm lo}(b),b,\theta)+\tfrac32(qb/d)^2>0.
\]
The latter rational expressions are used when $c>0,v_0>0$;
their cleared polynomials include the closed boundary.
\item On $q=1/x$, $c=hx$, $1/10\le x\le1/3$,
$0\le h\le1/2$, $b=dv_0/(q+2)$, $0\le v_0\le1$,
at each $\theta\in\{0,\theta_g\}$,
\[
 \partial_iK_H\ge0,\quad K_H(i_{\rm lo},b,\theta)+4(qb/d)^4\ge0,
 \qquad
 \partial_iK_J\ge0,\quad K_J(i_{\rm lo},b,\theta)+8(qb/d)^4\ge0.
\]
\item For $1/29\le q\le3$, $0\le c\le1/6$, $0\le z\le1$,
both $\mathcal G_1(\theta_{\rm lo})$ and
$\mathcal G_1(\theta_{\rm up})$ are strictly positive.
For $q=1/x$, $c=hx$, $1/10\le x\le1/3$,
$0\le h\le1/2$, $0\le z\le1$, both
$\mathcal G_2(\theta_{\rm lo})$ and
$\mathcal G_2(\theta_{\rm up})$ are strictly positive.
\end{enumerate}
\end{lemma}

All the tests are on rational functions explicitly defined above.
Their denominators are products of positive numerical constants and
$\lambda$, $q$, $q+1$, $q+2$, $q+2+c$, $2q+3+c$, $a_L$, $a_U$,
$1-b$, $\lambda+b$, $b+c$, $2b+c(d-qb)$, $1-e$, and $1-s$.
The speed endpoint denominators are positive by part~(i).
After substituting $b=dv_0/(q+2)$, the only additional vanishing
boundary factors are $dv_0+c(q+2)$ and
$(2-cq)v_0+c(q+2)$. Both are strictly positive when
$c>0$ and $v_0>0$, as required for clearing denominators.

For specificity, the finite coefficient operation used to check this
lemma is as follows.  Map each bounded variable to $y\in[0,1]$ by
the affine maps in the lemma.  If
$P(y_1,\ldots,y_d)=\sum_\nu p_\nu y^\nu$ has coordinate degrees
$D_1,\ldots,D_d$, its tensor Bernstein coefficients are
\begin{equation}\label{eq:low-Bernstein-transform}
 b_\kappa=\sum_{\nu\le\kappa}p_\nu
             \prod_{j=1}^d\frac{\binom{\kappa_j}{\nu_j}}
                                  {\binom{D_j}{\nu_j}}.
\end{equation}
The unbounded variable $j=i-2$ in part~(ii) stays in the ordinary
power basis.  Positive coefficients in those variables imply the
claimed sign because their basis functions are nonnegative and sum
to one in each bounded variable.  For the exceptional one-variable
groups, the arithmetic supplement gives the polynomial and its
Euclidean Sturm sequence: positive endpoint values and equal numbers
of endpoint sign variations imply positivity throughout $[0,1]$.
There are fifteen degree-seven groups in part~(iv) for
$\mathcal G_1$, and twenty-five degree-ten groups in part~(iii) for
$K_H$.  Vanishing factors at zero or one are removed explicitly
before applying Sturm's theorem.  The remaining part~(iv) tests have
$10150$ and $10850$ positive tensor coefficients.

We now prove that this finite lemma suffices.  Part~(i) and the
tail bounds for $\sum_{s>1}w_s$ imply
\[
 0<\theta_{\rm lo}\le\theta\le\theta_{\rm up}<\theta_g<1.
\]
For every integer $i\ge2$, Bernoulli's inequality gives
$b=cd/(\lambda^i-d)\le d/(q+i)$.  Also
$\lambda^i=d(1+c/b)$ and
$\log\lambda\le c$, $\log z\ge2(z-1)/(z+1)$ for $z\ge1$
give $i\ge i_{\rm lo}(b)$.  In the definition of $B_i$ its
$k$ coefficient is nonnegative, because $\tau_i>0$.
Hence $S_i/w_i\ge K_H(i,b,\theta)$.
For the square coefficient, expand~\eqref{eq:low-J} as
\[
 J_i/w_i=V_i/w_i+k(2i-1)\alpha_pr
       -\frac{kR}{2}\{(2i-1)(1/r_{\rm prev}-\lambda r)+2\lambda r\}.
\]
The bracket is positive.  Since $a_L/\lambda\le k\le a_U/\lambda$,
$J_i/w_i\ge K_J(i,b,\theta)$. Affinity in $\theta$ and the
nonnegative derivatives with respect to $i$ in the lemma give
\begin{equation}\label{eq:low-all-age-prices}
 \begin{cases}
 S_i\ge-\tfrac32w_i\beta_i^2,\quad J_i>0,&q\le3,\\
 S_i\ge-4w_i\beta_i^4,\quad J_i\ge-8w_i\beta_i^4,&3\le q\le10,
 \end{cases}\qquad i\ge2.
\end{equation}
For $q\le3$, positivity of $J_i$ follows from the lower bound in
\eqref{eq:low-positive-mixture},
\[
 V_i/w_i>i\lambda r_i^2
                      \{1+R(c/\lambda+2L_{i+1})\}.
\]
Indeed $A_N\ge\lambda^{1-N}$ and
$P_N\ge cN\lambda^{-N}$, and summing these two geometric
expectations gives the display.  Use $k\ge1/\lambda$ on the
positive $\alpha_p$ term and $k\le1$ on its negative bracket.
The result is greater than
$i\lambda r^2+(2i-1)\lambda ur+qK_R$, since
$\alpha_p>\lambda^2u$, $R<q$, and $K_R<0$.
Part~(ii) proves positivity.

To bound the negative terms in~\eqref{eq:low-master}, we use
the representation of all the $\Delta_i$ by the same measure $\sigma$. For
$p=2$ or $4$, the sequence
$m_i=w_i\beta_i^p$ is a positive mixture of powers:
\[
 m_i=\frac{h^{p+2}}\lambda
       \sum_{N\ge p+1}\binom{N}{p+1}d^{N-p-1}\lambda^{-Ni}.
\]
For $0\le t<1$, the function
$f_t(z)=z/(1-t+t\sqrt z)^2$ is concave on $[0,1]$, since
\[
 f_t''(z)=-\frac{3(1-t)t}
                       {2\sqrt z(1-t+t\sqrt z)^4}\le0
 \quad(z>0).
\]
In the integral representation of $\Delta_i$ take $\xi$ to be the
integration variable and set $z=(1-\xi)^2$. Sum over $i\ge1$ and
apply Jensen's inequality under $\sigma$ to each term of the mixture.
Subtracting the $i=1$ term gives
\begin{equation}\label{eq:low-marker-damping}
 \sum_{i\ge2}m_i\Delta_i
 \le e\sum_{i\ge2}i m_i(1-\sqrt e)^{i-1}.
\end{equation}
The sums are finite or positively dominated geometric sums, so
Tonelli justifies both operations.  Since $\rho_i\le1-e$ for
$i\ge2$, this bounds the negative part of the linear terms with $i\ge2$ in
\eqref{eq:low-master}.

For $q\le3$, discard $(1-s)^{i-2}\le1$ in
\eqref{eq:low-marker-damping}.  Bernoulli and monotonicity of
$y/(y-d)^{p+2}$ give
\[
 w_i\beta_i^p\le
        \frac{q^{p+2}(1+ci)}{\lambda(q+i)^{p+2}}.
\]
Expand $i=(q+i)-q$ and use~\eqref{eq:low-power-enclosures}, with
upper bounds on positive coefficients and lower bounds on negative
coefficients.  For $p=2$ this gives
$u^{-2}\sum_{i\ge2}i w_i\beta_i^2\le K_4$.
Thus the sum of the negative parts of $\Delta_i\rho_i S_i/u^2$ for $i\ge2$ is at most
$\tfrac32e(1-e)(1-s)K_4$.

For $3\le q\le10$, put $t_*=qs/(1-s)\le qd/(q+h)\le1$,
using~\eqref{eq:low-e-caps}.  Bernoulli and the convex chord
\[
 (1-t)+\frac{t}{1+y}-\frac1{1+ty}
       =\frac{t(1-t)y^2}{(1+y)(1+ty)}\ge0
\]
give
\[
 (1-s)^{i-1}\le\frac1{1-s}
       \left\{(1-t_*)+\frac{t_*q}{q+i}\right\}.
\]
Use the preceding Bernoulli weight bound with $p=4$.
Expanding $i$ and $i^2$ in powers of $q+i$, and applying
\eqref{eq:low-power-enclosures} term by term, gives the
bounds
\[
 \frac1{u^2}\sum_{i\ge2}i w_i\beta_i^4\le K_6,
 \quad\frac1{u^2}\sum_{i\ge2}\frac{qi w_i\beta_i^4}{q+i}\le K_7,
 \quad\frac1{u^2}\sum_{i\ge2}i^2w_i\beta_i^4\le K_{\rm sq}.
\]
Equation~\eqref{eq:low-marker-damping}, $\Delta_i\le i e$, and
\eqref{eq:low-all-age-prices} now bound the sums of the negative parts
of $\Delta_i\rho_i S_i/u^2$ and $\Delta_i^2J_i/u^2$, for $i\ge2$, by
$4e(1+s)\{(1-t_*)K_6+t_*K_7\}$ and $8e^2K_{\rm sq}$,
respectively.

It remains to bound the terms at $i=1$ and the term $C_{\rm cal}e$.
Combine the terms $F_{\rm pipe}+eS_1+e^2V_1$ before estimating:
\[
 F_{\rm pipe}+eS_1+e^2V_1
 =(1-e)^2V_1+u^2W_1(R)-2RuL_1+eB_1.
\]
The lower bound for the tail $V_1$ gives
$V_1/u^2\ge V_0+q\theta V_R$, and
$W_2^Y/(2u^2)=3k\lambda Rr_1/2\le3a_U Rr_1/2$.
To estimate $C_{\rm cal}$, replace $e$ by $E_v$ only in a
componentwise lower bound for the positive coefficient
$C_{\rm head}(2-e)$.  At $a_2=\beta_2(1-E_v)$,
\[
 A_0=E_v,\qquad
 2-(3\lambda+1)u+\lambda(\lambda+1)a_2=(2-h)E_v.
\]
The remaining factors are bounded below by
$u-\lambda\beta_2$ and
\[
 -k\log(\beta_2/u)-1+\beta_2/u\ge\phi>0.
\]
For the last inequality use
$-\log z\ge2(1-z)/(1+z)$ and $k\ge a_L/\lambda$, then insert
$\beta_2/u=(q+1)/(q+2+c)$.  This yields exactly the displayed
$\phi$. Its numerator is positive in both ranges of $q$:
for $q\le3$ this follows from $c\le1/6$, while for
$3\le q\le10$, $c=h/q$ gives
$c^2(q+2+c)\le31/216<6d$.
Finally, the identity
\[
 \frac{\kappa-1}{Z_p}=\frac{\theta(2-h)-1}{2q}>0
\]
gives $C_{\rm cal}/u^2\ge C_0\{\theta(2-h)-1\}$.
Thus~\eqref{eq:low-master} proves
\[
 F_{\rm floor}(v)/u^2>\mathcal G_1(\theta)\quad(q\le3),
 \qquad
 F_{\rm floor}(v)/u^2>\mathcal G_2(\theta)\quad(3\le q\le10).
\]
Both $\mathcal G_1$ and $\mathcal G_2$ are affine in $\theta$, so
part~(iv) of Lemma~\ref{lem:low-Q12-arithmetic} gives positivity
throughout $[\theta_{\rm lo},\theta_{\rm up}]$ in the respective
ranges of $q$. Since $1/15>1/29$, this covers $1/15<q\le10$
when $0<\alpha<1$.

\subsubsection{\texorpdfstring{Bounds for $F_{\rm pipe}$, $S_i$, and $J_i$ when $q\ge10$}{Bounds for F\_pipe, S(i), and J(i) when q >= 10}}
For $q\ge10$ we use the simpler bound
\begin{equation}\label{eq:low-large-master}
 F_{\rm floor}(v)>F_{\rm pipe}-eH_{\rm head}-e^2J_{\rm minus},
 \quad
 H_{\rm head}=\sum_{i\ge1}i[-S_i]_+,\quad
 J_{\rm minus}=\sum_{i\ge2}i^2[-J_i]_+.
\end{equation}
This follows from~\eqref{eq:low-master}, $\rho_i\le1$,
$\Delta_i\le ie$, and positivity of the coefficient of $e^2$ in braces
and of $C_{\rm cal}e$.
The next estimates will in particular show that the two negative-part
sums have finite support.

Put $x=1/q$, $c=hx$, $\lambda=1+hx$ and $g=q-R$.
We have $0<g<1$. Indeed, summing the lower
bound $z/(e^z-1)\ge1-z/2$ over the geometric series gives
\[
 T_p\ge\frac{h}{\lambda\log\lambda}-\frac1{2\lambda}
       >\frac{2q}{\lambda+1}-\frac1{2\lambda}.
\]
Here $k>2/(\lambda+1)$ follows by strict trapezoid comparison for
$\int_1^\lambda t^{-1}\dd t$.
Subtract the value $2q(q-1)/[(\lambda+1)q-\lambda]$ corresponding
to $R=q-1$.  The difference is
\[
 \frac{\lambda(\lambda+1)-qc^2}
 {2\lambda(\lambda+1)((\lambda+1)q-\lambda)}>0,
\]
since $qc^2=h^2/q\le1/40$.  The quotient is increasing in $T_p$,
so $R>q-1$; the upper bound $R<q$ was already proved.

We express $F_{\rm pipe}$ in terms of $T_p$ and its derivative,
\begin{equation}\label{eq:low-whole-source}
 F_{\rm pipe}=\frac{2q}{Z_p}\{-cT_p'-2T_p+T_p^2/q\},
 \qquad
 T_p=\frac{h^2}\lambda
                  \sum_{N\ge1}\frac{N d^{N-1}}{\lambda^N-1}.
\end{equation}
The derivative holds $h,d$ fixed.  To verify the identity, sum
$W_i^0$ to obtain $F_0=-cT_p'-(1+c/\lambda)T_p$.
Shifting the $W_i^p$ terms, retaining their first terms, gives
\[
 \sum_iw_iW_i^p=(\lambda+1)F_0+2T_p-2u^2-2(q-u)/\lambda.
\]
Insert this into the definition of $F_{\rm pipe}$ and use
$L_1=1/\lambda-u$.  The result is~\eqref{eq:low-whole-source}.
The same geometric majorants used in Lemma~\ref{lem:low-Lambert-tails}
justify the derivative and the summations.

Define
\begin{equation}\label{eq:low-source-rationals}
 \begin{gathered}
 f(h)=\frac{(1-h)(2h+1)}6,\quad
 s(a)=a-x/2+\frac{x^2(2-h)}{12a},\\
 s_l(a)=s(a)-\frac{x^4(24-36h+14h^2-h^3)}{720a^3},\\
 \begin{aligned}
 P_M(a)={}&2\{s(a)^2-(\lambda+1)s(a)+a^2-x^2(2-h)/12\}\\
 &-Mx\lambda\{2\lambda-(\lambda+1)s_l(a)\}.
 \end{aligned}
 \end{gathered}
\end{equation}
Applying~\eqref{eq:low-global-F}--\eqref{eq:low-global-C}
for each $N\ge1$ gives
\[
 \frac{s_l(a)}\lambda\le\frac{T_p}{q}\le\frac{s(a)}\lambda,
 \qquad
 K_{\rm actual}:=\frac{ch^2}{\lambda^2\delta^2}
       \sum_{N\ge1}d^{N-1}
         \frac{(N\delta)^2}{2(\cosh(N\delta)-1)}
 \ge\frac{a^2}{x\lambda^2}-\frac{x(2-h)}{12\lambda^2}.
\]
Writing $t_T=T_p/q$, equation~\eqref{eq:low-whole-source} becomes
\[
 F_{\rm pipe}=
 \frac{2\{t_T^2/x-(1+1/\lambda)t_T/x+K_{\rm actual}\}}
                 {2-(\lambda+1)t_T}.
\]
Since $s(a)\le1-x/4+x^2/6<1$, the quadratic in $t_T$ decreases
on the whole required interval.  Its lower bound uses $s(a)/\lambda$,
and its denominator is at most $2-(\lambda+1)s_l(a)/\lambda$.
For either positive target $M=f(h)-2x/5$ or $M=69/500$,
$P_M(a)>0$ therefore implies $F_{\rm pipe}>M$: the inequality
first proves positivity of the lower numerator, so replacing the
denominator preserves the lower bound. For $a_L\le a\le a_U$,
$s\ge19/20$, $\lambda+1\le41/20$, $0<s'<1$, and $s_l'>0$.
Consequently $P_M'(a)>37/10$.  It suffices to check $P_M(a_L)>0$.

To bound $S_i=B_i-2V_i$, set $t=i/q$, let
$z=d\lambda^{-i}$, and put $w=hz/(1-z)$.  Bernoulli gives
$0<w\le d/(1+t)$.  For a free point in this larger domain define
\begin{align}
 b&=xw,\quad r=(\lambda+b)^{-1},\quad r_{\rm prev}=(1-b)/\lambda,
 \nonumber\\
 R&=1/x-g,\qquad
 \alpha_p=\frac{\lambda^2+gd(\lambda+1)x^2}{1+x},
 \nonumber\\
 A&=1-b+(cb+b^2)/2,\nonumber\\
 P&=c+2b-4cb-5b^2+\tfrac92c^2b+\tfrac{25}2cb^2+8b^3,\nonumber\\
 C&=5b/3+cz/3,\nonumber\\
 U&=\{i(A+RP)+z/2+RC\}/\lambda,\nonumber\\
 \mathcal B&=2i(\alpha_p+cR)/\lambda+2iRrb/\lambda
             +2(i+1)Rr^2b\nonumber\\
 &\quad+\frac2{\lambda+1}(2i-1)
                \{1-(\alpha_p+\lambda R)r+(\lambda+1)R-R/r_{\rm prev}\},
 \nonumber\\
 \operatorname{shape}&=2ht+4(t-1/3)w-(2h+3)z/3,
 \qquad z=w/(h+w).\label{eq:low-large-price-rationals}
\end{align}
At $t=i/q$, $z=d\lambda^{-i}$ and $w=hz/(1-z)$, we have
$B_i/w_i\ge\mathcal B$ and $V_i/w_i\le U$.
The first follows from $k\ge2/(\lambda+1)$ and $\tau_i\ge0$.
For the second bound, the coefficients in
\eqref{eq:low-positive-mixture} satisfy, for every integer $N\ge1$,
\begin{align*}
 A_N&\le1-(N-1)c/2+(N^2-1)c^2/12,\\
 P_N&\le Nc-(N-1)(5N+2)c^2/6+(N^2-1)(4N+1)c^3/12,\\
 B_N&\le(N-1)/(2N),\qquad
 C_N\le c(N-1)(5N+2)/(6N).
\end{align*}
For the first bound, multiplication by $(1+c)^N-1$ leaves, at
power $c^j$, $j\ge4$, the nonnegative coefficient
\[
 \binom{N}{j-2}
 \frac{(j-3)(N+1)\{(j-4)N+5j-8\}}{12j(j-1)};
\]
the lower powers cancel.  The second bound follows by inserting the
first in $P_N$ and using the global third-order Taylor lower bound
for $(1+c)^{1-N}$, whose fourth derivative is nonnegative.
For the third, put
$D_N=(N-1)Q_N^2-2N\sum_{j=1}^{N-1}j\lambda^j$.
One has $D_N(1)=0$,
$D_N'(1)=N^2(N-1)(N-2)/3\ge0$.
The coefficients of $D_N'$ are negative only at degrees at most
$N-2$ and positive above that degree: the coefficient at degree
$j-1$ is
$j\{(N-1)\min(j+1,2N-1-j)-2Nj\mathbf1_{j<N}\}$.
Thus $D_N'/\lambda^{N-2}$ increases on $[1,\infty)$ and
$D_N\ge0$, proving the $B_N$ bound.
For the last bound we use the polynomial identity
\begin{align*}
 &(\lambda+1)N^2\lambda^{2N-1}-cN\lambda^NQ_N-(2+2Nc)Q_N^2
       =c^2\sum_{j=0}^{2N-2}a_j\lambda^j,\\
 &a_j=(j+1)(j+2)(N-1-j/3)\quad(0\le j\le N-2),\\
 &a_{2N-2-r}=(N-r-1)\{(2r+1)N+r(r+1)\}+2r(r+1)(r+2)/3,\\
 &\hspace{48mm}0\le r\le N-1.
\end{align*}
For $N=1$ both sides vanish.  For $N\ge2$ all $a_j\ge0$;
finite second differences, with the four joining coefficients,
verify the identity.  Their sum is
$N^2(N-1)(5N+2)/6$.  Since
$C_N=c\lambda^{1-N}\sum_j a_j\lambda^j/(NQ_N^2)$ and
$Q_N\ge N\lambda^{(N-1)/2}$, the required upper bound follows.
Summing all four inequalities under $\Pi$ gives exactly the
$A,P,z/2,C$ in~\eqref{eq:low-large-price-rationals} for every $i\ge1$.

To estimate the negative parts of $J_i$, we use one further rational function. For
$q\ge10$, $0<h\le1/2$ define
\begin{equation}\label{eq:low-AJ-rational}
 \begin{gathered}
 a_m=\frac{d}{\lambda(q+2)},\qquad
 r_m=\frac1{\lambda+d/(q+2)},\\
 N_m=\frac{(\lambda a_m+c)
       \{2\lambda^2a_m^2+(3\lambda^2+2\lambda+1)a_m+\lambda^2-1\}}
                 {\lambda^2(1-\lambda a_m)},\\
 A_J^{\rm lo}=\lambda r_m^2+
                   \frac{4\lambda^2u r_m}{\lambda+1}-qN_m.
 \end{gathered}
\end{equation}

\begin{lemma}[Rational inequalities for $q\ge10$]
\label{lem:low-Q3-arithmetic}
For $0<x\le1/10$, $0\le h\le1/2$,
$P_{f(h)-2x/5}(a_L)>0$.
For the domain
\[
 0<x\le1/10,\quad0<h\le1/2,\quad t\ge0,\quad
 0<w\le(1-h)/(1+t),\quad0\le g\le1,
\]
the rational functions in~\eqref{eq:low-large-price-rationals} satisfy
\begin{equation}\label{eq:low-Q3-price-check}
 \mathcal B-2U\ge\operatorname{shape}-\frac{28}{5}x(1+ht).
\end{equation}
Finally, $A_J^{\rm lo}>3-12/q$ for $q,h>0$ with $h\le1/2$.
\end{lemma}

These checks have the following finite specifications.
Put $A=12a_L$.  The polynomial
$1800A^3P_{f(h)-2x/5}(a_L)/x^3$ has constant coefficient $380160$;
the sum of the absolute values of its negative coefficients, weighted
by $10^{-j}2^{-k}$ for $x^jh^k$, is
$65345524702337817/204800000000<320000$.
For~\eqref{eq:low-Q3-price-check}, multiply
$[\mathcal B-2U-\operatorname{shape}]/x+28(1+ht)/5$ by the
positive denominator
\[
 3(h+w)(1+x)\lambda(\lambda+1)(1-xw)(\lambda+xw)^2.
\]
Substitute $w=(1-h)v_0/(1+t)$ and multiply by $(1+t)^7$.
The result is a polynomial of degrees $(6,8,8,7)$ in
$(x,h,t,v_0)$, affine in $g$.  At $g=0,1$, transform only
$x,h,v_0$ by~\eqref{eq:low-Bernstein-transform} on
$[0,1/10]\times[0,1/2]\times[0,1]$, keeping ordinary powers of
$t\ge0$.  Each endpoint has $4319$ nonzero positive coefficients,
with minimum $69/280$; the remaining coefficients are zero.
For~\eqref{eq:low-AJ-rational}, the numerator of
$A_J^{\rm lo}-(3-12/q)$ over
\[
 q(q+h)^3(2q+h)(q+1)(q+2)^2(q+h+1)(q^2+3q+2h)^2
\]
has $60$ nonnegative nonzero coefficients in $q,h$, including
$66q^{11}$.  These specifications prove the three statements by
finite rational arithmetic, with every denominator positive.

The first assertion gives
\begin{equation}\label{eq:low-Q3-source}
 F_{\rm pipe}>f(h)-2/(5q).
\end{equation}
To apply~\eqref{eq:low-Q3-price-check} to $S_i/w_i$, we first
show that for $t=i/q$, $z=d\lambda^{-i}$ and $w=hz/(1-z)$,
\begin{equation}\label{eq:low-Q3-shape}
 \operatorname{shape}\ge\frac{9t-7}{3(1+t)}+\frac34ht.
\end{equation}
For $t\le1/3$, use $w\le d/(1+t)$ and $z\le d$; the difference
in~\eqref{eq:low-Q3-shape} is at least
$h\{5t/4-4(t-1/3)/(1+t)+1/3\}+2h^2/3\ge0$.
For $t\ge1/3$ and $y=ht<3$, the exponential series, using
$n!\ge2\cdot3^{n-2}$ for $n\ge2$, gives
$e^y\le1+y+y^2/[2(1-y/3)]$.
Since $\lambda^i\le e^y$ and $\lambda^i\ge1+y$,
\[
 w\ge\frac{d(1-y/3)}{1+t-y/3+yt/6},\qquad z\le\frac d{1+y}.
\]
After insertion, the difference is at least
\[\begin{aligned}
 &\frac{h}{12(t+1)(y+1)((y+6)t+6-2y)}\\
 &\quad\times\Bigl\{t^3(15y^2-27y+18)+t^2(81y^2+13y+108)\\
 &\hspace{19mm}+t(42y^2-48y+18)+120+152y-40y^2
                  +\frac{48y-16y^2}{t}\Bigr\}>0.
\end{aligned}\]
The first and third quadratic factors have negative discriminants,
the second has positive coefficients, the constant quadratic has
positive endpoint values on $[0,3]$, and the last term is
nonnegative.  The case $y=3$ follows by continuity.
For $y\ge3$, $t\ge6$; drop the nonnegative $w$ term and use
$(2h+3)d\le3$ to get
$\operatorname{shape}\ge2y-1/(1+y)>3+3y/4$.
This proves~\eqref{eq:low-Q3-shape} in every case.  Thus, with
$\eta=28/(5q)$,
\begin{equation}\label{eq:low-Q3-head-envelope}
 \frac{S_i}{w_i}\ge\frac{9t-7}{3(1+t)}-\eta
                                  +ht(3/4-\eta).
\end{equation}
Since $\eta\le14/25<3/4$, $S_i<0$ is possible only if
$t<(7+3\eta)/(9-3\eta)$.

For $i\ge2$, the lower bound from the mixture and the expression
for $J_i$ give
\[
 J_i/w_i>iA_J-B_J,
 \quad A_J=\lambda r_i^2+2k\alpha_pr_i
 +R\{cr_i^2+2\lambda r_i^3L_i+k(\lambda r_i-1/r_{i-1})\},
\]
where $B_J=k\{C_\alpha r_i-R(1/r_{i-1}-\lambda r_i)/2\}<q+1$.
To bound $A_J$, set $a=L_i$.  At $k=1$ its bracket is exactly
\[
 -\frac{(\lambda a+c)
       \{2\lambda^2a^2+(3\lambda^2+2\lambda+1)a+\lambda^2-1\}}
       {\lambda^2(1+a)^3(1-\lambda a)}.
\]
Since $k\le1$, this is a lower bound for that bracket.
Bernoulli gives $a\le a_m$, $r_i\ge r_m$ for $i\ge2$.
Increase every positive numerator factor to $a_m$, omit the factor
$(1+a)^3\ge1$, and use $k\ge2/(\lambda+1)$,
$\alpha_p>\lambda^2u$, and $R<q$.  The result is
$A_J\ge A_J^{\rm lo}>3-12/q$ by Lemma~\ref{lem:low-Q3-arithmetic}.
Also $w_i\le w_1=u^2<1$, since
$y/(1-dy)^2$ increases in $y=\lambda^{-i}$.
Hence a finite sum bounds $J_{\rm minus}$. For $T\ge2$, writing
$n=\lfloor T\rfloor$ and $r=T-n$,
Faulhaber's formulas give
\begin{align*}
 \sum_{i=1}^n i^2(T-i)
 ={}&T^4/12+T^2(-r^2/2+r/2-1/12)\\
 &+T(2r^3/3-r^2+r/3)-r^2(1-r)^2/4
 \le(T^4+T^2)/12.
\end{align*}
The coefficient of $T^2$ is at most $1/24$, that of $T$ is at most
$1/12$ in absolute value, and $T/12\le T^2/24$ for $T\ge2$.
Taking $T=(q+1)/(3-12/q)$ gives
\begin{equation}\label{eq:low-Q3-square-cost}
 J_{\rm minus}<q^4J(q),\qquad
 J(q)=\frac{(q+1)^4}{324q(q-4)^3}
                       +\frac{(q+1)^2}{36q^3(q-4)}.
\end{equation}
For every integer $i>T$ the bound on $J_i/w_i$ gives $J_i>0$.

\subsubsection{Bounds for \texorpdfstring{$q\ge21$ and for $q\ge10$ with $h\ge1/4$}{q >= 21 and for q >= 10 with h >= 1/4}}
We first give a quadrature bound that includes a support endpoint
lying inside a mesh interval.  If $G$ is continuous, compactly
supported in $[0,\infty)$, $G(0)=0$, and $G'$ has bounded variation
on $(0,\infty)$, then
\begin{equation}\label{eq:low-BV-quadrature}
 \left|\Delta\sum_{i\ge1}G(i\Delta)-\int_0^\infty G(t)\dd t\right|
 \le\frac{\Delta^2}{8}\operatorname{TV}(G').
\end{equation}
On a mesh interval $[a,b]$ use $K(s)=(s-a)(b-s)/2$ and integrate
twice by parts:
\[
 \frac{\Delta}{2}\{G(a)+G(b)\}-\int_a^bG(s)\dd s
                  =\int_{(a,b)}K(s)\dd G'(s).
\]
Derivative jumps contribute their point masses; $K$ vanishes at
both mesh endpoints.  Since $0\le K\le\Delta^2/8$, summing proves
\eqref{eq:low-BV-quadrature}.  The first and last values of $G$ are
zero, so the complete trapezoid sum equals the right-endpoint sum.

For $t=i/q$, $i\ge1$, monotonicity of $s/(s-d)^2$ and
$\lambda^i\ge1+ht$ give
\begin{equation}\label{eq:low-Bernoulli-weight}
 w_i\le\frac{1+ht}{\lambda(1+t)^2}.
\end{equation}
If $q\ge21$, use $\eta_A=4/15$, $A=39/5$, $B=41/5$ in
\eqref{eq:low-Q3-head-envelope}. Every index for which $S_i<0$ satisfies
$t<A/B=39/41<1$.
Put $D_h=1-(\lambda+1)h$.
On that support $D_h(1+ht)\le(1-2h)(1+ht)\le1$.
The bound on $eq^2$ in~\eqref{eq:low-e-caps} and
\eqref{eq:low-Bernoulli-weight} imply
\[
 eH_{\rm head}\le\frac1q\sum_{i\ge1}G_A(i/q),\qquad
 G_A(t)=\left[\frac{t(A-Bt)}{3(1+t)^3}\right]_+.
\]
Inside its support,
\[
 G_A'=\frac{A-32t+Bt^2}{3(1+t)^4},\qquad
 G_A''=\frac{-32-4A+(96+2B)t-2Bt^2}{3(1+t)^5}.
\]
The numerator of $G_A''$ increases across the support, is negative
at zero and positive at its end.  Thus $G_A'$ has one minimum,
then jumps from a negative value to zero at the end.
Moreover $G_A'>-1/2$: after multiplying by $6(1+t)^4$ the
required inequality is
\[
 17-52t+36t^2+12t^3+3t^4+6\eta_A(1-t^2)>0.
\]
For $t\le1/3$ its first three terms are at least $11/3$; for
$t\ge1/3$, use $12t^3\ge4t^2$ and
$17-52t+40t^2>0$, whose discriminant is $-16$.
The last group is nonnegative because $t<1$.
Consequently $\operatorname{TV}(G_A')<13/5+1<4$.
An elementary antiderivative gives
\[
 I_A:=\int G_A=\frac83-\frac{B^2}{96}-\frac B3\log(16/B).
\]
For $0<s<1$ the exact geometric remainder in
$\int_0^s(1-r^2)^{-1}\dd r$ gives
\begin{equation}\label{eq:low-log-four}
 \log\frac{1+s}{1-s}>\ell(s):=2(s+s^3/3+s^5/5+s^7/7).
\end{equation}
Thus, with
\[
 I_A^{\rm up}=175/96+(9/16)\eta_A-(3/32)\eta_A^2
                    -(3-\eta_A)\ell(39/121),
\]
equation~\eqref{eq:low-BV-quadrature} gives
$eH_{\rm head}<I_A^{\rm up}+1/(2q^2)$.
The same bound for $J_{\rm minus}$, with $3-12/q\ge17/7$, yields
\[
 e^2J_{\rm minus}<
 J_A:=\frac{(22/21)^4}{12(17/7)^3}
               +\frac{(22/21)^2}{12(17/7)21^2}
       =\frac{2904605}{409493637},
\]
because $e^2q^4<1$.
Combining these bounds with $f(h)\ge1/6$ in~\eqref{eq:low-Q3-source}
gives
\[
 F_{\rm floor}(v)>
 \frac16-\frac2{105}-I_A^{\rm up}-\frac1{882}-J_A
 =\frac{28155382207511429143103}{124404112403316879470013600}
 >\frac1{5000}.
\]

Next let $q\ge10$ and $h\ge1/4$.  Use $\eta_B=14/25$,
$A=217/25$, $B=183/25$. When $S_i<0$ we have $t<A/B<2$.
Retain the factor $D_h$ in the bound on $eq^2$ in~\eqref{eq:low-e-caps}. Equations
\eqref{eq:low-Q3-head-envelope} and~\eqref{eq:low-Bernoulli-weight} give
\[
 eH_{\rm head}\le\frac{D_h}{q}\sum_{i\ge1}G_B(i/q),\qquad
 G_B(t)=\left[\frac{t(A-Bt)(1+t/2)}{3(1+t)^3}\right]_+.
\]
Inside the support,
\[\begin{aligned}
 G_B'&=\frac{2A-(2A+4B)t-(A+B)t^2}{6(1+t)^4},\\
 G_B''&=\frac{-10A-4B+(4A+10B)t+2(A+B)t^2}{6(1+t)^5}.
\end{aligned}\]
The second derivative changes sign once from negative to positive
before $t=1$, so $G_B'$ again has one minimum and its last jump is
to zero.  The inequality $G_B'>-1$ reduces to
\[
 584/25-(566/25)t+20t^2+24t^3+6t^4>0;
\]
the quadratic part has discriminant $-847644/625<0$.
Therefore $\operatorname{TV}(G_B')<217/75+2<5$.
The exact integral is $I_B=(8/3)\log(400/183)-73563/40000$.
For $s_B=217/583$ the remainder of the geometric series gives
\[
 \log(400/183)<2\{s_B+s_B^3/3+s_B^5/[5(1-s_B^2)]\}.
\]
Together with~\eqref{eq:low-BV-quadrature} and $q\ge10$, this proves
\[
 eH_{\rm head}<D_h C_B,\qquad
 C_B=\frac{122243248244149}{483498900280000}>\frac1{12}.
\]
The bound for $J_{\rm minus}$ with $3-12/q\ge9/5$ gives
\[
 e^2J_{\rm minus}<D_h^2J_B,\qquad
 J_B=\frac{(11/10)^4}{12(9/5)^3}
       +\frac{(11/10)^2}{12(9/5)10^2}
       =\frac{187913}{8748000}.
\]
Since $D_h\le1-2h$, equations~\eqref{eq:low-large-master} and
\eqref{eq:low-Q3-source} imply
\[
 F_{\rm floor}(v)>f(h)-\frac1{25}-(1-2h)C_B-(1-2h)^2J_B.
\]
The right side has derivative
$(1-4h)/6+2C_B+4(1-2h)J_B>-1/6+2C_B>0$ on $[1/4,1/2]$.
At $h=1/4$ it equals
\[
 \frac{16616808991576441}{1057412094912360000}>\frac1{100}.
\]
This proves positivity for $q\ge10$ and $h\ge1/4$, using the same
$e$ in the bounds for $eH_{\rm head}$ and $e^2J_{\rm minus}$.

\subsubsection{\texorpdfstring{The coefficients $S_1,\ldots,S_{24}$ for $10\le q\le21$ and $h\le1/4$}{The first 24 coefficients S(i) for 10 <= q <= 21 and h <= 1/4}}
Put $x=1/q\in[1/21,1/10]$. Since $25/21>217/183$,
\eqref{eq:low-Q3-head-envelope} gives $S_i\ge0$ for $i\ge25$.
It remains to bound $S_i$ for $1\le i\le24$.

We first bound $g$ and $F_{\rm pipe}$. With
$s,s_l,P_M$ as in~\eqref{eq:low-source-rationals}, let
\[
 t_g(\gamma)=\frac{2(1-\gamma x)}{x+(\lambda+1)(1-\gamma x)}.
\]
For $0<x\le1/10$, $0\le h\le1/4$, the finite inequalities
needed are
\begin{equation}\label{eq:low-Q3c-source-checks}
 s_l(a_L)/\lambda>t_g(1/4),\qquad
 t_g(1/8)>s(a_U)/\lambda,\qquad P_{69/500}(a_L)>0.
\end{equation}
They imply $1/8<g<1/4$ and $F_{\rm pipe}>69/500$ by monotonicity
of $R$ in $T_p$ and the bound using $P_M$ preceding
Lemma~\ref{lem:low-Q3-arithmetic}.

We next define a polynomial for each of the remaining indices.
To form $F_i(c,x,g)$ take $\lambda=1+c$, $d=1-c/x$ and
$a_U=1+c/2-c^2/12+c^3/24$; substitute $c=hx$ at the end.  For each
integer $0\le n\le25$ put
\begin{equation}\label{eq:low-Q3c-Dn}
 D_n(c,x)=1+x\sum_{j=1}^n\binom nj c^{j-1}.
\end{equation}
The empty sum is zero.  When $c=hx$ this is
$(\lambda^n-d)/h=\beta_n^{-1}$, with its continuous value at
$h=0$.  For $1\le i\le24$ let $n=i+1$, $p=\lambda^n$,
$b_0=1/D_n$, and set
\begin{align*}
 m_0&=b_0,\qquad m_1=pb_0^2,\qquad m_2=p(p+d)b_0^3,\\
 m_3&=p(p^2+4dp+d^2)b_0^4,\\
 m_4&=p(p^3+11dp^2+11d^2p+d^3)b_0^5,\\
 r&=D_i/D_{i+1},\qquad r_{\rm prev}=D_{i-1}/D_i,\\
 L_i&=dx/(\lambda D_i),\qquad w_i=\lambda^{i-1}/D_i^2.
\end{align*}
For a free $a\ge1$, use the same $R=1/x-g$ and $\alpha_p$ as in
\eqref{eq:low-large-price-rationals} and define
\begin{align}
 \operatorname{TL}_i(a)&=
       \{am_0/x+m_1/2+x m_2/(12a)-x^3m_4/(720a^3)\}/\lambda,
 \nonumber\\
 \operatorname{PU}_i(a)&=\lambda^{-2}
 \Bigl\{n(am_1+x m_2/2+x^2m_3/(12a))+a^2m_0/x\nonumber\\
 &\hspace{30mm}-xm_2/12+x^3m_4/(240a^2)\Bigr\},
 \nonumber\\
 \operatorname{VU}_i(a)&=-(1+cR)\operatorname{TL}_i(a)
              +\{1+(\lambda+1)R\}\operatorname{PU}_i(a)\nonumber\\
 &\quad-2Rn\lambda^i/D_{i+1}^2-2R/(x\lambda D_{i+1}),
 \nonumber\\
 \operatorname{BB}_i(a)&=2i(\alpha_p+cR)/\lambda+2iRrL_i
       +2nR\lambda r^2L_i\nonumber\\
 &\quad+(a/\lambda)(2i-1)
              \{1-(\alpha_p+\lambda R)r+(\lambda+1)R-R/r_{\rm prev}\},
 \nonumber\\
 K(t,x)&=\frac{(3-9x)t-7/3+3x}{1+t},\nonumber\\
 F_i(c,x,g)&=\operatorname{BB}_i(a_U)
                   -2\operatorname{VU}_i(a_U)/w_i-K(ix,x).
 \label{eq:low-Q3c-head-rational}
\end{align}
Each $F_i$ is an explicit rational function. To clear its denominators, set
\begin{equation}\label{eq:low-Q3c-head-denominator}
 \mathcal D_i=720\lambda^{i+1}x^2a_U^3D_{i+1}^5
                                (1+x)D_{i-1}(1+ix).
\end{equation}
For $\gamma\in\{1/8,1/4\}$ define the polynomial
\begin{equation}\label{eq:low-Q3c-head-polynomial}
 \mathcal P_{i,\gamma}(x,h)
 =x^{-2}\left.\{\mathcal D_i F_i(c,x,\gamma)\}\right|_{c=hx}.
\end{equation}
For each $1\le i\le24$, the arithmetic lemma checks that the numerator
formed using~\eqref{eq:low-Q3c-Dn} is divisible by $x^2$.

\begin{lemma}[Rational inequalities for $0<x\le1/10$ and $0\le h\le1/4$]
\label{lem:low-Q3c-arithmetic}
The three inequalities~\eqref{eq:low-Q3c-source-checks} hold on
$0<x\le1/10$, $0\le h\le1/4$.
Each of the $48$ polynomials~\eqref{eq:low-Q3c-head-polynomial}
has nonnegative tensor Bernstein coefficients on
$[0,1/10]\times[0,1/4]$, and is strictly positive when $x>0$.
Consequently $F_i(hx,x,g)>0$ for $1\le i\le24$,
$0<x\le1/10$, $0\le h\le1/4$, and $1/8\le g\le1/4$.
\end{lemma}

To specify the three smaller checks by positive denominators, put
$A=12a_L$, $A_u=24a_U$,
$F_4=4x+(\lambda+1)(4-x)$ and
$F_8=8x+(\lambda+1)(8-x)$.
The cleared polynomials are
\[
 \frac{60\lambda A^3F_4}{x^2}\{s_l(a_L)/\lambda-t_g(1/4)\},\quad
 \frac{24\lambda A_uF_8}{x^2}\{t_g(1/8)-s(a_U)/\lambda\},\quad
 \frac{90000A^3}{x^2}P_{69/500}(a_L).
\]
Their degrees in $(x,h)$ are $(8,9)$, $(6,7)$, $(9,10)$; all
$90$, $56$, and $110$ Bernstein coefficients are positive.  Their
respective minima are $25920$, $672/5$, and $9325152/125$.
For $\mathcal P_{i,\gamma}$, the degrees are $(27,20)$ at $i=1$ and
$(7i+20,7i+12)$ at $2\le i\le24$.
The coefficient operation is~\eqref{eq:low-Bernstein-transform};
the $636112$ nonzero coefficients are positive. In each array the
only zero coefficient is the one at index $(0,0)$, which proves
strict positivity for $x>0$. Checking the
identities~\eqref{eq:low-Q3c-head-rational}--\eqref{eq:low-Q3c-head-polynomial}
and these finite arrays is the entire arithmetic obligation for this
lemma.  Positivity of $D_n$, $a_U$, $\lambda$, $x$ and $1+ix$
proves the sign of~\eqref{eq:low-Q3c-head-denominator}. The factor
$x^2$ is positive for $x>0$. Affinity of $F_i$ in $g$
then covers the whole interval between the two checked endpoints.

We now use the polynomial inequalities to bound $S_i/w_i$. At
$a=c/\log\lambda$, Lemma~\ref{lem:low-Lambert-tails} gives
$V_i\le\operatorname{VU}_i(a)$ and
$B_i/w_i=\operatorname{BB}_i(a)$.
We must change this one common $a$ in both terms together.
For $i\le24$, $c\le1/40$, and $n=i+1\le25$,
\[
 p=\lambda^n<e^{5/8}
 \le1+5/8+\frac{(5/8)^2}{2(1-5/24)}=\frac{569}{304}<2.
\]
The exponential estimate includes the full tail, by the factorial
bound used in~\eqref{eq:low-Q3-shape}.  Moreover
$p-1<2nc$, so $b_0>1/(1+2n/q)\ge1/6$, while $b_0\le1$.
The displayed moments then give $m_3/m_1\le13$ and
$m_4/m_0\le150$.
Differentiating only the explicit rational functions yields, for all
$a\ge1$,
\[
 \operatorname{TL}_i'(a)<\frac{101m_0}{100x\lambda},\qquad
 \operatorname{PU}_i'(a)>\frac{19m_0}{10x\lambda^2}.
\]
Indeed, the derivative of the first has numerator
$m_0/x-xm_2/(12a^2)+x^3m_4/(240a^4)$; in the second its
$n$ term is positive since $x^2m_3/(12a^2m_1)\le13/1200<1$.
The remaining fourth-order ratios are at most $1/16000$ and
$1/8000$, respectively.  As $R>q-1/4$, $cR\le1/4$ and
$\lambda\le41/40$, it follows that
\[
 \operatorname{VU}_i'(a)>
 \frac{m_0}{x\lambda^2}
 \left\{\frac{19}{10}(2q+1/2)
                 -\frac{101}{100}\frac{41}{40}\frac54\right\}
 >3q^2m_0.
\]
For the last inequality, subtract $3q(41/40)^2$ from the bracket;
the result is $(1037/1600)q-1101/3200\ge19639/3200>0$.
We proved $\tau_i<1$ and $w_i<1$.  Hence
\[
 \frac{\dd}{\dd a}
       \{\operatorname{BB}_i(a)-2\operatorname{VU}_i(a)/w_i\}
 <47-6q^2b_0/w_i\le47-q^2<0.
\]
Now $a\le a_U$ and Lemma~\ref{lem:low-Q3c-arithmetic} imply
\begin{equation}\label{eq:low-Q3c-actual-head}
 \begin{split}
 S_i/w_i&\ge\operatorname{BB}_i(a)-2\operatorname{VU}_i(a)/w_i\\
 &\ge\operatorname{BB}_i(a_U)-2\operatorname{VU}_i(a_U)/w_i
 >K(ix,x),\qquad1\le i\le24.
 \end{split}
\end{equation}
This bounds $S_i/w_i$ at every index where $S_i$ can be negative.
For $i\ge25$ we already have $S_i\ge0$ by~\eqref{eq:low-Q3-head-envelope}.

To bound $H_{\rm head}$, we sharpen the estimate for $w_i$. Let
$y=i\log\lambda$ and $\rho=(1-h/2)/a\le1$.  Directly,
\[
 w_i=\frac{h^2}{\lambda
       \{h\cosh(y/2)+(2-h)\sinh(y/2)\}^2}
 \le\frac1{\lambda(1+\rho i/q)^2}.
\]
Also
\[
 \lambda(1-h/2)^2-d(1+c/2)^2
                  =\frac{h^2}{4}(\lambda-dx^2)>0.
\]
Since $a\le1+c/2$, this proves $\lambda\rho^2>d$ and
$\lambda\rho>d$.  All coefficients of
$\lambda(1+\rho t)^2-d(1+t)^2$ are therefore positive, so
\begin{equation}\label{eq:low-Q3c-weight}
 d w_i\le(1+i/q)^{-2}\qquad(i\ge1).
\end{equation}
Apply $e\le d/(q+1)^2$ to~\eqref{eq:low-Q3c-actual-head}.
Since $S_i\ge0$ for $i\ge25$, this gives
\[
 eH_{\rm head}\le\frac{u^2}{q}\sum_{i\ge1}G_x(i/q),\qquad
 G_x(t)=\left[\frac{t\{(7-9x)-(9-27x)t\}}{3(1+t)^3}\right]_+.
\]
For $0\le x\le1/10$ the support ends at
$T=(7-9x)/(9-27x)<1$.  With $A=7-9x$, $B=9-27x$,
\[
 G_x'=\frac{A-2(A+B)t+Bt^2}{3(1+t)^4},\qquad
 G_x''=\frac{-6A-2B+(6A+8B)t-2Bt^2}{3(1+t)^5}.
\]
The second derivative changes sign exactly once on the support.
Furthermore $G_x'>-1/2$.  The required numerator is
\[
 17-52t+36t^2+12t^3+3t^4+18x(-1+8t-3t^2).
\]
For $t\le1/4$, the quadratic part is at least $25/4$ and the
last term at least $-9/5$.  For $1/4\le t\le1$, the last
term is nonnegative; the same split at $1/3$ used for $G_A$
proves positivity.  Including the jump at $T$ gives
$\operatorname{TV}(G_x')\le7/3+1=10/3$.
Thus~\eqref{eq:low-BV-quadrature} gives
\[
 q^{-1}\sum_{i\ge1}G_x(i/q)\le I(x)+1/240,\qquad
 I(x)=\frac{A(A+2B)}{6(A+B)}-\frac B3\log\frac{A+B}{B}.
\]
For each fixed $t$, $G_x(t)$ is the positive part of an affine
function of $x$, so its integral is convex on $[0,1/10]$.
Its maximum is bounded by its two endpoint values.  Use
\eqref{eq:low-log-four} at $s=7/25$ and $61/187$ to obtain
\begin{align*}
 I(0)&<175/96-3\ell(7/25)=\frac{56737364551}{585937500000},\\
 I(1/10)&<11407/7440-(21/10)\ell(61/187)
              =\frac{33081533096174793233}{297463843858303887600}.
\end{align*}
The second rational bound exceeds the first by
\[
 \frac{696221180224885514072002509}{48415339169645815039062500000}>0.
\]
Consequently
\[
 eH_{\rm head}<u^2C_H,\qquad
 C_H=\frac{17160482889458863049}{148731921929151943800}.
\]
The same bound on $e$ and~\eqref{eq:low-Q3-square-cost} give
$e^2J_{\rm minus}<d^2u^4J(q)\le u^4J(q)$.
Together with $F_{\rm pipe}>69/500$, equation~\eqref{eq:low-large-master}
therefore proves
\[
 F_{\rm floor}(v)>L(q):=69/500-u^2C_H-u^4J(q).
\]
For the last continuous comparison, ordinary differentiation yields
\[
 \frac{(q+1)^3}{q}L'(q)
 =-2C_H+\frac{q(q+1)^3}{27(q-4)^4}
              +\frac{q^2-2q+2}{18q(q-4)^2}.
\]
Both nonconstant terms decrease strictly for $q\ge10$.
For the first, its logarithmic derivative is
$1/q+3/(q+1)-4/(q-4)<0$; the derivative of the second is
$(-q^3-6q+8)/[18q^2(q-4)^3]<0$.
Thus $L'$ can change sign only from positive to negative, and
$L$ has no interior minimum on $[10,21]$.  Its two endpoint
margins are
\[\begin{gathered}
 L(10)-1/40=\frac{195098675960795601986623}{65597470507242819053271000}>0,
 \\
 L(21)-1/40=\frac{713805528405149594753}{359931251068547703996000}>0.
\end{gathered}\]
This proves $F_{\rm floor}(v)>1/40$ for $10\le q\le21$ and $h\le1/4$.

\begin{proposition}[Negativity of $v'(\lambda)$ for $1<\lambda<\min(m,2/\sqrt3)$]\label{prop:entire-low-region}
For every fixed offspring law satisfying the assumptions of Theorem~\ref{thm:computer},
$v'(\lambda)<0$ for $1<\lambda<\min(m,2/\sqrt3)$.
\end{proposition}
\begin{proof}
Lemma~\ref{lem:calibration} covers $\alpha=1$. When $0<\alpha<1$,
we have $q>0$ and $0<h<1/2$. The six ranges
\[\begin{gathered}
 q\le1/15;\qquad1/15<q\le3;\qquad3<q\le10;\qquad q\ge21;\\
 10<q<21,\ h\ge1/4;\qquad10<q<21,\ 0<h<1/4
\end{gathered}\]
exhaust these parameters.  The preceding arguments give
$F_{\rm floor}(v)>0$ in the last five, and directly $D_\lambda>0$
in the first.  Lemma~\ref{lem:mellin-floor} gives
$D_\lambda\ge F_*>F_{\rm floor}(v)$ in all remaining cases.
Since $Z>0$, the derivative formula~\eqref{eq:full-sign} gives
$v'(\lambda)<0$. The finite arithmetic inputs are
exactly Lemmas~\ref{lem:low-Q0-arithmetic},
\ref{lem:low-Q12-arithmetic}, \ref{lem:low-Q3-arithmetic}, and
\ref{lem:low-Q3c-arithmetic}; their expressions, domains, positive
denominators, and uses in the proof have all been specified above.
\end{proof}

\subsection{Moment estimates for \texorpdfstring{$2/\sqrt3\le\lambda\le707/500$}{2/sqrt(3) <= lambda <= 707/500}}\label{sec:near-region}
We now prove $D_\lambda>0$ for
\begin{equation}\label{eq:near-domain}
 \frac2{\sqrt3}\le\lambda\le\frac{707}{500},\qquad \lambda<m.
\end{equation}
Lemma~\ref{lem:calibration} covers $\alpha=1$. We may therefore assume $0<\alpha<1$, so that $0<u<\lambda^{-2}$. We keep the moments $a_s=\E U^s$ of the given escape law fixed when varying $w$ in $\widetilde F(w)$. We write $a=a(w)$ and retain $h=cu/(1-u)$.

\subsubsection{A comparison that retains the first two moments}
Set
\begin{equation}\label{eq:near-profile-defs}
 \begin{gathered}
 Z_s=(1-u)\sum_{j=0}^{s-1}\lambda^j+u,\qquad r_s=\frac{Z_s}{Z_{s+1}},\\
 A=Z_2=1+\lambda(1-u),\qquad B=Z_3=1+\lambda(\lambda+1)(1-u).
 \end{gathered}
\end{equation}
Thus $\beta_s=u/Z_s$, $r_s=\beta_{s+1}/\beta_s$, and $A$ without arguments is the scalar $Z_2$, distinct from the pair kernel $A(x,y)$. Write
\[
 z=\frac{a_2}{u},\qquad \rho=\frac{a_2}{\beta_2}=Az.
\]
Here $\rho$ is the value at $s=2$ of the moment ratio $\rho(s)=a_s/\beta_s$. The recursion on $\{\nu=1\}$ rules out a positive constant value of $U$. Hence $a_2>u^2$, and Lemma~\ref{lem:canonical-moments} gives
\begin{equation}\label{eq:near-moment-domain}
 u<z<1/A,\qquad uA<\rho<1.
\end{equation}

To bound the contribution of the terms $\lambda^{n-1}\beta_n^2$ to the speed quotient, put
\begin{equation}\label{eq:near-cap}
 \begin{gathered}
 \ell_L=2\left(\frac c{\lambda+1}+\frac{c^3}{3(\lambda+1)^3}
                   +\frac{c^5}{5(\lambda+1)^5}\right),\qquad
 \ell_U=\ell_L+\frac{2c^7}{7(\lambda+1)^7[1-c^2/(\lambda+1)^2]},\\
 K_{\rm cap}=\frac c{\lambda\ell_L}+\frac{1-u}{2}
 +\frac{\ell_U}{12c}(\lambda+1-2\lambda u)(1-u).
 \end{gathered}
\end{equation}
All denominators in this formula are positive.

\begin{lemma}[Second-moment comparison]\label{lem:near-prefix}
For the moments of $U$ satisfying~\eqref{eq:near-moment-domain}, define
\begin{equation}\label{eq:near-assigned-pair}
 \begin{gathered}
 T_2=u^2\left[1+\frac{\rho^2(K_{\rm cap}-1+u)}{1-u}\right],\qquad
 Q_2=u+\lambda u^2z+\frac{Au^2z^2}{1-u},\\
 w=\frac{cT_2}{2Q_2-(\lambda+1)T_2},\qquad
 a=\frac{u-(w/c)A_0}{a_2}.
 \end{gathered}
\end{equation}
Then all displayed denominators are positive and
\begin{equation}\label{eq:near-assigned-range}
 v<w<h,\qquad \lambda^2u<a<\alpha<1.
\end{equation}
Moreover, $1<K_{\rm cap}<2/(1+\lambda u)$.
\end{lemma}
\begin{proof}
Let
\[
 \kappa=\frac{c^2}{1-u}\sum_{n\ge1}
               \frac{\lambda^{n-1}}{(\lambda^n-d)^2}.
\]
We first justify the bound $\kappa<K_{\rm cap}$. For $s>0$,
\[
 \frac1{1-e^{-s}}<\frac1s+\frac12+\frac{s}{12}.
\]
Indeed, after multiplication by the positive denominator, the difference is
\[
 (e^s-1)(1-s/2+s^2/12)-s
 =\sum_{n\ge5}\frac{(n-3)(n-4)}{12n!}s^n>0.
\]
The function $f(x)=\lambda^{x-1}/(\lambda^x-d)^2$ is the positive exponential mixture
\[
 f(x)=\lambda^{-1}\sum_{j\ge0}(j+1)d^j
                    e^{-(j+1)x\log\lambda}.
\]
Tonelli's theorem and the preceding one-variable inequality give
\[
 \sum_{n\ge1}f(n)<\int_1^\infty f(x)\dd x+\tfrac12 f(1)-\tfrac1{12}f'(1).
\]
Evaluation of these three terms, using $c+h=c/(1-u)$, yields
\[
 \kappa<\frac c{\lambda\log\lambda}+\frac{1-u}{2}
 +\frac{\log\lambda}{12c}(\lambda+1-2\lambda u)(1-u).
\]
The positive series for $2\operatorname{arctanh}(c/(\lambda+1))$ shows
\[
 \ell_L<\log\lambda<\ell_U<\frac{c(\lambda+1)}{2\lambda}.
\]
Since $\lambda+1-2\lambda u>0$, this proves $\kappa<K_{\rm cap}$.

The same inequalities give
\[
 K_{\rm cap}<k_0:=
 \frac{\lambda^2+13\lambda+6-\lambda(3\lambda+7)u+2\lambda^2u^2}{12\lambda}.
\]
For $23/20\le\lambda\le3/2$ and $0\le u\le\lambda^{-2}$, $(1+\lambda u)k_0<2$. Here is a finite verification of its sign which also records the relevant domain. After multiplying by $12\lambda^4$ and using the degree-three Bernstein basis in $\lambda^2u\in[0,1]$, the four coefficients are
\[
\begin{split}
 &-\lambda^3(\lambda^2-11\lambda+6),\\
 &-\tfrac13\lambda^2(3\lambda^3-32\lambda^2+28\lambda-1),\\
 &-\tfrac13\lambda c(3\lambda^3-28\lambda^2+10\lambda+5),\\
 &-c^2(\lambda^3-8\lambda^2-\lambda+2).
\end{split}
\]
Each is positive. The three cubics have negative derivatives on this interval and are negative at $23/20$; the quadratic is negative there as well. Also
\[
 K_{\rm cap}>\frac2{\lambda+1}+\frac{1-\lambda^{-2}}2>1
\]
on the same interval. This proves the asserted range of $K_{\rm cap}$.

Divide each moment $a_s$ by the corresponding term of $1,u,\rho\beta_2,\rho\beta_3,\ldots$. The ratios equal one through index two and decrease thereafter. The successive backward ratios of the comparison sequence are $1/u$, $1/z$, and then $1/r_s$, in decreasing order. The covariance argument used in Lemma~\ref{lem:harmonic-speed}, first for finite sums and then by absolute convergence, gives $Q/T>Q_2(\kappa)/T_2(\kappa)$. The formulas follow from
\[
 \sum_{n\ge1}\lambda^{n-1}\beta_n^2=\frac{u^2\kappa}{1-u},
 \qquad \sum_{n\ge1}\lambda^{n-1}\beta_n\beta_{n-1}=\frac u{1-u}.
\]
The second identity telescopes, since
$(\lambda^{n-1}-d)^{-1}-(\lambda^n-d)^{-1}
=c\lambda^{n-1}/[(\lambda^{n-1}-d)(\lambda^n-d)]$.

For completeness, the denominators needed for this comparison can be checked before taking any quotient. Set
\[
 R=1+\frac{A^2(K_{\rm cap}-1+u)z^2}{1-u},\qquad
 Q=2+2\lambda uz+\frac{2Auz^2}{1-u}.
\]
Both $Q-R$ and $Q-(\lambda+1)uR$ are positive for $0\le z\le1/A$. If their quadratic coefficient is nonnegative, every coefficient is nonnegative and the constant is positive. Otherwise they are concave, and their endpoint values are positive: at $z=1/A$ they are respectively $(2-K_{\rm cap})/(1-u)$ and $[2-(\lambda+1)uK_{\rm cap}]/(1-u)$. The range of $K_{\rm cap}$ proves these signs. Increasing $\kappa$ to $K_{\rm cap}$ increases the speed quotient, so $v<w$.

The comparison with the full reference sequence is equally explicit. With $T_p(K)=u^2K/(1-u)$ and $Q_p=u/(1-u)$,
\[
 T_p(K)Q_2-T_2(K)Q_p
 =\frac{(1-\rho^2)u^3(K-1)}{1-u}
   +\lambda(\rho-\rho^2)u\beta_2T_p(K)\ge0.
\]
Consequently
\[
 w\le\frac{cuK_{\rm cap}}{2-(\lambda+1)uK_{\rm cap}}<\frac{cu}{1-u}=h.
\]
The function $a(w)$ strictly decreases with $w$. The calculation for~\eqref{eq:hard-calibration} gives $a(h)>\lambda^2u$ and $a(v)=\alpha$, so $v<w<h$ yields $\lambda^2u<a<\alpha$.
\end{proof}

\subsubsection{An inequality for the first three moments of \texorpdfstring{$U$}{U}}
The next step uses the recursion itself. Moment inequalities for an arbitrary law on $(0,1/\lambda)$ would not suffice.

\begin{lemma}[Three-moment inequality]\label{lem:near-three-moments}
The moments of $U$ under the escape law satisfy
\begin{equation}\label{eq:near-three-moment-bound}
 \frac{a_3}{a_2}<\frac1{\lambda+1-\lambda a_2/u}.
\end{equation}
\end{lemma}
\begin{proof}
Put $\eta=(a_2-a_3)/(u-a_2)$ and $p(y)=(y-\eta)(1-y)$. Then $0<\eta\le z$ and $\E[Up(U)]=0$; the upper bound follows from $ua_3\ge a_2^2$. For a child list $x_i=U_i$, set
\[
 S=\sum_i\frac{x_i}{1-x_i},\quad F=\frac S{1+S},\quad
 A_i=\left(1+(1-x_i)\sum_{j\ne i}\frac{x_j}{1-x_j}\right)^{-1}.
\]
The recursion says $F\overset{\mathrm d}=\lambda U$. Directly,
\[
 F=1-A_i(1-x_i),\qquad \sum_i x_iA_i=F,
 \qquad p(F)=A_ip(x_i)+A_i(1-A_i)(1-x_i)^2.
\]
Replace $x_i$ by $\eta$ only in the definition of $A_i$, obtaining $A_{i,\eta}$. Independence of the $i$th child from the other children gives
\[
\begin{split}
 \E[Fp(F)]={}&\E\sum_i x_i(A_i^2-A_{i,\eta}^2)p(x_i)\\
 &+\E\sum_i x_iA_i^2(1-A_i)(1-x_i)^2>0.
\end{split}
\]
The first summands are nonnegative because $A_i$ increases with $x_i$. The second sum is positive when there are at least two children. Every absolute integrand is bounded by a constant times $\nu$, so finite mean justifies conditioning and summation. Substitution of $F\overset{\mathrm d}=\lambda U$ gives
\[
 \E[Fp(F)]
 =\frac{\lambda cu}{1-z}
 \left[z-(\lambda+1-\lambda z)\frac{a_3}{u}\right].
\]
Its strict positivity is precisely~\eqref{eq:near-three-moment-bound}.
\end{proof}

To deduce positivity of the moment differences $\Delta_s^{\mathrm{ass}}$ below from~\eqref{eq:near-three-moment-bound}, we need the following bounds on $K_{\rm cap}$.

\begin{lemma}[Upper bounds on $K_{\rm cap}$]\label{lem:near-cap-arithmetic}
Let $2/\sqrt3\le\lambda\le707/500$, $0<u\le\lambda^{-2}$, and define
\[
\begin{split}
 \overline K={}&\frac{\lambda+1}{2\lambda}-\frac{c^2}{6\lambda(\lambda+1)}
 +\frac{1-u}{2}
 +\frac{(\lambda^2+10\lambda+1)(\lambda+1-2\lambda u)(1-u)}{72\lambda(\lambda+1)},\\
 P={}&\lambda^2(\lambda+1)^2(1-u^2)+2\lambda+1,\\
 R_B={}&AP-(\lambda+1)^2(1-u)(1-\lambda u)^2,\qquad
 K_B=\frac{2\lambda R_B}{\lambda R_B+uBP}.
\end{split}
\]
Then $R_B>0$, and $K_{\rm cap}<\overline K<K_B<2\lambda A/(\lambda A+uB)$.
The finite coefficient input for the middle inequality is the following. Put
\[
 k_n=(2\lambda^3+20\lambda^2+2\lambda)u^2
 -(3\lambda^3+67\lambda^2+49\lambda+1)u
 +\lambda^3+71\lambda^2+143\lambda+25
\]
and
\[
 F_B=\frac{144\lambda^2(\lambda+1)R_B-k_n(\lambda R_B+uBP)}{\lambda+1}.
\]
This is a polynomial of degree six in $u$. On $1\le\lambda\le3/2$, its $\lambda$ derivative has positive degree-six Bernstein coefficients on $0\le u\le3/4$. If $f(u)=F_B(2/\sqrt3,u)=\sum_{i=0}^6a_i u^i$ and $J(u)=24a_4+60a_5u+120a_6u^2$, the exact finite comparisons are
\[
\begin{gathered}
 a_3<0,\quad a_6<0,\quad J'(3/4)>0,\quad J(3/4)<0,\\
 f(1/2)>2,\quad -18<f'(1/2)<0,\quad f''(3/5)>1000,\quad
 f'(3/5)>0,\quad f'(3/4)>0.
\end{gathered}
\]
The supplement verifies the signs of these Bernstein coefficients and evaluations. Their denominators are products of positive integers and powers of $\sqrt3$.
\end{lemma}
\begin{proof}
Since $a_6<0$, $J'$ decreases to its positive value at $3/4$. Thus $J$ increases and is negative throughout $[0,3/4]$. It follows that $f'''=6a_3+uJ<0$. In particular, $f''>1000$ on $[0,3/5]$, whence
\[
 f(u)>2-18|u-1/2|+500(u-1/2)^2\ge919/500.
\]
On $[3/5,3/4]$, $f'$ is concave with positive endpoint values, hence $f'>0$. Thus $f>919/500$ there too. Positivity of the Bernstein coefficients gives $F_B(\lambda,u)\ge f(u)>0$. The denominator in
\[
 K_B-\overline K=\frac{F_B}{72\lambda(\lambda R_B+uBP)}
\]
is positive: indeed
\[
 R_B=(\lambda+1)^2(1-u)[A\lambda^2(1+u)-(1-\lambda u)^2]+A(2\lambda+1)>0.
\]
The comparisons $K_{\rm cap}<\overline K$ follow from the two finite identities
\[
 (1-g^2/3)(1+g^2/3+g^4/5)-1=\frac{g^4(4-3g^2)}{45}>0,
\]
\[
 1+\frac{g^2}{3(1-g^2)}-\left(1+\frac{g^2}3+\frac{g^4}5+
                  \frac{g^6}{7(1-g^2)}\right)
 =\frac{g^4(2/15+2g^2/35)}{1-g^2}>0,
 \qquad g=\frac c{\lambda+1}.
\]
Finally,
\[
 \frac{\lambda A}{B}-\frac{\lambda R_B}{BP}
 =\frac{\lambda(\lambda+1)^2(1-u)(1-\lambda u)^2}{BP}>0.
\]
Since $2x/(x+u)$ increases in $x>0$, this gives the last comparison.
\end{proof}

\begin{lemma}[Positivity of the moment differences]\label{lem:near-condition-f}
For $(w,a)$ in~\eqref{eq:near-assigned-pair}, let
\[\begin{gathered}
 Q_a(U)=-\frac wc+\left[1+\frac{(\lambda+1)w}{c}\right]U
                  -\left(a+\frac{\lambda w}{c}\right)U^2,\\
 \Delta_s^{\mathrm{ass}}=\E[U^{s-1}Q_a(U)]-(1-a)a_s.
\end{gathered}\]
Then $\Delta_s^{\mathrm{ass}}>0$ for every real $s\ge2$.
\end{lemma}
\begin{proof}
Retain $R,Q$ from the denominator check in Lemma~\ref{lem:near-prefix}, and put $D_f=Q-R$, $D_w=Q-(\lambda+1)uR$. The only remaining scalar inequality is
\begin{equation}\label{eq:near-W}
 W:=\lambda D_f-uR(\lambda+1-\lambda z)>0.
\end{equation}
Since $a<1$ and $w<h$, the formula for $a$ gives
\[
 z>\frac{\lambda u}{1+cu}.
\]
If $u\le5/8$, solving $W>0$ for $K_{\rm cap}$ gives $K_{\rm cap}<K_{\rm cal}(z)$, with
\[
 K_{\rm cal}(z)=1-u+
 \frac{(1-u)[\lambda-(\lambda+1)u+\lambda u(2\lambda+1)z+
                       2\lambda Au z^2/(1-u)]}
      {A^2z^2[\lambda+(\lambda+1)u-\lambda uz]}.
\]
Its derivative has the sign of a cubic with coefficients
\[
\begin{split}
 &2[(\lambda+1)^2u^2-\lambda^2],\\
 &-2\lambda u[\lambda c+(\lambda+1)(\lambda+2)u],\\
 &2\lambda^2u^2(2\lambda+1),\qquad 2\lambda^2Au^2/(1-u).
\end{split}
\]
This cubic is negative for $z\le1/A$. To see the sign directly, use $1-u\ge3/8$ to obtain
\[
 (\lambda+1)(\lambda+2)-
 \frac{\lambda[(2\lambda+1)+(1-u)^{-1}]}A
 \ge\frac{9\lambda^3+3\lambda^2+2\lambda+48}{24+9\lambda}>\frac{15}{8}.
\]
The last difference has numerator $72\lambda^3+24\lambda^2-119\lambda+24$, positive at $\lambda=1$ and increasing for $\lambda\ge1$. The cubic is consequently less than
\[
 -2\left\{\lambda^2+u^2[\lambda^3-2\lambda^2-2\lambda-1+(15/8)\lambda u]\right\}<0.
\]
The expression in braces increases with $\lambda\ge9/8$ and decreases with $0\le u\le5/8$; its value at $(9/8,5/8)$ is $643/8192$. Thus $K_{\rm cal}$ decreases, and its value at $1/A$ is $2\lambda A/(\lambda A+uB)$. Lemma~\ref{lem:near-cap-arithmetic} proves~\eqref{eq:near-W}.

If $u\ge5/8$, the identity
\[
 2R-K_{\rm cap}Q
 =2(1-Az)[(1-K_{\rm cap})+(A-(\lambda+1)K_{\rm cap})z]<0
\]
gives $R/Q<K_{\rm cap}/2$. Also
\[
 \lambda+(\lambda+1)u-\lambda uz
 <\frac{\lambda+(\lambda^2+1)u-u^2}{1+cu}.
\]
It therefore suffices that
\[
 \overline K<\frac{2\lambda(1+cu)}{\lambda+(\lambda^2+1)u-u^2}.
\]
After multiplication by the positive denominator $72\lambda(\lambda+1)$, the difference is a quartic in $u-5/8\in[0,1/8]$. Its constant and linear coefficients are
\[
\begin{split}
 &\frac3{2048}(40\lambda^5+45744\lambda^4+24823\lambda^3
                     -56694\lambda^2-34745\lambda-3900),\\
 &\frac1{64}(26\lambda^5+8564\lambda^4-2688\lambda^3
                     -6539\lambda^2+2572\lambda+405),
\end{split}
\]
and are positive. For the constant coefficient the polynomial increases with $\lambda$ and its left-end value is $(16476-28378/\sqrt3)/9>0$, since $3\cdot16476^2-28378^2=9064844$. The linear coefficient is positive by grouping its negative terms with $8564\lambda^4$. The cubic coefficient is negative, the quartic coefficient is $2\lambda(\lambda^2+10\lambda+1)>0$, and the quadratic coefficient plus one eighth of the cubic coefficient is
\[
 \frac{\lambda^4(400-16\lambda)+389\lambda^3+258\lambda^2+1525\lambda+384}{16}>0.
\]
This proves the quartic's sign, and hence~\eqref{eq:near-W}, also in this case.

Now let $\theta=a_3/a_2$. The formula for $a$ and~\eqref{eq:near-three-moment-bound} give
\[
 \frac{\Delta_2^{\mathrm{ass}}}{a_2}
 =a(1-\theta)+\frac wc(\lambda+1-\lambda\theta-1/z)
 >\frac{(1-z)W}{(\lambda+1-\lambda z)zD_w}>0.
\]
The direction is correct because the coefficient of $\theta$ is
$-(a+\lambda w/c)=-D_f/(zD_w)<0$.
Finally $Q_a(y)-(1-a)y$ is a concave quadratic, negative at zero and positive at $1/\lambda$, where its value is $ac/\lambda^2$. It changes sign from negative to positive at its unique zero $\xi\in(0,1/\lambda)$. For every $s\ge2$,
\[
 \Delta_s^{\mathrm{ass}}-\xi^{s-2}\Delta_2^{\mathrm{ass}}
 =\E\{U[U^{s-2}-\xi^{s-2}][Q_a(U)-(1-a)U]\}\ge0.
\]
Thus the bound holds for every real exponent $s\ge2$.
\end{proof}

\subsubsection{A lower bound before truncating the series}
We bound each summand with index $s\ge2$ from below by a positive quantity. This allows us to keep only the first four terms in the subsequent estimate.

\begin{lemma}[A series lower bound for $D_\lambda$]\label{lem:near-absorption}
For the pair $(w,a)$ in~\eqref{eq:near-assigned-pair} and every integer $s\ge2$, put
\[
 e_s=1-\lambda r_s,\qquad L_s=\frac1\lambda-r_{s-1},\qquad
 B_s=\frac{1+w}{\lambda}e_s^2+2(1-a)(r_s-r_{s-1}),\qquad s\ge2.
\]
Then $B_s,L_s>0$, and
\begin{equation}\label{eq:near-root-floor}
 D_\lambda>
 \frac{(1+w)(u-\lambda a_2)^2-2(1-a)u(u-\lambda a_2)}\lambda
 +\sum_{s\ge2}s\lambda^{s-1}
       [a_s^2B_s+2a_sL_s\Delta_s^{\mathrm{ass}}].
\end{equation}
The series converges absolutely.
\end{lemma}
\begin{proof}
Use the integral representation of the moment ratios in Lemma~\ref{lem:canonical-marker} and the series expansion in Lemma~\ref{lem:mellin-floor}. With
\[\begin{gathered}
 \delta_s=a_sr_s-a_{s+1},\qquad \varepsilon_s=a_s-\lambda a_{s+1},\\
 \Theta_s=\frac c{\lambda\log\lambda}\beta_s[-\rho'(s)],\quad
 C_s=\frac c{\lambda\log\lambda}\frac{1-1/(2s)}{r_s},
\end{gathered}\]
they give $\delta_s>0$, $\Theta_s>C_s\delta_s$, and
\[
 \gamma_{s-1}=a_s(r_s-r_{s-1})-\delta_s+\Theta_s,
 \qquad \gamma_{s-1}+\varepsilon_s/\lambda=a_sL_s+\Theta_s>a_sL_s.
\]
We claim that
\begin{equation}\label{eq:near-age-sign}
 e_s+(1-\lambda^2u)(C_s-1)>0\qquad(s\ge2).
\end{equation}
To prove it, set $N_s=cZ_s=\lambda^s(1-u)-1+\lambda u$. The recursion
$N_{s+1}=\lambda N_s+c(1-\lambda u)$ and $N_1=c$ give
\[
 e_s\ge\frac{1-\lambda^2u}{\sum_{j=0}^s\lambda^j}.
\]
Indeed, $1-\lambda u\ge1-\lambda^2u$ and $N_{s+1}\le\lambda^{s+1}-1=c\sum_{j=0}^s\lambda^j$. Convexity of $1/x$ on $[1,\lambda]$ gives $c/(\lambda\log\lambda)>2/(\lambda+1)$, while $r_s<1/\lambda$. Finally
\[
 s(\lambda+1)-[s-(s-1)\lambda]\sum_{j=0}^s\lambda^j
 =\lambda\left[(s-1)+c\sum_{j=1}^{s-1}j\lambda^j\right]>0.
\]
These three inequalities imply~\eqref{eq:near-age-sign}. If $C_s\ge1$, then $e_s>0$ gives the next inequality directly. If $C_s<1$, use $a>\lambda^2u$ in the same bound. In both cases
\[
 (1+w)e_s+(1-a)(C_s-1)>0.
\]
Writing $\xi_s=\Theta_s-C_s\delta_s>0$, expansion now gives
\[
\begin{split}
 \frac{1+w}{\lambda}\varepsilon_s^2+2(1-a)a_s\gamma_{s-1}
 ={}&a_s^2B_s+2a_s\delta_s[(1+w)e_s+(1-a)(C_s-1)]\\
 &+(1+w)\lambda\delta_s^2+2(1-a)a_s\xi_s>a_s^2B_s.
\end{split}
\]
Here $r_s-r_{s-1}>0$ follows from $Z_s^2-Z_{s-1}Z_{s+1}=(1-u)\lambda^{s-1}(1-\lambda u)>0$.

To identify the term with $s=1$ and bound those with $s\ge2$, write the coefficient $J_{s-1}$ as
\[
 J_{s-1}=2ct_{s-1}+c[(1+w)\varepsilon_s-2(1-a)a_s],
 \qquad t_{s-1}=(1-a)a_s+\Delta_s^{\mathrm{ass}}.
\]
For $s\ge2$ the term in the series for $\widetilde F(w)$, before multiplication by $s\lambda^{s-1}$, is therefore
\[
 \frac{1+w}{\lambda}\varepsilon_s^2+2(1-a)a_s\gamma_{s-1}
       +2(\gamma_{s-1}+\varepsilon_s/\lambda)\Delta_s^{\mathrm{ass}}.
\]
At $s=1$, $t_0=0$, and the summand is the first term of~\eqref{eq:near-root-floor}. For $s\ge2$, Lemma~\ref{lem:near-condition-f} gives $\Delta_s^{\mathrm{ass}}>0$, so we may bound its coefficient from below by $2a_sL_s$. Lemma~\ref{lem:assigned-floor} gives
$D_\lambda\ge F_*=\widetilde F(v)>\widetilde F(w)$, proving the strict inequality.

All these rearrangements can first be made in finite sums. At fixed $\lambda>1$, $a_s\le\lambda^{-s}$ and $\beta_s\le C_\lambda\lambda^{-s}$. Moreover $\E[U^s|\log U|]\le C_\lambda\lambda^{-s}$ for all sufficiently large $s$, by maximizing $x^s|\log x|$ on $[0,1/\lambda]$. The same bounds hold for $\Theta_s,\gamma_{s-1},\Delta_s^{\mathrm{ass}}$ up to a fixed factor, and $t_{s-1}=O_\lambda(\lambda^{-s})$. Every term above is thus bounded by $C_\lambda s\lambda^{-s}$. Dominated summation proves the asserted convergence and all passages to the infinite sum.
\end{proof}

\begin{lemma}[A lower bound using the first four terms]\label{lem:near-four-age}
Let $T=1-\sqrt{1-\rho}$ and $f_s=T^{s-1}[s-(s-1)T]$. Define
\[
 M_s=a(1-r_s)+\frac wc\left[\lambda+1-\lambda r_s-
                              \frac1{r_{s-1}f_s}\right],\qquad s\ge2,
\]
and
\begin{equation}\label{eq:near-S4}
 \frac{S_4}{u^2}=
 \frac{(1+w)(1-\lambda\rho/A)^2-2(1-a)(1-\lambda\rho/A)}\lambda
 +\sum_{s=2}^4\frac{s\lambda^{s-1}f_s^2}{Z_s^2}[B_s+2L_sM_s].
\end{equation}
Then $D_\lambda>S_4$.
\end{lemma}
\begin{proof}
The integral representation of $\rho_s$ with respect to $\sigma$ gives
\[
 \rho_s=\E_\sigma[t^{s-1}(s-(s-1)t)],\qquad
 \E_\sigma(1-t)^2=1-\rho.
\]
As a function of $x=(1-t)^2$, the first integrand is
$(1-\sqrt x)^{s-1}[1+(s-1)\sqrt x]$. Its derivative is
$-s(s-1)(1-\sqrt x)^{s-2}/2$, an increasing function for integer $s\ge2$.
Jensen's inequality therefore gives $a_s\ge\beta_sf_s$, with equality at $s=1,2$. Also $a_{s-1}\le\beta_{s-1}$ and $a_{s+1}/a_s<r_s$, so direct substitution in $\Delta_s^{\mathrm{ass}}$ gives $\Delta_s^{\mathrm{ass}}/a_s>M_s$. Lemma~\ref{lem:near-condition-f} gives $\Delta_s^{\mathrm{ass}}>0$. For $s\ge2$ the expression in square brackets in~\eqref{eq:near-root-floor} is therefore bounded below by
\[
 (\beta_sf_s)^2[B_s+2L_s\max\{0,M_s\}].
\]
Since $B_s,L_s>0$, the bracket is positive when we use $a_s\ge\beta_sf_s$. We can now drop the terms with $s\ge5$ and replace $\max\{0,M_s\}$ by $M_s$ for $s=2,3,4$. This gives~\eqref{eq:near-S4}.
\end{proof}

\subsubsection{Positivity of \texorpdfstring{$S_4$ for $0<u\le1/2$}{S4 for 0 < u <= 1/2}}
We prove $S_4>0$ by estimating~\eqref{eq:near-S4} on a larger set of scalar parameters. The comparison $D_\lambda>S_4$ from Lemma~\ref{lem:near-four-age} uses the moments of $U$ under the given escape law.

Temporarily replace $K_{\rm cap}$ in~\eqref{eq:near-assigned-pair} by a scalar $K$. In the variables $\lambda,u,\rho$, put
\begin{equation}\label{eq:near-finite-calibration}
 \begin{gathered}
 R=1-u+\rho^2(K-1+u),\qquad
 Q=A(1-u)+\lambda u\rho(1-u)+u\rho^2,\\
 D_K=2Q-(\lambda+1)uAR,\qquad
 A_0^\mathrm{num}=[1-(\lambda+1)u]A+\lambda u\rho,\\
 w=\frac{cuAR}{D_K},\qquad
 a=\frac{A(D_K-RA_0^\mathrm{num})}{\rho D_K},\qquad
 \Phi(K)=\rho D_K S_4(K)/u^2.
 \end{gathered}
\end{equation}
Define the explicit quadratic bound
\begin{equation}\label{eq:near-Kb}
 K_b=\frac{\lambda+1}{2\lambda}+\frac{1-u}{2}
       +\frac{(\lambda+1)(\lambda+1-2\lambda u)(1-u)}{24\lambda}.
\end{equation}
The inequalities for $\ell_L,\ell_U$ give $K_{\rm cap}<K_b$.

\begin{lemma}[Comparison of $\Phi(K_{\rm cap})$ and $\Phi(K_b)$]\label{lem:near-K-affinity}
For $23/20\le\lambda\le3/2$, $0<u\le1/2$, $0<\rho\le1$, the denominators $D_{K_b}$ and $D_{K_{\rm cap}}$ are positive, and
\[
 \Phi(K_{\rm cap})>\Phi(K_b).
\]
\end{lemma}
\begin{proof}
First $(1+\lambda u)K_b<2$. Indeed this product increases with $u$ on $[0,1/2]$, as follows by differentiating its displayed cubic; its derivative is at least $(-2\lambda^3+9\lambda^2-2\lambda-1)/(24\lambda)>0$. At $u=1/2$ the remaining difference is
\[
 2-(1+\lambda/2)K_b(\lambda,1/2)
 =\frac{6+c(19-37c)}{96\lambda}>0.
\]
Direct subtraction now gives
\begin{equation}\label{eq:near-positive-D}
 \frac{D_{K_b}}{A(1-u)}-R
 =(1-\rho)\left[(1-\lambda u)(1+\rho)+\frac{2\lambda u\rho}{A}\right]
 +\frac{2-(1+\lambda u)K_b}{1-u}\rho^2>0.
\end{equation}
It also proves positivity at the smaller $K_{\rm cap}$.

Here is the full dependence on $K$. Put $E=1-\lambda\rho/A$ and $W_s=s\lambda^{s-1}/Z_s^2$. The cancellation
$r_s-r_{s-1}=L_s(1-r_s)$ in~\eqref{eq:near-S4} gives
\[
 S_4/u^2=F_0+aF_a+wF_w,
\]
where
\[
\begin{split}
 F_0&=(E^2-2E)/\lambda+
       \sum_{s=2}^4W_sf_s^2[e_s^2/\lambda+2(r_s-r_{s-1})],\\
 F_a&=2E/\lambda,\\
 F_w&=E^2/\lambda+
 \sum_{s=2}^4W_s\left\{f_s^2\left[\frac{e_s^2}{\lambda}
       +\frac{2L_s}{c}(\lambda+1-\lambda r_s)\right]
       -\frac{2L_sf_s}{cr_{s-1}}\right\}.
\end{split}
\]
Set $\widehat C=\rho F_0+AF_a$. Substitution of~\eqref{eq:near-finite-calibration} gives
\[
 \Phi(K)=2Q\widehat C-AR\widehat H,
 \qquad \widehat H=(\lambda+1)u\widehat C+A_0^\mathrm{num}F_a-cu\rho F_w.
\]
Its sign is exposed by the exact cancellation
\[
 \widehat H=\frac{2E}{\lambda}\left(A-\frac{\lambda u\rho^2}{A}\right)
 +2u\rho\sum_{s=2}^4\frac{W_sL_sf_s}{r_{s-1}}
          [1-\lambda r_{s-1}r_s^2f_s]>0.
\]
Here $0<f_s\le1$, $\lambda r_s<1$, and $E>0$. Thus
$\partial_K\Phi=-A\rho^2\widehat H<0$, proving the claim. Only the scalar argument $K$ of $\Phi$ varies in this comparison.
\end{proof}

Henceforth $K=K_b$ and $\eta=K_b-1+u$. Define
\begin{equation}\label{eq:near-complete-finite-profile}
\begin{gathered}
 \Phi_R(\rho)=\frac2\lambda\left(1-\frac{\lambda\rho}{A}\right)
 \left\{Q\left(2A-\rho-\frac{\lambda\rho^2}{A}\right)
       -AR\left(A-\frac{\lambda u\rho^2}{A}\right)\right\},\\
 \omega_s=\frac{s\lambda^{s-2}(1-\lambda u)}{Z_s^3},\qquad
 V_s=\frac{\omega_suAR}{r_{s-1}},\\
 U_s=\omega_s[Q(2-r_s-\lambda r_s^2)+uAR\lambda r_s^2],\qquad s=2,3,4.
\end{gathered}
\end{equation}
All $U_s,V_s$ are positive. Substituting the formulas for $a,w$ in~\eqref{eq:near-S4} and multiplying both sides by $\rho D_{K_b}$ gives
\begin{equation}\label{eq:near-complete-Phi}
 \Phi(K_b)=\Phi_R+2\rho\sum_{s=2}^4(U_sf_s^2-V_sf_s).
\end{equation}
Here $\Phi_R$ comes from the term with $s=1$. Since
$\rho^2-f_3=T^2(1-T)^2$ and $f_3^2-\rho f_4=T^4(1-T)^2$, it also gives
\begin{equation}\label{eq:near-merge}
\begin{split}
 \Phi(K_b)=\Phi_R+2\rho\{&(U_2-V_3)\rho^2-V_2\rho
 +(\rho U_3-V_4)f_4+U_4f_4^2\\
 &+T^2(1-T)^2[V_3+U_3T^2]\}.
\end{split}
\end{equation}

\begin{lemma}[Bounds on $U_2-V_3$ and $\rho U_3-V_4$]\label{lem:near-age-comparisons}
In the present range,
\[
 U_2-V_3>\frac{(1-\lambda u)Q}{A^3}>0.
\]
For $1/2\le\rho\le1$, one also has $\rho U_3-V_4>0$.
\end{lemma}
\begin{proof}
For the first claim, let $x=uAR/Q$. Equation~\eqref{eq:near-positive-D} implies
$x<2u/(1+\lambda u)\le40/63$. The identities in~\eqref{eq:near-profile-defs} give
\[
 U_2-V_3=\frac{2(1-\lambda u)Q}{A^3}
      [2-r_2-\lambda r_2^2(1+x/2)].
\]
The derivative $\partial_u r_2=\lambda^2/Z_3^2$ is positive. At $u=1/2$, direct multiplication gives $r_2<5/7$ and $\lambda r_2^2<7/12$. For the latter sign, the cleared difference is
\[
 7(\lambda^2+\lambda+2)^2-12\lambda(\lambda+2)^2
 =4-12c+35c^2+30c^3+7c^4>0,
\]
because $35c^2-12c+4=35(c-6/35)^2+104/35$. For the former, $5\lambda^2-2\lambda-4>0$ at $23/20$ and increases. The bracket is therefore larger than
$2-5/7-(7/12)(1+20/63)=391/756>1/2$.

For the second claim, $r_3$ increases with $u$, and at $u=1/2$ its value is
$(\lambda^2+\lambda+2)/(\lambda^3+\lambda^2+\lambda+2)$. Thus $r_3<3/4$ and $\lambda r_3^2<2/3$: the first cleared difference is $3\lambda^3-\lambda^2-\lambda-2>0$, and the second is
\[
 2+29c^2+47c^3+35c^4+13c^5+2c^6>0.
\]
Consequently $2-r_3-\lambda r_3^2>7/12$. The ratio $K_b/(1-u)$ decreases as $1-u$ increases, giving
\begin{equation}\label{eq:near-K-ratio}
 \frac{K_b}{1-u}\le\frac{31}{12},\qquad
 \frac R{1-u}\le1+\frac{19}{12}\rho^2.
\end{equation}
At $\rho=1/2$,
\[
 \frac{uAR}{Q}\le\frac{67}{16}\frac{u(1-u)}{3-2u}\le\frac{67}{128}.
\]
The last inequality is $(2u-1)(4u-3)\ge0$. The exact comparison is
\[
 \rho U_3-V_4=\omega_3\left\{\rho Q(2-r_3-\lambda r_3^2)
             -uA\lambda r_3^2R(4/3-\rho)\right\}.
\]
At $1/2$, the expression in braces divided by $Q$ exceeds
$7/24-(5/6)(67/128)(2/3)=1/1152$. As $\rho$ increases, $\rho Q$ increases, whereas
\[
 \frac1{1-u}\frac\dd{\dd\rho}[R(4/3-\rho)]
 =-1+\frac\eta{1-u}(8\rho/3-3\rho^2)\le-11/144<0.
\]
Indeed the quadratic is at most $7/12$ on $[1/2,1]$, and $0<\eta/(1-u)\le19/12$. This proves the claim throughout that interval.
\end{proof}

\begin{lemma}[The range $0<\rho\le2/3$]\label{lem:near-low-rho}
One has $\Phi(K_b)>0$ whenever $0<\rho\le2/3$.
\end{lemma}
\begin{proof}
First suppose $\rho\le1/2$. The negative terms in~\eqref{eq:near-complete-Phi} are bounded using $f_s\le\rho^{s-1}$ and
\begin{equation}\label{eq:near-price-constants}
 \frac{W_2L_2}{r_1}\le\frac12,\qquad
 \frac{W_3L_3}{r_2}\le\frac14,\qquad
 \frac{W_4L_4}{r_3}\le\frac4{25}.
\end{equation}
Here $W_s=s\lambda^{s-1}/Z_s^2$. To verify all three bounds, put $k=1-\lambda u$ only in this calculation. Then $A\ge1+k$ and $A\ge2-k$. The first bound is $2k/A^2\le1/2$. For the next two,
\[
 Z_3^2-\lambda(A+k)^2=c(\lambda A^2-k^2)>0,
 \qquad Z_4-\lambda(A+2k)=c(\lambda A-k)>0.
\]
They reduce to $3k/[A(A+k)^2]\le1/4$ and $4k/[(A+k)(A+2k)^2]\le4/25$. For $k\le1/2$, use $A\ge2-k$ and monotonicity of the resulting one-variable fractions. For $k\ge1/2$, use $A\ge1+k$ and
\[
 (1+2k)^2(1+k)-12k=(2k-1)(2k^2+5k-1),
\]
\[
 (1+3k)^2(1+2k)-25k=(2k-1)(9k^2+15k-1).
\]

For the term $\Phi_R$ we have
\[
 \frac{\Phi_R}{A(1-u)}\ge
 \frac2\lambda\left[\frac{29}{48}(A-\lambda\rho)-\rho+
                              \frac{\lambda^2\rho^3}{A^2}\right].
\]
This follows from $Q\ge A(1-u)$, $R/(1-u)\le67/48$, and the identity
$(1-\lambda\rho/A)(AD-\rho-\lambda\rho^2/A)
=D(A-\lambda\rho)-\rho+\lambda^2\rho^3/A^2$.
The bracket decreases with $\rho\le1/2$ and, at $\rho=1/2$, increases with $A\ge1+\lambda/2$. Hence
\[
 \frac{\Phi_R}{A(1-u)}\ge\frac5{24\lambda}+\frac\lambda{(\lambda+2)^2}
 \ge\frac{461}{1764}.
\]
The last function decreases on $[1,3/2]$, as is seen by differentiation. By~\eqref{eq:near-price-constants}, the sum of the absolute values of the negative terms satisfies
\[
 \frac{2\rho\sum_{s=2}^4V_sf_s}{A(1-u)}
 \le\frac{67}{48}\left(\frac18+\frac1{32}+\frac1{100}\right)
 =\frac{8911}{38400}.
\]
Deleting only the positive squares in~\eqref{eq:near-complete-Phi} gives
$\Phi/[A(1-u)]>1/40$.

Now suppose $1/2\le\rho\le2/3$. In~\eqref{eq:near-merge}, Lemma~\ref{lem:near-age-comparisons} gives $U_2-V_3>0$ and $\rho U_3-V_4>0$. The only negative term left is $-2\rho^2V_2$, with
\[
 \frac{2\rho^2V_2}{A(1-u)}<\frac{322}{1215},
\]
because
\[
 \frac1{\lambda^2(\lambda+1)}-
       2u\frac{W_2L_2}{r_1}
 =\frac{[\lambda-(2\lambda+1)(1-\lambda u)]^2}
       {\lambda^2(\lambda+1)A^2}\ge0,
\]
$1/[\lambda^2(\lambda+1)]<7/20$, and $R/(1-u)\le46/27$.

To bound $\Phi_R$, we retain the positive terms of $Q$ that depend on $u$. The explicit formulas give
\[
 \frac Q{A(1-u)}\ge1+\frac u{1-u}\frac{\rho(1+2\rho)}3,
 \qquad \frac R{1-u}\le1+\rho^2\left(\frac{17}{24}+\frac u{1-u}\right).
\]
After substitution in~\eqref{eq:near-complete-finite-profile}, the coefficient of $u/(1-u)$ inside the braces for $\Phi_R/[A(1-u)]$ is at least
\[
 \frac\rho3\{A(2+\rho)-\rho(1+2\rho)(1+\rho)\}>\frac7{20}.
\]
To check the last bound, set $A=63/40$, where the expression is concave in $\rho\in[1/2,2/3]$ and has endpoint values $13/32$ and $434/1215$. Both exceed $7/20$. Since $E\ge3/7$ and $2/\lambda\ge4/3$, this part contributes at least $u/[5(1-u)]$.

For the other part, use
\[
 g(A,\rho)=\frac{37}{54}(A-\lambda\rho)-\rho+\frac{\lambda^2\rho^3}{A^2}.
\]
Its $\rho$ derivative is negative and its $A$ derivative exceeds $3/8$. It follows that
\[
 \frac{\Phi_R}{A(1-u)}\ge
 \frac1{27\lambda}-\frac{37}{162}+\frac{64\lambda}{27(\lambda+2)^2}
 +\frac34(1/2-u)+\frac15\frac u{1-u}.
\]
The first three terms increase with $\lambda\in[1,3/2]$ and are at least $35/486$. For the monotonicity, $\lambda^2(2-\lambda)/(\lambda+2)^3$ decreases to $9/343$, so its contribution to the derivative is still larger than the negative term. Finally, with $x=1/2-u$,
\[
 \frac34(1/2-u)+\frac15\frac u{1-u}-\frac{199}{1000}
 =\frac{(3/2)(x-2/125)^2+77/125000}{1+2x}>0.
\]
Thus $\Phi/[A(1-u)]>35/486+199/1000-322/1215
=1457/243000>1/200$.
\end{proof}

\subsubsection{Bounds for \texorpdfstring{$2/3\le\rho\le5/6$}{2/3 <= rho <= 5/6}}
Define the finite functions
\begin{equation}\label{eq:near-intermediate-functions}
 F(\rho,\lambda,u)=\Phi_R(\rho)-\frac{4\rho^2u(1-\lambda u)R}{A},
 \qquad g(\rho,\lambda,u)=\frac{(1-\lambda u)Q}{A^3}.
\end{equation}
Equation~\eqref{eq:near-merge} and Lemma~\ref{lem:near-age-comparisons} imply
\begin{equation}\label{eq:near-intermediate-lower}
 \Phi(K_b)>F(\rho,\lambda,u)+2\rho^3g(\rho,\lambda,u)
 \qquad(1/2\le\rho<1).
\end{equation}

\begin{lemma}[Bounds for $F$ and $g$ at $\rho=3/4$ and $\rho=5/6$]\label{lem:near-intermediate-arithmetic}
On $23/20\le\lambda\le3/2$, $0\le u\le1/2$, the explicitly defined rational functions in~\eqref{eq:near-intermediate-functions} satisfy
\[
 -\partial_uF(3/4,\lambda,u)>0,\quad
 -\partial_uF(5/6,\lambda,u)>0,\quad
 -\partial_ug(3/4,\lambda,u)>0.
\]
Their boundary values satisfy
\[
 F(3/4,\lambda,1/2)>1/8,\qquad
 F(5/6,\lambda,1/4)>0,
\]
\[
 F(5/6,\lambda,1/2)+\frac{27}{32}g(3/4,\lambda,1/2)>1/100.
\]
These are finite polynomial signs: the supplement forms each displayed difference from~\eqref{eq:near-intermediate-functions}, clears its positive powers of $\lambda$ and $A$, maps the stated rational rectangle to $[0,1]^2$, and lists its complete Bernstein coefficients. The exact positive factors, degrees and coefficient transform are recorded in Lemma~\ref{lem:near-finite-certificate}; no numerical evaluation is used to infer an interval sign.
\end{lemma}

\begin{lemma}[The range $2/3\le\rho\le5/6$]\label{lem:near-intermediate-rho}
One has $\Phi(K_b)>0$ throughout this range.
\end{lemma}
\begin{proof}
We first show that the expression in braces in $\Phi_R$ decreases with $\rho$. The explicit quadratic $K_b$ gives
\begin{equation}\label{eq:near-eta-bounds}
 0<\eta<3/4,\qquad 1<A\eta<4/3.
\end{equation}
For details, as $1-u$ increases, $\eta$ decreases, and its value at $u=1/2$ is $(13\lambda+25)/(48\lambda)<3/4$. The derivative of $A\eta$ with respect to $1-u$, multiplied by $24$, is
\[
 (11\lambda^2+1)/\lambda
 -2(\lambda^2+10\lambda-3)(1-u)+6\lambda(\lambda+1)(1-u)^2.
\]
Its discriminant is negative because
$-\lambda^4+46\lambda^3-28\lambda^2+66\lambda-3\ge159/2$.
Its minimum endpoint value gives
$A\eta-1\ge(13\lambda^2-45\lambda+50)/(96\lambda)>0$.
Its maximum endpoint value is
$(\lambda+1)(\lambda^2+2\lambda+13)/(24\lambda)
\le(35-3\lambda)/24<4/3$.

The bracket in $\Phi_R$ has exactly the five $\rho$ coefficients
\[
\begin{split}
 &A^2(1-u),\quad A(1-u)(2\lambda u-1),\quad
 2Au-\lambda(1-u)-A^2\eta,\\
 &-u[1+\lambda^2(1-u)/A],\qquad
 (\lambda u/A)(A\eta-1).
\end{split}
\]
Its derivative is negative on $[2/3,5/6]$. Indeed the quadratic coefficient is less than $-\lambda(1-u)\le-\lambda/2$. If the linear coefficient is positive, it is less than $2/3$, whereas twice $\rho$ times the quadratic coefficient is at most $-23/30$. The cubic and quartic contributions to the derivative sum to
\[
 u\rho^2\{-3[1+\lambda^2(1-u)/A]+4\lambda\rho(\eta-1/A)\}\le0,
\]
using $\lambda\rho\le5/4$, $1/A\ge2/5$, and~\eqref{eq:near-eta-bounds}. The first pair is already strictly negative, including $u=0$.

For $2/3\le\rho\le3/4$, Lemma~\ref{lem:near-intermediate-arithmetic} gives $F(3/4,\lambda,u)>1/8$. Since $F\le\Phi_R$, the bracket in $\Phi_R$ is positive at $3/4$ and therefore positive to its left. Both that bracket and $1-\lambda\rho/A$ decrease, so $\Phi_R$ decreases. The term $4\rho^2u(1-\lambda u)R/A$ subtracted from $\Phi_R$ increases, so $F$ decreases as well. Thus
$\Phi>F(\rho)\ge F(3/4)>1/8$.

For $3/4\le\rho\le5/6$, the same finite lemma gives
\[
 F(5/6,\lambda,u)+\frac{27}{32}g(3/4,\lambda,u)>1/100.
\]
We also need $\Phi_R(5/6)>0$. If $u\le1/4$, this follows from $F(5/6,\lambda,1/4)>0$ and the $u$ derivative. If $u\ge1/4$, concavity of $uA(1-u)$ and its two endpoint values give $uA(1-u)>1/3$. Since $R(5/6)\ge1-u$ and $Q(3/4)\le Q(1)=A$,
\[
 \frac{4(5/6)^2u(1-\lambda u)R(5/6)}A
 >\frac{27}{32}g(3/4).
\]
After division by $2(1-\lambda u)/A^3$, this follows from $25/54>27/64$. The preceding mixed bound now gives $\Phi_R(5/6)>0$.
As before, $\Phi_R(5/6)>0$ implies that $F$ decreases with $\rho$ on $[3/4,5/6]$. The function $g$ increases with $\rho$, since $Q$ does. Equation~\eqref{eq:near-intermediate-lower} therefore gives
\[
 \Phi(K_b)>F(5/6)+\frac{27}{32}g(3/4)>1/100.
\]
\end{proof}

\subsubsection{Bounds for \texorpdfstring{$5/6\le\rho<1$}{5/6 <= rho < 1}}
For $5/6\le\rho<1$, we need a lower bound that also applies when the quadratic term in $1-T$ is negative. We estimate it together with the value of $\Phi(K_b)$ at $\rho=1$ and the cubic and quartic terms in $1-T$.

Put $n_s=s(s-1)/2$ and $B_s^{\rm age}=2U_s-V_s$. The following are finite polynomials in $\rho$:
\begin{equation}\label{eq:near-core-defs}
\begin{gathered}
 \Psi(\rho)=\Phi_R(\rho)+2\rho\sum_{s=2}^4(U_s-V_s),\qquad p=\Psi(1),\\
 L_{\rm core}(\rho)=\frac{\Psi(\rho)-p}{1-\rho}
                       -2\rho\sum_{s=2}^4n_s B_s^{\rm age}(\rho)
                   =\sum_{j=0}^4c_j\rho^j,\qquad C_0=L_{\rm core}(1),\\
 H(\rho)=-c_2-(1+\rho)c_3-(1+\rho+\rho^2)c_4.
\end{gathered}
\end{equation}
The quotient is taken by polynomial division, including at $\rho=1$. In particular,
\begin{equation}\label{eq:near-core-chord}
 L_{\rm core}(\rho)=(1-\rho)L_{\rm core}(0)+\rho C_0+
                         \rho(1-\rho)H(\rho).
\end{equation}
Define the three further rational functions
\begin{equation}\label{eq:near-recovery-defs}
\begin{gathered}
 b_{\rm env}=\frac53\left[2B_3^{\rm age}(5/6)+\frac{27}{4}B_4^{\rm age}(5/6)\right],\\
 j_{\rm env}=\frac53\left[U_2(5/6)+\frac{169}{36}U_3(5/6)
                                  +\frac{2601}{256}U_4(5/6)\right],\\
 K_{5/6}=L_{\rm core}(0)-C_0+\frac56H(5/6),\\
 J_1=\frac{25p+C_0}{(1-\lambda u)A},\qquad
 J_2=4b_{\rm env}+3C_0,\qquad
 J_3=K_{5/6}+j_{\rm env}+4C_0.
\end{gathered}
\end{equation}
Here $U_s$ and $V_s$ for $s=2,3,4$ are evaluated at the same $K_b$ and $u$.

The following Bernstein coefficient bounds also give the inequalities of Lemma~\ref{lem:near-intermediate-arithmetic}. To clear the denominators, write
\[
 D_{\rm profile}=A^3Z_3^3Z_4^3Z_5^2.
\]
For the two $J_1$ boundary derivatives, use
\[
\begin{split}
 D_0&=12\lambda^3 Z_3(\lambda,0)^4Z_4(\lambda,0)^4Z_5(\lambda,0)^3,\\
 D_{1/2}&=24\lambda^3(2-\lambda)^2(\lambda+2)^5
 (\lambda^2+\lambda+2)^4(\lambda^3+\lambda^2+\lambda+2)^4
 (\lambda^4+\lambda^3+\lambda^2+\lambda+2)^3.
\end{split}
\]
Every one of these factors is strictly positive on the intervals below.

\begin{lemma}[Finite coefficient signs]\label{lem:near-finite-certificate}
For each row of the following table, let $P$ be the function in its first column multiplied by the positive denominator in its second column. Then $P$ is a polynomial of the indicated bidegree, and every Bernstein coefficient of $P$ on the indicated rectangle $[a,b]\times[0,1/2]$ is strictly positive. The functions $F,g$ are exactly~\eqref{eq:near-intermediate-functions}, and $p,c_j,J_i$ are exactly~\eqref{eq:near-core-defs}--\eqref{eq:near-recovery-defs}.
\begin{center}\small
\setlength{\tabcolsep}{3pt}
\begin{tabular}{@{}llll@{}}
Function & Positive denominator & Bidegree & $[a,b]$\\ \hline
$-\partial_uF(3/4)$ & $4096\lambda^2A^3$ & $(7,6)$ & $[23/20,3/2]$\\
$F(3/4,\lambda,1/2)-1/8$ & $8192\lambda^2(\lambda+2)^2$ & $(5,0)$ & $[23/20,3/2]$\\
$-\partial_uF(5/6)$ & $93312\lambda^2A^3$ & $(7,6)$ & $[23/20,3/2]$\\
$-\partial_ug(3/4)$ & $16A^4$ & $(3,2)$ & $[23/20,3/2]$\\
$F(5/6,\lambda,1/4)$ & $497664\lambda^2(3\lambda+4)^2$ & $(6,0)$ & $[23/20,3/2]$\\
$\begin{gathered}F(5/6,\lambda,1/2)\\+\frac{27}{32}g(3/4,\lambda,1/2)-\frac1{100}\end{gathered}$
 & $4665600\lambda^2(\lambda+2)^3$ & $(6,0)$ & $[23/20,3/2]$\\
$-c_2$ & $12\lambda D_{\rm profile}$ & $(29,14)$ & $[23/20,10/7]$\\
$-(c_3+4c_4)/u$ & $12\lambda D_{\rm profile}$ & $(28,13)$ & $[23/20,10/7]$\\
$-\partial_u^2J_1$ & $6\lambda(1-\lambda u)^3A^6Z_3^5Z_4^5Z_5^4$ & $(51,23)$ & $[23/20,707/500]$\\
$\partial_\lambda J_1(\lambda,0)-1/4$ & $D_0$ & $(35,0)$ & $[23/20,707/500]$\\
$\partial_\lambda J_1(\lambda,1/2)$ & $D_{1/2}$ & $(40,0)$ & $[23/20,707/500]$\\
$J_2$ & $216\lambda^2D_{\rm profile}$ & $(30,15)$ & $[23/20,3/2]$\\
$J_3-(1-\lambda u)/20$ & $2488320\lambda^2D_{\rm profile}$ & $(30,15)$ & $[23/20,10/7]$\\
$p-1/1000$ & $3000\lambda^2D_{\rm profile}/A$ & $(29,14)$ & $[15/13,707/500]$
\end{tabular}
\end{center}
The division by $u$ in the eighth row is exact polynomial cancellation before setting $u=0$. The two additional rational evaluations are
\[
 J_1(23/20,0)>57/100,
 \qquad J_1(23/20,1/2)>21/25.
\]
More explicitly, the finite coefficient obligation in each row is as follows. If
\[
 P(a+(b-a)x,y/2)=\sum_{k=0}^M\sum_{l=0}^Np_{kl}x^ky^l,
\]
then for every $0\le i\le M$, $0\le j\le N$,
\begin{equation}\label{eq:near-Bernstein-obligation}
 \sum_{k=0}^i\sum_{l=0}^j
 p_{kl}\frac{\binom{i}{k}}{\binom{M}{k}}
        \frac{\binom{j}{l}}{\binom{N}{l}}>0.
\end{equation}
The supplement supplies the full coefficients and the exact rational verification of these finite inequalities.
\end{lemma}
\begin{proof}
Only finite addition, multiplication and division by positive integers are needed to generate~\eqref{eq:near-Bernstein-obligation} from the displayed definitions. The complete generated data are provided in the supplement. Once these finite signs are checked, the passage to the entire rectangle is immediate from
\[
 P(a+(b-a)x,y/2)=\sum_{i=0}^M\sum_{j=0}^N b_{ij}
 \binom Mi x^i(1-x)^{M-i}\binom Nj y^j(1-y)^{N-j}.
\]
The basis functions are nonnegative and sum to one, including on the boundary. Positivity of all $b_{ij}$ proves every claimed sign. All denominators listed in the table are positive; thus no multiplication changes an inequality's direction.
\end{proof}

\begin{lemma}[The range $5/6\le\rho<1$]\label{lem:near-high-rho}
One has $\Phi(K_b)>0$ throughout this range.
\end{lemma}
\begin{proof}
For $s=2,3,4$ and $u>0$, all three coefficients of $B_s^{\rm age}$ as a polynomial in $\rho$ are positive. Dividing them by $\omega_s>0$ gives
\[
\begin{split}
 &[\rho^0]:\ A(1-u)[2(2-r_s-\lambda r_s^2)+u(2\lambda r_s^2-1/r_{s-1})],\\
 &[\rho^1]:\ 2\lambda u(1-u)(2-r_s-\lambda r_s^2),\\
 &[\rho^2]:\ u[2(2-r_s-\lambda r_s^2)+A\eta(2\lambda r_s^2-1/r_{s-1})].
\end{split}
\]
Here $r_2<5/7$, $r_3<3/4$, $r_4<5/6$, $1/r_2\le19/10$, $1/r_3\le65/38$, and $r_s>1/2$. These bounds follow from monotonicity in $u$ and $\lambda$ and the finite recursion for $Z_s$; for $r_4<5/6$, the endpoint difference is $c(5\lambda^3+4\lambda^2+3\lambda+2)>0$. The constant brackets are respectively larger than $1/7,1/20,7/114$. The linear coefficients are positive.

For $s=3,4$, if $2\lambda r_s^2-1/r_{s-1}<0$, then~\eqref{eq:near-eta-bounds} bounds the brackets in the quadratic coefficients below by $2/15$ and $1/19$, respectively. If $2\lambda r_s^2-1/r_{s-1}\ge0$, both terms in the bracket are nonnegative and the first is positive. At $s=2$, put $\overline H=(35-3\lambda)/24$. Since $A\eta\le\overline H$, its bracket is at least
\[
 f(A)=4-2r_2+2\lambda(\overline H-1)r_2^2-\overline H A,
 \qquad r_2=\frac A{(\lambda+1)A-\lambda}.
\]
This function decreases with $A$, since
$f'(A)\le2\lambda/Z_3^2-\overline H<3/4-61/48<0$.
At $A=\lambda+1$, its numerator after multiplication by $24(1+\lambda+\lambda^2)^2$ is
\[
 3c^6-8c^5-85c^4-206c^3-169c^2+14c+64>0.
\]
Indeed $0<c\le3/7$, and deleting the positive terms and using
$c^2<1/5$, $c^3<2/25$, $c^4<7/200$, $c^5<3/200$ leaves $85/8>0$. This proves the claimed signs of the coefficients of $B_s^{\rm age}$.

The coefficient of $\rho^5$ in $\Phi_R$ gives
\[
 c_4=-[\rho^5]\Phi_R=\frac{2\lambda u}{A^2}(A\eta-1)>0.
\]
Lemma~\ref{lem:near-finite-certificate} gives $c_2<0$ and $c_3+4c_4<0$. Therefore $H$ in~\eqref{eq:near-core-defs} is positive and increases on $[0,1]$. By~\eqref{eq:near-core-chord}, the chord joining the endpoint values of $L_{\rm core}$ is a lower bound for $L_{\rm core}$.

Let $x=1-T$, so $\rho=1-x^2$ and $0\le x\le1/\sqrt6<5/12$. Directly,
\[
 f_2=1-x^2,\quad f_3=1-x^2(3-2x),\quad
 f_4=1-x^2(6-8x+3x^2).
\]
Substitution in the complete expression~\eqref{eq:near-complete-Phi} gives the exact identity
\begin{equation}\label{eq:near-recovery-identity}
 \Phi=p+x^2[L_{\rm core}(\rho)+xB_*+x^2J_*],
\end{equation}
where
\[
 B_*=2\rho[2B_3^{\rm age}+(8-3x)B_4^{\rm age}],
\]
\[
 J_*=2\rho[U_2+U_3(3-2x)^2+U_4(6-8x+3x^2)^2].
\]
Their coefficients are positive. Since $\rho\ge5/6$ and $x\le5/12$, they satisfy $B_*\ge b_{\rm env}$ and $J_*\ge j_{\rm env}$. Combining~\eqref{eq:near-core-chord} and the increase of $H$ gives
\begin{equation}\label{eq:near-quartic}
 \Phi\ge p+C_0x^2+b_{\rm env}x^3+(K_{5/6}+j_{\rm env})x^4.
\end{equation}

The finite signs now have three distinct roles. First $p>1/1000$; its rectangle begins at $15/13<2/\sqrt3$. Second, $J_1$ is concave in $u$, and both boundary functions increase with $\lambda$. Its two rational endpoint values therefore imply
\[
 J_1(\lambda,u)\ge(1-2u)J_1(\lambda,0)+2uJ_1(\lambda,1/2)>57/100.
\]
In particular $25p+C_0>0$. Third, $J_2>0$ and $J_3>(1-\lambda u)/20>0$ give
\[
 b_{\rm env}>-3C_0/4,\qquad K_{5/6}+j_{\rm env}>-4C_0.
\]
The complete bound~\eqref{eq:near-quartic} is consequently larger than
\[
 p+C_0(x^2-3x^3/4-4x^4).
\]
On $0\le x\le1/\sqrt6$, set $G(x)=x^2-3x^3/4-4x^4$. Then
\[
 0\le G(x)<1/25,\qquad 1-25G(x)\ge37639/640000.
\]
Here $G\ge0$ follows from $x<33/80$ and $1-3(33/80)/4-4(33/80)^2=1/100$. For the other bound, $1-25G$ decreases up to $1/4$ and has value $31/256$ there. Above $1/4$ its second derivative exceeds $50$. Its value and first derivative at $3/10$ are $53/800$ and $69/80$; the quadratic Taylor lower bound has minimum $37639/640000$. This proves the two assertions without an extremum assumption for $\Phi$.

If $C_0\ge0$, the preceding bound is at least $p$. If $C_0<0$, $C_0>-25p$ gives
\[
 \Phi>p[1-25G(x)]>\frac{37639}{640000000}>0.
\]
This proves the remaining moment range.
\end{proof}

Lemmas~\ref{lem:near-low-rho}, \ref{lem:near-intermediate-rho} and~\ref{lem:near-high-rho} cover $0<\rho<1$. Lemma~\ref{lem:near-K-affinity}, the positivity of its denominator, and Lemma~\ref{lem:near-four-age} therefore prove $D_\lambda>0$ whenever $u\le1/2$ in~\eqref{eq:near-domain}.

\subsubsection{The first two terms for \texorpdfstring{$1/2<u<\lambda^{-2}$}{1/2 < u < lambda^(-2)}}
It suffices to keep the terms with $s=1,2$ in the lower bound for $D_\lambda$. We begin with $K_{\rm cap}$ on the domain
\begin{equation}\label{eq:near-high-u-domain}
 1/2<u<\lambda^{-2}\le3/4,\qquad uA<\rho<1,\qquad z=\rho/A.
\end{equation}
Lemma~\ref{lem:near-absorption}, with the exact second moment $a_2=uz$, and the bound $a_3/a_2<r_2$ imply
\begin{equation}\label{eq:near-H2}
 D_\lambda>u^2H_2,
 \quad H_2=\frac{(1+w)(1-\lambda z)^2-2(1-a)(1-\lambda z)}\lambda
          +2\lambda z^2[B_2+2L_2M_2],
\end{equation}
where
\[
 M_2=a(1-r_2)+\frac wc(\lambda+1-\lambda r_2-1/z).
\]
All discarded terms, those with $s\ge3$, are positive. The term with $s=1$ and the value $a=a(w)$ are the same as in the full series.

For the polynomial calculations below we use $k$, previously $c/(\lambda\log\lambda)$, for the coordinate $1-\lambda u$. Thus
\[
 A=\lambda+k,\qquad B=\lambda^2+(\lambda+1)k,
 \qquad u=(1-k)/\lambda.
\]
Put
\begin{equation}\label{eq:near-high-denominator}
 E=\frac{K_{\rm cap}-1+u}{1-u},\qquad
 D_\rho=2-(\lambda+1)u+\frac{2\lambda u}{A}\rho
 +\left[\frac{2u}{A(1-u)}-(\lambda+1)uE\right]\rho^2.
\end{equation}
Then $w=cu(1+E\rho^2)/D_\rho$. The denominator $D_\rho$ is positive on $uA\le\rho\le1$. More precisely, with
\[
 E_+=\frac{2/(1+\lambda u)-(1-u)}{1-u},
\]
one has
\begin{equation}\label{eq:near-high-denominator-gap}
 D_\rho-(1-u)(1+E\rho^2)
 =(1-\rho)\left[k(1+\rho)+\frac{2\lambda u\rho}{A}\right]
 +(1+\lambda u)(E_+-E)\rho^2>0.
\end{equation}
Thus $0<w<h$ for $uA\le\rho\le1$.

We first exclude $\lambda\ge11/8$ when $u>1/2$ and $\alpha<1$. Differentiating $w$ with respect to $\rho$ gives
\[
 w_\rho=\frac{cu}{D_\rho^2}
 \left[-\frac{2\lambda u}{A}+4\left(E-\frac{u}{A(1-u)}\right)\rho
                      +\frac{2\lambda uE}{A}\rho^2\right]>0.
\]
Indeed $E\ge u/(1-u)$ and $\rho\ge uA>1/2$ give the positive lower bracket
$(2\lambda u/A)[-1+2\rho+u\rho^2/(1-u)]$.
Writing $J=1/u-(\lambda+1)+\lambda z$, we have $a=(1-wJ/c)/z$. Here $J\ge(1-u)k/u>0$ and
\[
 a+\lambda w/c=\frac{1+(w/c)(\lambda+1-1/u)}z>0.
\]
Consequently
\[
 a_\rho=-\frac{a+\lambda w/c+AJw_\rho/c}{\rho}<0.
\]
At $\rho=1$ a direct substitution gives
\[
 a(1)-1=\frac{(1-u)[2\lambda-(\lambda+1)K_{\rm cap}]}
                     {2-(\lambda+1)uK_{\rm cap}}.
\]
Since $a<\alpha<1$ and $a$ decreases with $\rho$, we must have
$K_{\rm cap}>2\lambda/(\lambda+1)$. On the other hand $K_{\rm cap}$ decreases with $u>1/2$, and the logarithmic bounds in~\eqref{eq:near-cap} give
\[
 K_{\rm cap}<\frac{25(\lambda+1)}{48\lambda}+\frac14.
\]
For $\lambda\ge11/8$ this is smaller than $2\lambda/(\lambda+1)$, since
\[
 59\lambda^2-62\lambda-25
 =59(\lambda-11/8)^2+\frac{401}{4}(\lambda-11/8)+\frac{83}{64}>0.
\]
Thus $u>1/2$ and $\alpha<1$ imply $\lambda<11/8$.

It remains to prove $H_2>0$ for $2/\sqrt3\le\lambda\le11/8$. Collecting the terms in~\eqref{eq:near-H2} according to their dependence on $w$ gives
\[
 H_2=H_0(z)+wH_v(z),
\]
where
\begin{equation}\label{eq:near-H0-Hv}
\begin{split}
 H_0(z)&=\frac2{\lambda z}-2-\frac1\lambda
       +z^2\left[\lambda+\frac{2k^2}{B^2}
                       +\frac{4\lambda^2(1-u)k}{AB}\right],\\
 H_v(z)&=\frac{(1-\lambda z)^2}{\lambda}
  -\frac{2(1-\lambda z)J}{\lambda cz}+\frac{2z^2k^2}{B^2}
  +\frac{4kz^2}{Ac}(\lambda+1-\lambda A/B-1/z).
\end{split}
\end{equation}
In particular $N(\rho)=\rho D_\rho H_2(\rho)$ is a polynomial of degree five. Its coefficients $N=\sum_{j=0}^5n_j\rho^j$ are
\begin{equation}\label{eq:near-N-coefficients}
\begin{gathered}
 P=\lambda+2k^2/B^2+4\lambda^2(1-u)k/(AB),\\
 C_2=\lambda+2k^2/B^2+4k(\lambda+1-\lambda A/B)/(Ac),\\
 W=\lambda-2k/B+2k^2/B^2,\qquad
 S_0=[2-(\lambda+1)u]P+cuC_2,\\
 n_0=2A/\lambda,\qquad n_1=4u-2-2/\lambda,\\
 n_2=-2AE/\lambda+4u/[\lambda(1-u)]-4uB/A^2,\\
 n_3=2E+S_0/A^2-2u(2\lambda+1)/[\lambda A(1-u)],\\
 n_4=(2u/A^3)[\lambda P+EA(A-2k)],\qquad
 n_5=(2u/A^2)[P/(A(1-u))-EW].
\end{gathered}
\end{equation}
These are obtained by multiplying the quadratic $D_\rho$ by $\rho H_0$ and the quadratic $cu(1+E\rho^2)$ by $\rho H_v$; thus the displayed six coefficients include all terms.

\begin{lemma}[Positivity of $N'''(\rho)$]\label{lem:near-high-curvature}
For $2/\sqrt3\le\lambda\le11/8$ and~\eqref{eq:near-high-u-domain}, $N'''(\rho)>0$ on $uA\le\rho\le1$.
\end{lemma}
\begin{proof}
First $n_4>0$, since $A-2k=c+\lambda u>0$. Also $0<W\le\lambda$. Put $G=1+\lambda u$ and $C=A(2/G-1+u)$. The exact identities
\[
 P-W=\frac{2k}{B}\left(1+\frac{2\lambda^2(1-u)}A\right),\quad
 C-1=\frac{k[\lambda(1-u)^2+u]}G
\]
give
\[
 P-CW=\frac kG\left\{\frac{2G}{B}
       \left(1+\frac{2\lambda^2(1-u)}A\right)
       -[\lambda(1-u)^2+u]W\right\}>0.
\]
Indeed $u\le3/4$, $1/4\le1-u\le1/2$, $A\le27/16$, $B\le337/128$, $G>3/2$, and $\lambda^2\ge4/3$. The first term in braces is at least $(384/337)(113/81)>385/256$, whereas the second is at most $385/256$. Since $E<E_+$, this proves $n_5>0$.

Using $P,C_2\ge\lambda$, $S_0\ge2\lambda(1-u)$, and $\rho\ge uA$ in~\eqref{eq:near-N-coefficients}, we obtain
\begin{equation}\label{eq:near-third-derivative-bound}
 N'''>\frac{12u}{1-u}\frac{R_0}{\lambda A}
              +\frac{12\lambda(1-u+4\lambda u^2)}{A^2},
\end{equation}
where
\[
 R_0=\lambda^2-\lambda-1+\lambda^2(4u^3-u)+4\lambda cu^2.
\]
In obtaining this bound, replace $E$ by $u/(1-u)$ before combining it with the negative term: its coefficient is $12[1+4u^2(A-2k)/A]>0$. The sign of $R_0$ is not used in this replacement.

The function $R_0$ increases with both variables, since
\[
 (R_0)_u=\lambda^2(12u^2-1)+8\lambda cu>0,
\]
\[
 (R_0)_\lambda=2\lambda(1-u+4u^2+4u^3)-(1+4u^2)
 \ge1-2u+4u^2+8u^3\ge2.
\]
On $1/2\le u\le3/5$, $R_0\ge(5-4\sqrt3)/3>-2/3$, $u/(1-u)\le3/2$ and $\lambda A>5/3$. The first term of~\eqref{eq:near-third-derivative-bound} is therefore greater than $-36/5$; the second is at least
$12(23/20)(23/20+2/5)/(27/16)^2>36/5$.
On $3/5\le u\le2/3$, use $R_0\ge-5749/25000>-1/4$, $u/(1-u)\le2$ and $\lambda A>3/2$. The two terms are respectively greater than $-4$ and $4$. Finally, on $u\ge2/3$,
$R_0\ge(275-150\sqrt3)/81>0$. This proves the result in every case.
\end{proof}

The next lemma gives the signs needed to bound $N'(1)$ and $N(1)$. It verifies inequalities on a rectangle in $\lambda,k$, with $u=(1-k)/\lambda$, which also includes parameter values outside~\eqref{eq:near-high-u-domain}.

\begin{lemma}[Polynomial inequalities for $N'(1)$ and $N(1)$]\label{lem:near-high-arithmetic}
Let $23/20\le\lambda\le11/8$, $0\le k\le1/2$, $u=(1-k)/\lambda$, $A=\lambda+k$, $B=\lambda^2+(\lambda+1)k$, and let $n_j(E)$ be exactly~\eqref{eq:near-N-coefficients}. Define
\[
 Y(E)=\sum_{j=1}^5j n_j(E),\qquad
 D=\lambda A^4(c+k)B^2>0,
\]
\[
 P_-=-\tfrac12D\,Y\left(\frac u{1-u}\right),\qquad
 P_+=-\tfrac12(2-k)D\,Y\left(\frac{2/(2-k)-(1-u)}{1-u}\right).
\]
These are polynomials in $(\lambda,k)$ of bidegrees $(10,7)$ and $(10,8)$, respectively, and both are strictly positive. Also define
\begin{equation}\label{eq:near-upper-S}
\begin{split}
 S(K)={}&k(1-uK)(1+2\lambda/B^2)-A(1-u)[(\lambda+1)K-2\lambda]\\
 &+\frac{2\lambda^2(1-u)}{AB}[2\lambda-(\lambda+1)uKA].
\end{split}
\end{equation}
The polynomial $24\lambda^2AB^2S(K_b)$ has bidegree $(10,7)$ and is strictly positive on the same rectangle.

The exact finite obligations are positivity of all Bernstein coefficients of $P_-$ and $24\lambda^2AB^2S(K_b)$ on that rectangle, and positivity of all Bernstein coefficients of $P_+$ on the three rectangles
\[
 [23/20,101/80]\times[0,1/4],\quad
 [23/20,101/80]\times[1/4,1/2],\quad
 [101/80,11/8]\times[0,1/2].
\]
The coefficients are generated by~\eqref{eq:near-Bernstein-obligation} with the corresponding affine substitution in each coordinate. They and their exact verification are in the supplement. Every clearing factor here is positive, including $2-k$, $c+k$, and $1-u=(c+k)/\lambda$.
\end{lemma}

Because $Y$ is affine in $E$, Lemma~\ref{lem:near-high-arithmetic} and
$u/(1-u)<E<E_+$ give $N'(1)=Y(E)<0$. Combined with Lemma~\ref{lem:near-high-curvature}, this reduces the sign of $N$ to its two endpoints. Indeed $N'$ is strictly convex, and its nonpositive set is an interval containing the right endpoint. Thus $N'$ is either nonpositive throughout or is positive first and nonpositive afterwards; $N$ cannot have an interior minimum. Consequently
\begin{equation}\label{eq:near-high-endpoint-reduction}
 \min_{uA\le\rho\le1}N(\rho)=\min\{N(uA),N(1)\}.
\end{equation}

At the lower endpoint $z=u$, the expression $H_2$ is affine in $w$. Its values at $w=0$ and $w=h$ are, respectively,
\[
 H_0(u)=\frac{k^2}{\lambda}+\frac{2(1-u)k}{\lambda u}
       +2\lambda u^2\left[\frac{k^2}{\lambda B^2}
                              +\frac{2\lambda(1-u)k}{AB}\right]>0,
\]
\[
 H_h(u)=\frac k\lambda\left\{(1+h)k(1+2\lambda u^2/B^2)+2c
       +\frac{4\lambda u^2[\lambda c-(\lambda^2+1)k]}{AB}\right\}.
\]
We verify the latter sign. The positive function
\[
 R(\lambda,k)=\frac{\lambda^2+1}{\lambda(\lambda+k)[\lambda^2+(\lambda+1)k]}
\]
decreases with $\lambda$: its logarithmic derivative is
\[
 \frac{2\lambda}{\lambda^2+1}-\frac1\lambda-\frac1{\lambda+k}
                        -\frac{2\lambda+k}{\lambda^2+(\lambda+1)k}<0,
\]
since $\lambda k<1$. Also $h=c(1-k)/(c+k)$ increases with $c$. Dropping two positive terms gives
\[
 \frac{\lambda H_h(u)}k>
 2c+(1+h)k-4k(1-k)^2R(\lambda,k)
\]
\[
 \ge\frac3{10}+\left[1+\frac{(3/20)(1-k)}{3/20+k}\right]k
       -4k(1-k)^2R(23/20,k).
\]
After multiplication by the positive denominator
$(23/20)(23/20+k)[(23/20)^2+(43/20)k](3/20+k)$, the last expression equals
\[
 Q(k)=\frac{2518569}{32000000}-\frac{29541}{200000}k
 -\frac{7316903}{3200000}k^2+\frac{1790369}{80000}k^3-\frac{57507}{8000}k^4.
\]
For $0\le k\le1/2$, coefficient comparison and $k^4\le k^3/2$ give
\[
 Q(k)>\frac3{40}-\frac3{20}k-\frac52k^2+18k^3
 =18(k+1/9)(k-1/8)^2+\frac{11}{160}k+\frac7{160}>0.
\]
Hence $H_h(u)>0$. Equation~\eqref{eq:near-high-denominator-gap} gives $0<w<h$, so at $z=u$ the value of $H_2$ is a convex combination of $H_0(u)$ and $H_h(u)$ with positive weights. Thus $N(uA)>0$.

At the upper endpoint, direct substitution in~\eqref{eq:near-H2} gives
\[
 N(1)=\frac{2k}{\lambda A^2(1-u)}S(K_{\rm cap}),
\]
with $S$ from~\eqref{eq:near-upper-S}. Its coefficient of $K$ is negative. More explicitly, if $\delta=(\lambda+1)K-2\lambda$, then $S=P-\delta Q$ with
\[
 P=k\left[\frac{c+2k}{\lambda+1}(1+2\lambda/B^2)
                     +\frac{4\lambda^3(1-u)^2}{AB}\right]>0,
\]
\[
 Q=A(1-u)+\frac{uk}{\lambda+1}(1+2\lambda/B^2)
                          +\frac{2\lambda^2u(1-u)}B>0.
\]
The specific logarithmic bound $K_{\rm cap}<K_b$ therefore gives
$S(K_{\rm cap})>S(K_b)>0$, the latter sign being Lemma~\ref{lem:near-high-arithmetic}. Thus $N(1)>0$.
Equation~\eqref{eq:near-high-endpoint-reduction} proves $H_2>0$ and, by~\eqref{eq:near-H2}, $D_\lambda>0$ for $1/2<u<\lambda^{-2}$.

The bounds for $u\le1/2$ and $u>1/2$ cover $0<\alpha<1$ under~\eqref{eq:near-domain}. Lemma~\ref{lem:calibration} covers $\alpha=1$. Thus $D_\lambda>0$ throughout~\eqref{eq:near-domain}, and Lemma~\ref{lem:full-response} gives $v'(\lambda)<0$ there.

\subsection{Kernel and integral bounds for \texorpdfstring{$707/500\le\lambda<(3+\sqrt5)/2$}{707/500 <= lambda < (3+sqrt(5))/2}}\label{sec:middle-high}

We prove $v'(\lambda)<0$ on the two intervals
\[
 \frac{707}{500}\le\lambda\le\frac{39}{20},
 \qquad
 \frac{39}{20}\le\lambda<\frac{3+\sqrt5}{2},
 \qquad \lambda<m.
\]
Fix the offspring law. Lemma~\ref{lem:canonical-moments} gives the moment inequalities, and Lemma~\ref{lem:full-response} gives the derivative formula
\[
 D=2\eta(\Psi)-C,\qquad
 \eta(f)=\E[b(X)(\mathcal Rf)(X)],\qquad v'(\lambda)=-D/Z,
\]
where $Z>0$. Lemma~\ref{lem:calibration} gives $D>0$ when $(\lambda+1)h\ge1$. It therefore suffices to consider
\begin{equation}\label{eq:mh-hard}
 0<h<\frac1{\lambda+1},\qquad
 0<a_1=\E U<\lambda^{-2},\qquad U=\frac{X}{c+X},\quad c=\lambda-1.
\end{equation}
We use $\inf\Psi>0$ and $v<h$ from Lemma~\ref{lem:harmonic-speed}. Recall that $a_n=\E U^n$, $\beta_n=h/(\lambda^n-1+h)$ and $\rho_n=a_n/\beta_n$. The polynomial bounds below hold on domains that contain the parameter values obtained from the escape law.

\subsubsection{Kernel bounds for \texorpdfstring{$707/500\le\lambda\le39/20$}{707/500 <= lambda <= 39/20}}
Define
\[
 P_v(y)=y^2+c\{(1+v)y-v\},\qquad
 P_c(y)=(1-y)\{2v-(1+v)y\}.
\]
Thus $\Psi=\cJ P_v$, and~\eqref{eq:C-kernel-identity} gives
\begin{equation}\label{eq:mh-C-kernel}
 C=\E[b(X)(\cJ P_c)(X)].
\end{equation}
The measure comparison of Lemma~\ref{lem:entropy-response} gives
\begin{equation}\label{eq:mh-max-source}
 \eta\ge\max(g,b)\pi_\lambda,
 \qquad
 D\ge\E\left[2\max\{g(X),b(X)\}\Psi(X)-b(X)(\cJ P_c)(X)\right].
\end{equation}
All functions in the last expectation are bounded. The identity for one child
\[
 g(F_\lambda(z))-b(F_\lambda(z))=F_\lambda'(z)g(z)
\]
and the signs $g(z)>0$ for $0<z<1$, $g(1)=0$, and $g(z)<0$ for $z>1$ show that
\begin{equation}\label{eq:mh-source-split}
 g(x)\ge b(x)\quad\left(0<x\le\frac1{\lambda+1}\right),
 \qquad
 b(x)\ge g(x)\quad\left(\frac1{\lambda+1}\le x<1\right).
\end{equation}
Consequently, for $707/500\le\lambda\le39/20$, it suffices to prove
\begin{align}\label{eq:mh-two-kernels}
 2\frac{g(x)}{b(x)}\Psi(x)-(\cJ P_c)(x)&>0,
 &&0<x\le\frac1{\lambda+1},\notag\\
 2\Psi(x)-(\cJ P_c)(x)&>0,
 &&\frac1{\lambda+1}\le x<1.
\end{align}
Multiplication by $b$ gives the bounded integrand in~\eqref{eq:mh-max-source}, even if the ratio $g/b$ is unbounded.

\subsubsection{A four-moment bound on \texorpdfstring{$v/(1+v)$}{v/(1+v)}}
Put $s=a_2/a_1$ for the kernel estimates on $707/500\le\lambda\le39/20$. The moment comparison and Jensen's inequality give
\begin{equation}\label{eq:mh-moment-domain}
 \begin{gathered}
 a_1\le s\le\frac1{D_2},\qquad D_2=\lambda+1-\lambda a_1,\\
 D_3=\lambda^2+\lambda+1-\lambda(\lambda+1)a_1,\\
 D_4=\lambda^3+\lambda^2+\lambda+1-\lambda(\lambda^2+\lambda+1)a_1.
 \end{gathered}
\end{equation}
and
\[
 a_3\le a_2D_2/D_3,\qquad a_4\le a_2D_2/D_4.
\]
All three denominators are positive; in particular $D_2>\lambda$. Write
\begin{equation}\label{eq:mh-four-envelope}
 \begin{split}
 A(a_1)&=c\lambda+c\lambda^2(D_2/D_3)^2+\lambda^4(D_2/D_4)^2,\\
 E_4&=c+A(a_1)s^2,\qquad
 B=1-\lambda a_1+\lambda a_1s(1-\lambda s)>0.
 \end{split}
\end{equation}
Here $A(a_1)$ is a scalar moment bound, distinct from the pair denominator $A_\lambda(x,y)$.

To bound the terms of order four and higher, set
\[
 T=\sum_{n\ge1}\lambda^{n-1}a_n^2,
 \qquad Q=\sum_{n\ge1}\lambda^{n-1}a_na_{n-1},\qquad a_0=1.
\]
Since $0<U<1/\lambda$, both series converge absolutely. Expanding the pair denominator gives
\begin{equation}\label{eq:mh-moment-speed}
 v=\frac{cT}{2Q-(\lambda+1)T},\qquad
 \frac{v}{1+v}=\frac{cT}{2(Q-T)}.
\end{equation}
For $n\ge4$, $a_n\le a_4\lambda^{-(n-4)}$, whence
\[
 \sum_{n\ge4}\lambda^{n-1}a_n^2
 \le\frac{\lambda^3a_4^2}{1-1/\lambda}
 =\frac{\lambda^4a_4^2}{c}.
\]
Retaining the first three terms of $T$ and the first two terms of $Q-\lambda T$ therefore gives
\begin{align*}
 cT&\le c a_1^2+c\lambda a_2^2+c\lambda^2a_3^2+\lambda^4a_4^2
 \le a_1^2E_4,\\
 Q-\lambda T&\ge a_1-\lambda a_1^2+\lambda a_1a_2-\lambda^2a_2^2=a_1B.
\end{align*}
For the second bound, every omitted term is nonnegative because $a_{n-1}\ge\lambda a_n$. The function $cT/[2(a_1B+cT)]$ increases with $T$. Thus
\begin{equation}\label{eq:mh-speed-upper}
 \frac{v}{1+v}\le\frac{a_1E_4}{2(B+a_1E_4)}<\frac12.
\end{equation}
No moment of the offspring number beyond its first moment has entered these bounds.

\subsubsection{The kernel for \texorpdfstring{$1/(\lambda+1)\le x<1$}{1/(lambda+1) <= x < 1}}
We prove the second inequality in~\eqref{eq:mh-two-kernels} on the slightly larger bias interval $7/5\le\lambda\le39/20$. At $x=1/(\lambda+1)$,
\[
 A_\lambda(x,y)=\frac{\lambda(\lambda+y)}{\lambda+1}.
\]
Substitution of $y=cU/(1-U)$ yields the exact identity
\begin{equation}\label{eq:mh-high-FG}
 \frac{2P_v(y)-P_c(y)}{(1+v)(\lambda+y)^2}
 =F(U)-2\frac{v}{1+v}G(U),
\end{equation}
where
\[
 F(U)=\frac{cU(2c+1-\lambda U)}{(\lambda-U)^2},\qquad
 G(U)=\frac{\lambda-(3\lambda-1)U+\lambda^2U^2}{(\lambda-U)^2}>0.
\]
The positivity follows alternatively from $G(U)=(y^2-y+\lambda)/(\lambda+y)^2$.
Set
\[
 f_1=\frac{c(2c+1)}{\lambda^2},\qquad
 f_2=\frac{c\{(4-\lambda)c+2-\lambda\}}{\lambda^3}>0,
 \qquad \gamma_0=\frac{3c}{\lambda^2},\quad
 \gamma_1=\frac{2\lambda+1}{(\lambda+1)^2}.
\]
For $1<\lambda\le2$ and $0\le U\le1/\lambda$,
\begin{equation}\label{eq:mh-high-FG-bounds}
 F(U)\ge f_1U+f_2U^2,\qquad
 G(U)\le\frac1\lambda-\min(\gamma_0,\gamma_1)U.
\end{equation}
Indeed the first remainder is
\[
 \frac{cU^3}{\lambda^3(\lambda-U)^2}
 \left[\lambda\{3(2c+1)-2\lambda^2\}
       -U\{2(2c+1)-\lambda^2\}\right].
\]
Its bracket decreases with $U$ and at $U=1/\lambda$ equals
$-2c(\lambda^3-2\lambda^2-\lambda+1)/\lambda>0$, since
$\lambda^3-2\lambda^2-\lambda+1=(\lambda-2)(\lambda^2-1)-1\le-1$.
For the second inequality, the function
\[
 \frac{G(0)-G(U)}U
 =\frac{3c-(\lambda^2-1/\lambda)U}{(\lambda-U)^2}
\]
has derivative whose numerator decreases with $U$. Its minimum is therefore at one of the two endpoints, where its values are $\gamma_0,\gamma_1$.

For $\ell=0,1$ consider the explicit cubic
\begin{equation}\label{eq:mh-high-sigma}
 \Sigma_\ell(s)=(f_1+f_2s)(B+a_1E_4)
                 -E_4(1/\lambda-\gamma_\ell a_1).
\end{equation}
Its coefficients in increasing powers of $s$ are
\begin{align*}
 C_0&=f_1(1-\lambda a_1)+c\{a_1(f_1+\gamma_\ell)-1/\lambda\},\\
 C_1&=f_2(1-a_1)+\lambda a_1f_1,\\
 C_2&=\lambda a_1(f_2-\lambda f_1)+A(a_1)\{a_1(f_1+\gamma_\ell)-1/\lambda\},\\
 C_3&=a_1f_2\{A(a_1)-\lambda^2\}.
\end{align*}
The following finite algebra lemma states the exact arithmetic used here.

\begin{lemma}[Coefficient inequalities for the kernel at $x=1/(\lambda+1)$]\label{lem:mh-high-arithmetic}
For $\lambda\ge7/5$ the five rational expressions
\[
 \lambda^2-A(\lambda^{-2}),\quad
 \lambda-f_1-\gamma_0,\quad\lambda-f_1-\gamma_1,\quad
 \lambda^3-\lambda^2+\lambda-2,\quad
 \lambda^5+\lambda^4-\lambda^3-3\lambda^2+1
\]
are strictly positive. For an independent coordinate $0\le u\le1$, define
\begin{align*}
 H_2&=\lambda^2+\lambda-u,\\
 H_3&=\lambda^3+\lambda^2+\lambda-(\lambda+1)u,\\
 H_4&=\lambda^4+\lambda^3+\lambda^2+\lambda-(\lambda^2+\lambda+1)u,\\
 \mathrm{Den}&=H_2^2H_3^2H_4^2,\\
 E_n&=c\mathrm{Den}+c\lambda^3H_3^2H_4^2
             +c\lambda^4H_2^2H_4^2+\lambda^6H_2^2H_3^2,\\
 V_n&=\lambda(\lambda-u)\mathrm{Den}
       +\lambda^2uH_2H_3^2H_4^2-u\lambda^4H_3^2H_4^2+uE_n,\\
 L_n&=\lambda H_2(2c+1)+\lambda\{(4-\lambda)c+2-\lambda\},\\
 N_0&=cL_nV_n-E_nH_2(\lambda^4-3c\lambda u),\\
 N_1&=cL_nV_n(\lambda+1)^2
       -E_nH_2\{\lambda^3(\lambda-u)(\lambda+1)^2+\lambda^5u\}.
\end{align*}
Then $N_0,N_1>0$. Every $H_i$ and $\mathrm{Den}$ is strictly positive on this domain. At $u=\lambda^2a_1$ the exact identities are
\[
 \Sigma_0(1/D_2)=\frac{N_0}{\lambda^5H_2\mathrm{Den}},\qquad
 \Sigma_1(1/D_2)=\frac{N_1}{\lambda^5H_2\mathrm{Den}(\lambda+1)^2}.
\]
\end{lemma}

Only finite coefficient arithmetic is used in this lemma. More explicitly, after writing $\tau=\lambda-7/5$, the numerators and denominators of the first five rational expressions have nonnegative coefficients and positive constants. Each $N_\ell$ has degree at most eight in $u$. Its nine degree-eight Bernstein coefficients are polynomials in $\tau$ with nonnegative coefficients and positive constants. The exact rational definitions above and the conversion
\begin{equation}\label{eq:mh-bernstein}
 [u^j]_{\mathrm{Bern},d}\left(\sum_{b=0}^dp_bu^b\right)
 =\sum_{b=0}^j p_b\frac{\binom jb}{\binom db}
\end{equation}
specify all eighteen coefficient checks in the supplement. Positivity follows on $0\le u\le1$ because the Bernstein basis is nonnegative and sums to one. Finally $H_i$ decrease in $u$, and their values at $u=1$ are respectively $\lambda^2+\lambda-1$, $\lambda^3+\lambda^2-1$, and $\lambda^4+\lambda^3-1$, all positive.

The ratios $D_2/D_3,D_2/D_4$ increase with $a_1$, so $A(a_1)\le A(\lambda^{-2})<\lambda^2$. Lemma~\ref{lem:mh-high-arithmetic} and
\[
 \lambda f_1-f_2=\frac{2c^2(\lambda^2+\lambda-1)}{\lambda^3}>0
\]
give $C_2<0,C_3<0$. Hence $\Sigma_\ell$ is strictly concave on $[0,1/D_2]$. At $s=0$ its value satisfies
\begin{align*}
 C_0(\gamma_0)&\ge\frac{c(\lambda^3-\lambda^2+\lambda-2)}{\lambda^4}>0,\\
 C_0(\gamma_1)&\ge\frac{c(\lambda^5+\lambda^4-\lambda^3-3\lambda^2+1)}{\lambda^4(\lambda+1)^2}>0.
\end{align*}
The other endpoint is positive by Lemma~\ref{lem:mh-high-arithmetic}. Concavity therefore proves~\eqref{eq:mh-high-sigma} positive for every admissible $s$.
Taking expectations in~\eqref{eq:mh-high-FG-bounds} and choosing $\gamma=\min(\gamma_0,\gamma_1)$ gives
\[
 \frac{v}{1+v}\le\frac{a_1E_4}{2(B+a_1E_4)}
 <\frac{a_1(f_1+f_2s)}{2(1/\lambda-\gamma a_1)}
 \le\frac{\E F(U)}{2\E G(U)}.
\]
The denominator is positive because $\gamma a_1<1/\lambda$. This proves the kernel inequality at $x=1/(\lambda+1)$.

To extend it to larger $x$, the polynomial
\[
 2P_v(y)-P_c(y)
 =(1-v)y^2+\{2c(1+v)+1+3v\}y-2\lambda v
\]
is strictly increasing for $y\ge0$, has negative value at zero and positive value $2\lambda$ at one. Let its unique zero be $y_0\in(0,1)$. For $x\ge1/(\lambda+1)$ the ratio
\[
 W_x(y)=\left[\frac{A_\lambda(1/(\lambda+1),y)}{A_\lambda(x,y)}\right]^2
\]
increases in $y$, since
\[
 \frac{\dd}{\dd y}\log W_x(y)
 =\frac{2\lambda\{x-1/(\lambda+1)\}}
 {A_\lambda(1/(\lambda+1),y)A_\lambda(x,y)}\ge0.
\]
Multiplying $[W_x(y)-W_x(y_0)][2P_v(y)-P_c(y)]\ge0$ by the positive reciprocal square at $1/(\lambda+1)$ and integrating proves the required strict inequality for every larger $x$.

\subsubsection{The kernel for \texorpdfstring{$0<x\le1/(\lambda+1)$}{0 < x <= 1/(lambda+1)}}
For $0<x\le1/(\lambda+1)$ put $\theta=\lambda x/(c+x)$, so $0<\theta\le1/\lambda$. The formula for $g$ gives
\[
 \frac{g(x)}{b(x)}=-\frac1c-\frac{\log\theta}{\log\lambda(1-\theta)}.
\]
We use the rational lower bound
\begin{equation}\label{eq:mh-Rrat}
 R=1+\frac{9\lambda(1-\lambda\theta)
                  \{28\lambda+17+(17\lambda+19)\theta\}}
 {c(\lambda+1)(2\lambda+7)(2\lambda+1)(2+7\theta)(2+\theta)}
 \le\frac{g(x)}{b(x)},\qquad R\ge1.
\end{equation}
To bound the logarithmic term, observe that for $0<\theta\le1/\lambda$,
\begin{align*}
 \frac{g(x)}{b(x)}-1
 &=\frac{1/\lambda-\theta}{\log\lambda}
   \int_0^1\frac{1-t}
   {\{\theta+(1-\theta)t\}\{1/\lambda+(1-1/\lambda)t\}}\,\dd t.
\end{align*}
The reciprocal product in this integrand is convex in $t$. Apply Jensen's inequality separately on $[0,1/2]$ and $[1/2,1]$ with weight $(1-t)\dd t$. The masses are $3/8,1/8$ and the means $2/9,2/3$, respectively. The integral is therefore at least
\[
 \lambda\left\{
 \frac{243}{8(2+7\theta)(2\lambda+7)}
 +\frac9{8(2+\theta)(2\lambda+1)}\right\}.
\]
The trapezoid bound for the convex function $1/t$ gives
$\log\lambda\le c(\lambda+1)/(2\lambda)$. Substitution and addition of the two fractions give exactly~\eqref{eq:mh-Rrat}. Equality holds at $\theta=1/\lambda$.

For fixed $R\ge1$, the signed kernel divided by $1+v$ decreases as $v/(1+v)$ increases: its numerator has derivative
\[
 -2R(y^2+c)-2(1-y)<0.
\]
Since $\Psi>0$, we may first replace $g(x)/b(x)$ by $R$, and then replace $v/(1+v)$ by the upper bound in~\eqref{eq:mh-speed-upper}. Both operations give lower bounds for the expected signed kernel.

We record the exact quadratic reduction before proving its sign. In this calculation only, write $r=v/(1+v)$ and denote the three polynomial coefficients by
\[
 C=2r(1+cR),\quad
 L=c+2(\lambda+1)r+2Rc(c+2r),\quad K=\lambda(c+C).
\]
Substitution of $y=cU/(1-U)$ gives
\begin{equation}\label{eq:mh-low-transform}
 \frac{2RP_v(y)-P_c(y)}{(1+v)A_\lambda(x,y)^2}
 =\frac{-C+LU-KU^2}{(c+x)^2(1-\theta U)^2}.
\end{equation}
Put $D_1=L-2\theta C$ and $E_1=-K+2\theta L-3\theta^2C$. The expression
\[
 D_1=c+2r\{\lambda+1-2\theta\}+2Rc\{c+2r(1-\theta)\}
\]
is strictly positive. Polynomial division yields
\[
 \frac{-C+LU-KU^2}{(1-\theta U)^2}
 =-C+D_1U+U^2j(U),\qquad
 j(U)=\frac{E_1-\theta^2D_1U}{(1-\theta U)^2}.
\]
Its derivative is
\[
 j'(U)=\frac{\theta\{2E_1-\theta D_1-\theta^2D_1U\}}{(1-\theta U)^3}.
\]
The numerator decreases with $U$, so $j$ attains its minimum at $0$ or $1/\lambda$. Consequently the expectation in~\eqref{eq:mh-low-transform} is strictly positive if both
\begin{equation}\label{eq:mh-B01}
 B_0=-C+D_1a_1+E_1a_1s>0,\qquad
 B_1=(-C+D_1a_1)(1-\theta/\lambda)^2
             +(E_1-\theta^2D_1/\lambda)a_1s>0.
\end{equation}
We now verify these two inequalities at $r=a_1E_4/[2(B+a_1E_4)]$ and $R$ from~\eqref{eq:mh-Rrat}, throughout $707/500\le\lambda\le39/20$. The ratio $r$ and the coefficients $C,L,K$ in this reduction will not denote the response function or parameter term outside this calculation.

\subsubsection{Concavity for \texorpdfstring{$a_1\le s\le1/D_2$}{a1 <= s <= 1/D2}}
The dependence on $s$ in~\eqref{eq:mh-B01} can be handled before the endpoint arithmetic. Set
\begin{align*}
 T&=1+cR,\qquad M=1+2cR,\qquad m_b=cM,\\
 e_b&=-\lambda c+2\theta m_b,\\
 l_d&=2(\lambda+1)+4cR-4\theta T,\\
 e_r&=-2\lambda T+4\theta(\lambda+1)+8\theta cR-6\theta^2T.
\end{align*}
Here $T$ is a rational coefficient, rather than the moment sum in~\eqref{eq:mh-moment-speed}. For $i=0$ take
\[
 (w,m_0,m_1,l_0,l_1)=(1,m_b,e_b,l_d,e_r),
\]
and for $i=1$ take
\[
 w=(1-\theta/\lambda)^2,\quad
 (m_0,m_1,l_0,l_1)
 =(m_bw,e_b-\theta^2m_b/\lambda,l_dw,e_r-\theta^2l_d/\lambda).
\]
Define, for either choice,
\begin{equation}\label{eq:mh-hybrid-sigma}
 \Sigma_i=(B+a_1E_4)(m_0+m_1s)
               +\frac{a_1E_4}{2}(l_0+l_1s)-E_4Tw.
\end{equation}
Direct substitution into~\eqref{eq:mh-B01} gives
\begin{equation}\label{eq:mh-hybrid-normalization}
 B_i=\frac{a_1\Sigma_i}{B+a_1E_4}.
\end{equation}
All denominators are positive, including $w$, because $\theta/\lambda\le\lambda^{-2}<1$.

We prove $\Sigma_i''(s)<0$ for $a_1\le s\le1/D_2$ and $7/5\le\lambda\le39/20$. With
\[
 K_0=a_1(m_0+l_0/2)-Tw,\qquad K_1=a_1(m_1+l_1/2),
\]
the four cubic coefficients of~\eqref{eq:mh-hybrid-sigma} are
\begin{align*}
 C_0&=m_0(1-\lambda a_1)+cK_0,\\
 C_1&=m_1(1-\lambda a_1)+\lambda a_1m_0+cK_1,\\
 C_2&=\lambda a_1m_1-\lambda^2a_1m_0+A(a_1)K_0,\\
 C_3&=-\lambda^2a_1m_1+A(a_1)K_1.
\end{align*}
Thus $\Sigma_i''=2(C_2+3C_3s)$, and its sign is affine in $A(a_1)$. It suffices to prove negativity at $A=0$ and $A=\lambda^2$.

We first bound $R$ from~\eqref{eq:mh-Rrat}. The derivative of
$\{28\lambda+17+(17\lambda+19)\theta\}/\{(2+7\theta)(2+\theta)\}$ has numerator
\[
 -4(95\lambda+49)-14(28\lambda+17)\theta
                    -7(17\lambda+19)\theta^2<0.
\]
Moreover the value of $R-1$ at $\theta=0$ decreases with $\lambda$: the logarithmic derivative of its numerator is less than $2/\lambda$, while the two denominator factors $c,\lambda+1$ alone have logarithmic derivative $2\lambda/(\lambda^2-1)>2/\lambda$. Its value at $7/5$ gives
\begin{equation}\label{eq:mh-price-bounds}
 1\le R\le1+5(1-\lambda\theta),\qquad
 \lambda\le T,\quad 2\lambda-1\le M\le1+12c-10c\lambda\theta<13.
\end{equation}
Also $D_2\ge59/35>3/2$, so $s<2/3$. Put $\zeta=c/T\in[0,c/\lambda]$.

\begin{lemma}[Finite algebra used in the curvature proof]\label{lem:mh-curvature-arithmetic}
For $7/5\le\lambda\le39/20$, all four values
\[
 1-\lambda^{-2}\left[2-3\lambda s+
 \frac{\{(12+6\zeta)s-2(2-\lambda+\zeta)/\lambda\}^2}{36s}\right]
\]
are strictly positive when $\zeta\in\{0,c/\lambda\}$ and
$s\in\{\lambda^{-2},(\lambda+1-1/\lambda)^{-1}\}$.
The polynomial
\[
 F(\lambda,t)=\lambda^4-(1+12c-10ct)(4\lambda^2t-2t^2-\lambda^4)
\]
is positive for $\lambda\ge7/5$ and $0\le t\le1$.
Finally set
\[
 h_0=4-2\lambda-2/\lambda+2\zeta+
        \frac{2(c-\zeta)}{(1-\lambda^{-2})^2},\qquad
 g=2+\lambda(\lambda-2-\zeta)
       -\frac{\lambda(c-\zeta)}{(1-\lambda^{-2})^2}.
\]
The following five rational quantities are strictly positive:
\begin{gather*}
 6\frac{\lambda^2}{\lambda^2+\lambda-1}-2
 -\frac{2c}{\lambda(1-\lambda^{-2})^3}
       \left(3\frac{\lambda^2}{\lambda^2+\lambda-1}-2\right),\\
 -g\big|_{\zeta=0},\quad
 (-12g-h_0^2)\big|_{\zeta=0},\quad
 (1/10-g)\big|_{\zeta=c/\lambda},\quad
 (\lambda^2-1/5-h_0)\big|_{\zeta=c/\lambda}.
\end{gather*}
All displayed denominators are positive for $\lambda>1$. These assertions are finite coefficient inequalities: after $\lambda=7/5+\tau$, the nine rational numerators and denominators have nonnegative coefficients and positive constants; the six degree-five Bernstein coefficients of $F$ in $t$ have the same property as polynomials in $\tau$.
\end{lemma}

We give the remaining continuous curvature argument. For $i=0,A=0$, divide $C_2+3C_3s$ by the positive number $a_1c\lambda$ and put $x=1-3\lambda s$. The result is
\[
 x(2\theta M-\lambda)-\lambda M.
\]
When $0\le x\le1$, its maximum is at $x=0$ or $1$. The first value is negative; the second is at most $M(2/\lambda-\lambda)-\lambda<0$, since $2/\lambda-\lambda\le1/35$ and $M<13$. When $x<0$, omit the nonpositive term $2\theta Mx$ and use $M\ge2\lambda-1$ to obtain the strict upper bound $\lambda^2(3s-2)<0$.

For $i=0,A=\lambda^2$, division by $\lambda^2T$ gives
\[
 a_1(h_0+3sg)-1,\qquad
 h_0=2-\frac{2\theta}{\lambda}(2-\lambda+\zeta),\quad
 g=-\lambda+4\theta+2\zeta\theta-3\theta^2.
\]
Completing the square in $\theta$ bounds this by
\begin{equation}\label{eq:mh-curvature-square}
 a_1\left[2-3\lambda s+
 \frac{\{(12+6\zeta)s-2(2-\lambda+\zeta)/\lambda\}^2}{36s}\right]-1.
\end{equation}
The bracket is convex first in $\zeta$, and then, with $\zeta$ fixed, in $s$. We need only $\zeta=0,c/\lambda$ and $s=a_1,1/D_2$. To see that these endpoint bounds increase with $a_1$, put
\[
 p=12+6\zeta,\quad q=2(2-\lambda+\zeta)/\lambda,\quad
 L=p^2/36-3\lambda,\quad b=2-pq/18.
\]
The expressions in~\eqref{eq:mh-curvature-square}, before subtracting one, are
$ba_1+La_1^2+q^2/36$ and
$a_1[b+L/D_2+q^2D_2/36]$. At $\zeta=0$ their derivatives are bounded below, respectively, by
\[
 \frac{10\lambda^2-26\lambda+24}{3\lambda^2}>0,
 \qquad\frac{\lambda^2-5\lambda+12}{3\lambda^2}>0.
\]
Both numerators have negative discriminants. At $\zeta=c/\lambda$, one has $0<q\le62/49$, $p<18$, $b>2-q>0$, and
\[
 L=9-3\lambda-6/\lambda+1/\lambda^2>1/4.
\]
The last function decreases on $[7/5,2]$ and equals $1/4$ at two. Both derivatives are again positive. We may therefore replace $a_1$ by $\lambda^{-2}$; the first four inequalities of Lemma~\ref{lem:mh-curvature-arithmetic} prove strict negativity.

For $i=1,A=0$, the corresponding expression divided by $a_1c\lambda$ is
\[
 -\lambda x+M\{(x+1)J_\theta-\lambda\},\qquad
 x=1-3\lambda s,\quad J_\theta=2\theta-\theta^2/\lambda.
\]
For $0\le x\le1$, check the two endpoints. At zero it is $-\lambda Mw<0$. At one, if $2J_\theta-\lambda>0$, use~\eqref{eq:mh-price-bounds}; multiplication of the required strict inequality by $\lambda^3$ gives precisely $F(\lambda,\lambda\theta)>0$ from Lemma~\ref{lem:mh-curvature-arithmetic}. If the bracket is nonpositive there is nothing to prove. For $x\le-1$, use $J_\theta\ge0$ and $M\ge2\lambda-1$ to obtain $\lambda^2(3s-2)<0$. For $-1\le x\le0$, use $J_\theta\le2/\lambda-1/\lambda^3<\lambda$ and $M\ge2\lambda-1$. The resulting affine function of $x$ is negative at both endpoints: its values there are $-2\lambda c$ and $(2\lambda-1)(2/\lambda-1/\lambda^3-\lambda)$.

For $i=1,A=\lambda^2$, divide by $\lambda^2Tw$. The result is $a_1(h_0+3sg)-1$, where now
\[
 h_0=4-2\lambda-2\theta+2\zeta+2(c-\zeta)/w,
 \qquad g=2\lambda\theta+\lambda(\lambda-2-\zeta)-\lambda(c-\zeta)/w.
\]
It is affine in $\zeta$, so take $\zeta=0,c/\lambda$. Its remaining dependence is
\[
 h_0+3sg=4-2\lambda+2\zeta+3\lambda s(\lambda-2-\zeta)
 +(6\lambda s-2)\theta+(c-\zeta)(2-3\lambda s)/w.
\]
If $\lambda s\le2/3$, this is convex in $\theta$ and its maximum is at $0$ or $1/\lambda$. If $\lambda s>2/3$, it is concave. At $\theta=1/\lambda$ its derivative is bounded below by
\[
 6\lambda s-2-\frac{2c}{\lambda(1-\lambda^{-2})^3}(3\lambda s-2)>0.
\]
Indeed this expression is affine in $\lambda s$, equals two at $2/3$, and is positive at the other endpoint $\lambda^2/(\lambda^2+\lambda-1)$ by Lemma~\ref{lem:mh-curvature-arithmetic}. Thus the maximum is again at an endpoint in $\theta$.
At $\theta=0$, $h_0=2,g=-\lambda$ and
\[
 a_1(h_0+3sg)-1\le2a_1-3\lambda a_1^2-1
 \le1/(3\lambda)-1<0.
\]
At $\theta=1/\lambda,\zeta=0$, the finite lemma gives $g<0$ and $-12g>h_0^2$; since $s\ge a_1$, completing the square in $a_1$ gives an upper bound $h_0^2/(-12g)-1<0$.
At $\theta=1/\lambda,\zeta=c/\lambda$, the same lemma gives $g<1/10$ and $h_0+1/5<\lambda^2$. Since $s<2/3$, it follows that $h_0+3sg<\lambda^2$, and $a_1\le\lambda^{-2}$ again gives strict negativity. Affinity in $A$ proves
\begin{equation}\label{eq:mh-s-concavity}
 \Sigma_0''(s)<0,\qquad\Sigma_1''(s)<0
 \quad(a_1\le s\le1/D_2).
\end{equation}

\subsubsection{Positivity at \texorpdfstring{$s=a_1$ and $s=1/D_2$}{s = a1 and s = 1/D2}}
We next prove the four endpoint inequalities required by~\eqref{eq:mh-s-concavity}. At $s=a_1$, taking $R=1$ gives a lower bound after dividing $\Sigma_i$ by the positive factor $Tw$. At $s=1/D_2$ we use the dependence of $A(a_1)$ on the first moment to narrow the interval $0\le A\le\lambda^2$ used in the concavity argument.

Set $J_\theta=2\theta-\theta^2/\lambda$ and $\Phi_i=\Sigma_i/(Tw)$. Direct expansion of~\eqref{eq:mh-hybrid-sigma} gives
\begin{align*}
 \Phi_i&=BP_i+E_4(a_1Q_i-1),\\
 P_0&=2c(1+2\theta s)-\zeta\{1+(\lambda+2\theta)s\},\\
 Q_0&=2\lambda-2\theta+s\{-\lambda(1+\zeta)+4\lambda\theta-3\theta^2\},\\
 P_1&=2c-\zeta+\frac{s}{w}\{2cJ_\theta-\zeta(\lambda+J_\theta)\},\\
 Q_1&=2\lambda-2\theta+
 \frac{s}{w}\{-\lambda(1+\zeta)+4\lambda\theta-5\theta^2+2\theta^3/\lambda\}.
\end{align*}
In these expressions $B,E_4$ are independent of $\zeta$. Thus
\begin{align*}
 \partial_\zeta\Phi_0&=-B\{1+(\lambda+2\theta)s\}-\lambda a_1sE_4<0,\\
 \partial_\zeta\Phi_1&=-B\{1+s(\lambda+J_\theta)/w\}
                      -\lambda a_1sE_4/w<0.
\end{align*}
Since $R\ge1$ gives $\zeta\le c/\lambda$, the value of $\Phi_i$ at $R=1$ is a lower bound for $\Phi_i$ at the value of $R$ in~\eqref{eq:mh-Rrat}. The positive factor $Tw$ then transfers strict positivity to $\Sigma_i$.

\begin{lemma}[Coefficient inequalities at $s=a_1$]\label{lem:mh-lower-arithmetic}
In~\eqref{eq:mh-hybrid-sigma}, replace $A(a_1)$ temporarily by $b\in\{0,\lambda^2\}$, and set $R=1$, $a_1=s=u/\lambda^2$, $\theta=t/\lambda$. The explicit products
\begin{equation}\label{eq:mh-lower-polynomials}
 F_i^b(\lambda,u,t)
 =\lambda^{8+2i}\left.\frac{\Sigma_i}{T}\right|_{R=1,\ a_1=s=u/\lambda^2,\ A=b,\ \theta=t/\lambda}
 \quad(i=0,1)
\end{equation}
are polynomials and are strictly positive for $\lambda\ge7/5$, $0\le u,t\le1$. Their respective degrees in $(u,t)$ are at most $(4,2)$ and $(4,3)$. Every coefficient in the corresponding product Bernstein bases is a polynomial in $\lambda-7/5$ with nonnegative coefficients and positive constant.
\end{lemma}

Equation~\eqref{eq:mh-lower-polynomials} and the definitions of $B,E_4,T,m_0,m_1,l_0,l_1$ specify the four polynomials. Formula~\eqref{eq:mh-bernstein} in each variable reduces the lemma to $2\cdot15+2\cdot20=70$ finite coefficient checks. All factors divided out are positive powers of $\lambda$ and $T=\lambda$. Because $\Sigma_i$ is affine in $A$, the lemma covers $0<A(a_1)<\lambda^2$. The comparison with $R=1$ therefore gives
\begin{equation}\label{eq:mh-lower-positive}
 \Sigma_0(a_1)>0,\qquad\Sigma_1(a_1)>0.
\end{equation}

At $s=1/D_2$ we first bound $A(a_1)$. Differentiation gives
\begin{align*}
 (D_2/D_3)'&=\lambda^2/D_3^2,&
 (D_2/D_3)''&=2\lambda^3(\lambda+1)/D_3^3,\\
 (D_2/D_4)'&=\lambda^2(\lambda+1)/D_4^2,&
 (D_2/D_4)''&=2\lambda^3(\lambda+1)(\lambda^2+\lambda+1)/D_4^3.
\end{align*}
Thus $A'>0,A''>0$ on $0\le a_1\le\lambda^{-2}$. The endpoint values are
\begin{align*}
 A_0=A(0)&=c\lambda+
 \frac{c\lambda^2(\lambda+1)^2}{(\lambda^2+\lambda+1)^2}
 +\frac{\lambda^4}{(\lambda^2+1)^2},\\
 A_{\max}=A(\lambda^{-2})&=c\lambda+
 \frac{c\lambda^2(\lambda^2+\lambda-1)^2}{(\lambda^3+\lambda^2-1)^2}
 +\frac{\lambda^4(\lambda^2+\lambda-1)^2}{(\lambda^4+\lambda^3-1)^2}.
\end{align*}
\begin{lemma}[Finite chord bound]\label{lem:mh-chord-arithmetic}
The rational difference $1/3-A_{\max}+A_0$ is positive for $\lambda\ge7/5$. Its denominator is
\[
 3(\lambda^2+1)^2(\lambda^2+\lambda+1)^2
   (\lambda^3+\lambda^2-1)^2(\lambda^4+\lambda^3-1)^2>0,
\]
and its numerator is the explicit polynomial
\begin{align*}
 &\lambda^{22}+6\lambda^{21}+13\lambda^{20}+16\lambda^{19}+5\lambda^{18}
 -35\lambda^{17}-67\lambda^{16}-59\lambda^{15}-13\lambda^{14}\\
 &\quad+45\lambda^{13}+37\lambda^{12}+11\lambda^{11}-22\lambda^{10}-9\lambda^9
 +22\lambda^8+18\lambda^7+15\lambda^6\\
 &\quad-19\lambda^5-8\lambda^4-8\lambda^3+3\lambda^2+2\lambda+1.
\end{align*}
All its coefficients after translation by $7/5$ are positive.
\end{lemma}
The displayed numerator makes Lemma~\ref{lem:mh-chord-arithmetic} a finite, degree-twenty-two coefficient check. Increasing convexity of $A$ and that lemma give the bound
\begin{equation}\label{eq:mh-A-chord}
 A_0\le A(a_1)\le A_0+\frac{\lambda^2a_1}{3}.
\end{equation}

To bound $\Sigma_i(1/D_2)$, we use independent polynomial variables $u,t\in[0,1]$. We will evaluate the polynomials at $u=\lambda^2a_1$, $t=\lambda\theta$. Put
\begin{align*}
 H&=\lambda^2+\lambda-u,\\
 D_r&=(\lambda+1)(2\lambda+7)(2\lambda+1)(2\lambda+7t)(2\lambda+t),\\
 Z_r&=cD_r+9\lambda^2(1-t)\{28\lambda^2+17\lambda+(17\lambda+19)t\},\\
 T_n&=D_r+Z_r,\qquad M_0=c(D_r+2Z_r),\\
 N_m&=-\lambda^2cD_r+2tM_0,\\
 N_d&=2\lambda(\lambda+1)D_r+4\lambda Z_r-4tT_n,\\
 N_e&=-2\lambda^3T_n+4t\lambda(\lambda+1)D_r+8t\lambda Z_r-6t^2T_n,\\
 \delta_0&=(\lambda^2+\lambda+1)^2(\lambda^2+1)^2,\\
 \alpha_0&=c\lambda\delta_0+c\lambda^2(\lambda+1)^2(\lambda^2+1)^2
                         +\lambda^4(\lambda^2+\lambda+1)^2.
\end{align*}
For the lower and upper bounds in~\eqref{eq:mh-A-chord}, respectively, use
\[
 (\alpha,\delta)=(\alpha_0,\delta_0)
 \quad\hbox{or}\quad
 (\alpha,\delta)=(3\alpha_0+u\delta_0,3\delta_0).
\]
In either case $A=\alpha/\delta$ and $\delta$ is independent of $u$. Define
\begin{align*}
 E_n&=c\delta H^2+\alpha\lambda^2,\\
 V_n&=\delta\{\lambda(\lambda-u)H^2+\lambda^2uH-\lambda^4u\}+uE_n,\\
 N_0&=\lambda HV_nM_0+\lambda V_nN_m
       +\tfrac12uE_n(HN_d+N_e)-\lambda^3HE_nT_n,\\
 N_1&=\{\lambda HV_nM_0+\tfrac12uE_nHN_d-\lambda^3HE_nT_n\}(\lambda^2-t)^2\\
    &\qquad+\lambda^3V_n(\lambda^2N_m-t^2M_0)
       +\tfrac12\lambda^2uE_n(\lambda^2N_e-t^2N_d).
\end{align*}
These products satisfy
\begin{equation}\label{eq:mh-upper-numerators}
 \Sigma_0(1/D_2)=\frac{N_0}{\lambda^3\delta H^3D_r},\qquad
 \Sigma_1(1/D_2)=\frac{N_1}{\lambda^7\delta H^3D_r}.
\end{equation}
For example the coefficient substitutions are
\[
 \begin{gathered}
 cR=Z_r/D_r,\quad T=T_n/D_r,\quad m_b=M_0/D_r,\\
 e_b=N_m/(\lambda D_r),\quad
 l_d=N_d/(\lambda D_r),\quad e_r=N_e/(\lambda^2D_r).
 \end{gathered}
\]
and $E_4=E_n/(\delta H^2)$, $B+a_1E_4=V_n/(\lambda^2\delta H^2)$.
They verify~\eqref{eq:mh-upper-numerators} term by term. Every displayed denominator is positive: $H\ge\lambda^2+\lambda-1>0$, $D_r>0$, and $\delta>0$.

\begin{lemma}[Finite algebra for the two integrations]\label{lem:mh-integral-arithmetic}
For each of the two pairs $(\alpha,\delta)$ above, set
\[
 K_i=H\partial_uN_i+4N_i\qquad(i=0,1).
\]
Then $K_i>0$ for $707/500\le\lambda\le39/20$ and $0\le u,t\le1$. The $u^4$ terms cancel in this definition, and the degrees in $(u,t)$ are at most $(3,4)$ and $(3,5)$, respectively. Their product Bernstein coefficients, viewed as polynomials in $\lambda-7/5$, are nonnegative with positive constants, with one exception: for the upper pair and $i=0$, the coefficient with Bernstein indices $(0,4)$ has this property after translation by $707/500$ instead.

For the lower pair define the exact polynomial quotients
\[
 W_0(\lambda,t)=\frac{N_0(\lambda,0,t)}{\lambda^5(\lambda+1)},\qquad
 W_1(\lambda,t)=\frac{N_1(\lambda,0,t)}{\lambda^6(\lambda+1)}.
\]
Their derivatives in $\lambda$, with $t$ fixed, are strictly positive for $\lambda\ge7/5$, $0\le t\le1$. Their degree-three and degree-four Bernstein coefficients in $t$, respectively, have nonnegative coefficients and positive constants after translation by $7/5$. At $\lambda=707/500$, all four degree-three Bernstein coefficients of $W_0$ are positive. Define $b_j$ to be the degree-five Bernstein coefficients of the explicitly defined polynomial $W_1(707/500,t)$, so that
\begin{equation}\label{eq:mh-boundary-P}
 W_1(707/500,t)=\sum_{j=0}^5\binom5j b_jt^j(1-t)^{5-j}.
\end{equation}
These exact rational coefficients satisfy
\[
 72b_0+25b_1>0,\qquad10b_2-40b_0>0,\qquad10b_3-187b_0>0,
 \qquad b_0,b_4,b_5>0.
\]
\end{lemma}

Lemma~\ref{lem:mh-integral-arithmetic} reduces to coefficient calculations for the displayed products. Formula~\eqref{eq:mh-bernstein} gives the $88$ Bernstein coefficients for the four choices of $K_i$, nine for the two derivatives, and four for $W_0(707/500,t)$. The monotonicity in $u$ applies to $N_i/H^4$.

For completeness, the finite comparisons in Lemma~\ref{lem:mh-integral-arithmetic} prove the sign of the polynomial in~\eqref{eq:mh-boundary-P} by an estimate on the whole half-line. The exact rational values of the $b_j$ belong to the coefficient data; their defining polynomial and the inequalities used here are explicit.
For $0\le t<1$ set $y=t/(1-t)$. Multiplication by $(1+y)^5/b_0$ bounds the polynomial from below by
\[
 1-\frac{72}{5}y+40y^2+187y^3.
\]
For $0\le y\le1/12$, omit the cubic term and use that the resulting quadratic decreases on this interval; its value at $1/12$ is $7/90>0$. For $y\ge1/12$, use $187y^3\ge(187/12)y^2$ to obtain
\[
 1-\frac{72}{5}y+\frac{667}{12}y^2
 =\frac{667}{12}\left(y-\frac{432}{3335}\right)^2
     +\frac{1123}{16675}>0.
\]
At $t=1$ the value is $b_5>0$. Thus $W_1(707/500,t)>0$ on the entire closed interval.

We can now perform the two integrations. First, for $\lambda\ge707/500$ and fixed $t$,
\[
 W_i(\lambda,t)=W_i(707/500,t)+
       \int_{707/500}^{\lambda}\partial_\ell W_i(\ell,t)\,\dd\ell>0.
\]
This proves $N_i(\lambda,0,t)>0$ for the lower pair. For the upper pair both $N_i$ and $\delta$ at $u=0$ are exactly three times their lower-pair values, so the same positivity holds.
Second,
\[
 \frac{\partial}{\partial u}\frac{N_i}{H^4}=\frac{K_i}{H^5}>0.
\]
Since $H>0$ for the entire integration interval,
\[
 \frac{N_i(\lambda,u,t)}{H(\lambda,u)^4}
 =\frac{N_i(\lambda,0,t)}{H(\lambda,0)^4}
  +\int_0^u\frac{K_i(\lambda,z,t)}{H(\lambda,z)^5}\,\dd z>0.
\]
Both integrations use rational functions of the independent variables. By~\eqref{eq:mh-upper-numerators}, $\Sigma_i(1/D_2)>0$ at both bounds for $A$ in~\eqref{eq:mh-A-chord}. Since $\Sigma_i$ is affine in $A$, the inequality holds at $A=A(a_1)$ as well. Together with~\eqref{eq:mh-lower-positive} and strict concavity~\eqref{eq:mh-s-concavity}, this proves $\Sigma_i(s)>0$ for $a_1\le s\le1/D_2$. Equation~\eqref{eq:mh-hybrid-normalization} now proves~\eqref{eq:mh-B01}.

Returning to~\eqref{eq:mh-low-transform}, the comparisons in $R$ and $v/(1+v)$ prove the first inequality in~\eqref{eq:mh-two-kernels}. The kernel estimate for $x\ge1/(\lambda+1)$ gives the second. At $x=1/(\lambda+1)$ both inequalities are strict even though $g=b$ and $R=1$. Since $0<X<1$ almost surely and $b(X)>0$, equations~\eqref{eq:mh-max-source}--\eqref{eq:mh-source-split} give
\begin{equation}\label{eq:mh-middle-conclusion}
 D>0\qquad\left(707/500\le\lambda\le39/20,\quad\lambda<m\right).
\end{equation}
\subsubsection{An integral formula for \texorpdfstring{$39/20\le\lambda<(3+\sqrt5)/2$}{39/20 <= lambda < (3+sqrt(5))/2}}
For $39/20\le\lambda<(3+\sqrt5)/2$ we retain the full series for $\mathcal R$ and express $D$ as a signed integral using the response measure $\eta$.
For every continuous $f$,
\[
 \eta((I-\cL)f)=\E[b(X)f(X)].
\]
Apply this with $f=\cJ P_c$ and use~\eqref{eq:mh-C-kernel}. It gives
\begin{equation}\label{eq:mh-direct-signed}
 D=\eta(2\cJ P_v-\cJ P_c+\cL\cJ P_c).
\end{equation}
Conditional on $\nu$, the escape probabilities $X_i$ at the children of the root are independent with law $\pi_\lambda$. Write $S=\sum_{i=1}^\nu X_i$. Cancellation of denominators and exchangeability give
\[
 (\cL\cJ P_c)(x)
 =\frac1\lambda\E\frac{\sum_{i=1}^\nu P_c(X_i)}{(c+x+S)^2}.
\]
The recursion also gives
\[
 (\cJ P)(x)=\frac1{\lambda^2}
 \E\frac{(\lambda+S)^2P(F_\lambda(S))}{(c+x+S)^2}
\]
for either polynomial $P=P_v,P_c$. Expanding both sides gives
\begin{align}\label{eq:mh-root-polynomial}
 &(\lambda+S)^2\{2P_v(F_\lambda(S))-P_c(F_\lambda(S))\}
       +\lambda\sum_{i=1}^\nu P_c(X_i)\notag\\
 &\quad=\lambda\left[
 2S^2+2\{c-(\lambda+1)v\}S
 +(1+v)\sum_{i=1}^\nu X_i^2+2v(\nu-\lambda^2)\right].
\end{align}
Divide the bracket in~\eqref{eq:mh-root-polynomial} by $\lambda(c+x+S)^2$ and take expectation over the root and its children. Applying $\eta$ to the resulting function of $x$ gives $D$ by~\eqref{eq:mh-direct-signed}. The quotient is absolutely integrable uniformly in $x\in[0,1]$ because $S^2/(c+S)^2\le1$, $\sum_iX_i^2\le S$ and $\E\nu<\infty$.

Let $f(z)=\E z^\nu$ be the offspring generating function, and set
\[
 \phi(t)=\E e^{-tX},\qquad
 w(t)=\eta(e^{-t\cdot}),\qquad
 \xi(\dd t)=t e^{-ct}f(\phi(t))\,\dd t.
\]
The function $w$ is positive, bounded and decreasing, because $\eta$ is a finite positive measure with $\eta\ge b\pi_\lambda$. The measure $\xi$ is finite, positive and nonzero. For $t>0$ define the exponentially weighted means
\[
 u_t=\frac{\E[Xe^{-tX}]}{\phi(t)},\quad
 z_t=\frac{\E[X^2e^{-tX}]}{\phi(t)},\quad
 n_t=\frac{\phi(t)f'(\phi(t))}{f(\phi(t))},\quad
 j_t=\frac{\phi(t)^2f''(\phi(t))}{f(\phi(t))}.
\]
They are finite even when the offspring second moment is infinite, since $0<\phi(t)<1$. Define
\begin{equation}\label{eq:mh-B-source}
 B(t)=(3+v)n_tz_t+2j_tu_t^2
       +2\{c-(\lambda+1)v\}n_tu_t+2v(n_t-\lambda^2).
\end{equation}
The identities
\begin{align*}
 \E e^{-tS}&=f(\phi),&
 \E[Se^{-tS}]&=f'(\phi)\E[Xe^{-tX}],\\
 \E[S^2e^{-tS}]&=f'(\phi)\E[X^2e^{-tX}]
                         +f''(\phi)(\E[Xe^{-tX}])^2,&
 \E[\nu e^{-tS}]&=\phi f'(\phi)
\end{align*}
and $\int_0^\infty te^{-(c+x+S)t}\dd t=(c+x+S)^{-2}$ now give
\begin{equation}\label{eq:mh-signed-integral}
 \begin{gathered}
 D=\frac1\lambda\int_0^\infty w(t)B(t)\,\xi(\dd t),\qquad
 I:=\int_0^\infty B(t)\,\xi(\dd t),\\
 I=\E\frac{2S^2+2\{c-(\lambda+1)v\}S
 +(1+v)\sum_iX_i^2+2v(\nu-\lambda^2)}{(c+S)^2}.
 \end{gathered}
\end{equation}
The preceding absolute bounds justify applying Tonelli's theorem separately to the nonnegative terms before subtracting. They also give $\int|B|\dd\xi<\infty$.

If $c-(\lambda+1)v\ge0$, then $B(t)$ decreases strictly. The weighted moments $u_t,z_t$ of $X$ decrease by the covariance inequality. The weighted offspring law assigns mass proportional to $\PP(\nu=n)\phi(t)^n$. Its means of $n$ and $n(n-1)$ decrease as $\phi(t)$ decreases. The mean $n_t$ decreases strictly, because the offspring law has positive mass at one and at some larger value. Thus each term in~\eqref{eq:mh-B-source} is nonincreasing, and the term $2vn_t$ gives strict decrease. The covariance inequality for the two decreasing functions $w,B$ yields
\begin{equation}\label{eq:mh-source-Chebyshev}
 D\ge\frac1\lambda\frac{\int w\dd\xi}{\int\dd\xi}\,I.
\end{equation}
For an unbounded $B$, first apply the inequality to $\max(-k,\min(B,k))$ and then use $|\max(-k,\min(B,k))|\le|B|$ and boundedness of $w$ to pass to the limit. Hence $I>0$ and $c-(\lambda+1)v\ge0$ imply $D>0$.

\subsubsection{A speed bound for \texorpdfstring{$39/20\le\lambda<(3+\sqrt5)/2$}{39/20 <= lambda < (3+sqrt(5))/2}}
Assume $39/20\le\lambda<(3+\sqrt5)/2$ and~\eqref{eq:mh-hard}. From here on, write $s=h/c$. Then
\begin{equation}\label{eq:mh-coarse-domain}
 c\ge19/20,\quad
 0<s<\frac1{c(c+2)}\le\frac{400}{1121}<\frac9{25},\quad
 0<\rho_2<1,\quad d=1-cs>0,
\end{equation}
and the first two moments of $U$ are
\[
 a_1=\frac{s}{1+s},\qquad a_2=\frac{\rho_2s}{c+2+s}.
\]
Retain $T,Q$ from~\eqref{eq:mh-moment-speed}. The full moment ratio comparison gives
\begin{align}\label{eq:mh-coarse-TQ}
 Q&\ge\lambda T+D_0,\notag\\
 D_0&=a_1(1-\lambda a_1)+\lambda a_2(a_1-\lambda a_2)>0.
\end{align}
The two factors are strictly positive: $a_1<\lambda^{-2}<\lambda^{-1}$, and $\beta_1>\lambda\beta_2$, together with $\rho_2<1$, gives $a_1>\lambda a_2$. For every $n\ge2$, $a_n\le\rho_2\beta_n\le\rho_2h/(\lambda^n-1)$, so
\[
 T\le a_1^2+\rho_2^2h^2\sum_{n\ge2}\frac{\lambda^{n-1}}{(\lambda^n-1)^2}.
\]
We bound this entire series. If $\ell=\log\lambda$, its $n$th term after multiplication by $c^2$ is
\[
 \frac{c^2\lambda^{n-1}}{(\lambda^n-1)^2}
 =\left[\frac{\sinh(\ell/2)}{\sinh(n\ell/2)}\right]^2.
\]
It decreases with $\ell$, because its logarithmic derivative is
$\coth(\ell/2)-n\coth(n\ell/2)<0$; the function $x\coth x$ increases on $(0,\infty)$. The latter follows by differentiating its numerator $\sinh x\cosh x-x$, whose derivative is $2\sinh^2x>0$.
For every $\lambda>1$, the tail with $n\ge5$ satisfies
\[
 c^2\sum_{n\ge5}\frac{\lambda^{n-1}}{(\lambda^n-1)^2}
 \le\frac{c^2}{(1-\lambda^{-5})^2}\sum_{n\ge5}\lambda^{-n-1}
 =\frac{c\lambda^5}{(\lambda^5-1)^2}.
\]
At $\lambda=19/10$, the terms with $n=2,3,4$ and the bound for $n\ge5$ are respectively
\[
 \frac{190}{841}<\frac{23}{100},\quad
 \frac{36100}{423801}<\frac9{100},\quad
 \frac{6859000}{178730161}<\frac1{25},\quad
 \frac{24760990000}{627316273089}<\frac1{25}.
\]
Their right-hand sides sum to $2/5$. Termwise comparison and the limit of nonnegative partial sums therefore prove, for every $\lambda\ge19/10$,
\begin{equation}\label{eq:mh-coarse-V}
 T\le\overline T=s^2\{(1+s)^{-2}+(2/5)\rho_2^2\},\qquad
 v\le V=\frac{c\overline T}{c\overline T+2D_0}.
\end{equation}
The speed bound follows first from~\eqref{eq:mh-coarse-TQ}, then from monotonicity of $cT/(cT+2D_0)$ in $T$.

For $39/20\le\lambda<2$ we have $D_0>\overline T$. Indeed $D_0-\overline T$ is a concave quadratic in $\rho_2$, and hence is at least the smaller of its two endpoint values. At zero its value is $s(1-\lambda s)/(1+s)^2>0$. At one, with $L=c+2+s$, its value is
\[
 s\left\{\frac{1-\lambda s}{(1+s)^2}-\frac25s\right\}
 +\frac{\lambda s^2d}{(1+s)L^2}>0.
\]
The bracket is positive because
\[
 \frac{1-\lambda s}{(1+s)^2}\ge\frac{175}{1156},\qquad
 \frac25s\le\frac{18}{125},\qquad
 \frac{175}{1156}-\frac{18}{125}=\frac{1067}{144500}>0.
\]
The inequality $D_0>\overline T$ is equivalent to $V<c/(\lambda+1)$. For $2\le\lambda<(3+\sqrt5)/2$, the simpler bound $v<h<1/(\lambda+1)$ gives
$c-(\lambda+1)v>c-1\ge0$. We have thus proved $c-(\lambda+1)v>0$ on $39/20\le\lambda<(3+\sqrt5)/2$, as needed for~\eqref{eq:mh-source-Chebyshev}.

\subsubsection{A lower bound on \texorpdfstring{$I$}{I} from the offspring number at the root}
To prove $I>0$, we distinguish whether the root has one child or at least two. Put
\[
 W=(c+X)^{-1},\quad w_0=\E W=(c+h)^{-1},\quad
 H=(c+S)^{-1},\quad z=\E W^2,\quad z_H=\E H^2.
\]
The recursion implies $\lambda W\overset{\mathrm d}=1+H$, so
\[
 \E H=d w_0,\qquad z_H=\lambda^2z-2\lambda w_0+1.
\]
Define the two root expectations
\[
 A_{\rm root}=\E\frac{\sum_iX_i^2}{(c+S)^2},\qquad
 N_{\rm root}=\E\frac{\nu-1}{(c+S)^2}.
\]
Expanding the root expression in~\eqref{eq:mh-signed-integral} gives
\begin{equation}\label{eq:mh-I-root}
 I=2-2\{c+(\lambda+1)v\}d w_0
             +(1+v)A_{\rm root}+2vN_{\rm root}.
\end{equation}
Let $p=\PP(\nu=1)$. Conditional on $\nu=1$, the sum $S$ has the law of $X$. On $\{\nu\ge2\}$ we have $\nu-1\ge1$. Thus
\[
 A_{\rm root}\ge p a_2,\qquad
 N_{\rm root}\ge z_H-pz\ge0,\qquad 0<p\le p_{\max}=z_H/z.
\]
The resulting lower bound in~\eqref{eq:mh-I-root} is affine in $p$. Its two endpoints on $[0,p_{\max}]$ therefore give
\begin{align}\label{eq:mh-F-endpoints}
 I&\ge\min\{F_0(v),F_1(v)\},\notag\\
 F_0(v)&=2-2\{c+(\lambda+1)v\}d w_0+2vz_H,\notag\\
 F_1(v)&=2-2\{c+(\lambda+1)v\}d w_0+(1+v)a_2p_{\max}.
\end{align}
These functions are affine in $v$, and
\[
 F_0(0)=\frac{2\lambda s}{1+s}>0,\qquad F_1(0)=F_0(0)+a_2p_{\max}>0.
\]
Since $0<v\le V$, it remains to check $F_0(V)>0,F_1(V)>0$. No sign assumption on their slopes is necessary.

\begin{lemma}[The sixty-three coefficient inequalities]\label{lem:mh-coarse-arithmetic}
For independent real variables $c\ge19/20$, $0\le s\le9/25$, and $0\le r\le1$, put $\lambda=c+1$, $L=c+2+s$, and define the explicit polynomials
\begin{align*}
 D_N&=s(1-cs)L^2+\lambda r s^2(1+s)L
                         -\lambda^2r^2s^2(1+s)^2,\\
 A_N&=\{1-(2c+1)s\}L+\lambda^2rs(1+s),\\
 B_N&=(1-s)L+rs(1+s),\\
 E_N&=c(2c+3)s-(c+2),\\
 T_N&=s^2L^2\{5+2r^2(1+s)^2\},\\
 s^2P_0&=T_N(A_N+cLE_N)+10\lambda csLD_N,\\
 s^2P_1&=(T_NE_N+10\lambda sD_N)LB_N
                   +(1+s)rsA_N(cT_N+5D_N).
\end{align*}
The numerators in the last two lines are divisible by $s^2$, so $P_0,P_1$ are polynomials, and
\begin{equation}\label{eq:mh-coarse-positive}
 P_0\ge151,\qquad P_1\ge40
\end{equation}
on the entire stated rectangular domain.
\end{lemma}

Here is the exact finite arithmetic content of the lemma. Put $e=c-19/20$. The degrees in $(e,s,r)$ of $P_0,P_1$ are $(6,6,3)$ and $(6,8,4)$. If
\[
 [e^k]P_\ell=\sum_{i,b}a_{kib}s^ir^b,
\]
its degree-$D$ Bernstein coefficients in $r$, where $D=3$ for $\ell=0$ and $D=4$ for $\ell=1$, are
\[
 u_{kj}(s)=\sum_i\left[\sum_{b\le j}a_{kib}
                \frac{\binom jb}{\binom Db}\right]s^i.
\]
There are $7\cdot4+7\cdot5=63$ such polynomials. For each nonzero polynomial, let $q$ be its first nonzero power and write $u_{kj}(s)=\sum_{i\ge q}a_is^i$. Direct rational arithmetic gives $a_q>0$ and
\begin{equation}\label{eq:mh-coarse-payment}
 a_q+\sum_{i>q,\ a_i<0}a_i(9/25)^{i-q}>0.
\end{equation}
For $k=0,\ldots,5$, $q=0$, and the following lower bounds hold for every Bernstein coefficient in the indicated block:
\[
\begin{array}{c|rrrrrr}
 k&0&1&2&3&4&5\\ \hline
 P_0&151&197&173&99&28&3\\
 P_1&40&235&278&133&25&2/5
\end{array}
\]
Four of the polynomials vanish identically. The coefficients of $e^6$ factor as
\[
 [e^6]P_0=4r^2s(1+s)^2,\qquad
 [e^6]P_1=4r^2s(1-s)(1+s)^2.
\]
Inequality~\eqref{eq:mh-coarse-payment} bounds each nonzero polynomial for $0\le s\le9/25$. After factoring out $s^q$, we discard the positive higher-order terms and bound each negative term in absolute value at $s=9/25$. Since the Bernstein basis is nonnegative and sums to one and $e\ge0$, the coefficient inequalities give~\eqref{eq:mh-coarse-positive}.

We now substitute $r=\rho_2$ and $s=h/c$. The moment formulas give
\begin{align*}
 D_0&=\frac{D_N}{(1+s)^2L^2},&
 \overline T&=\frac{T_N}{5(1+s)^2L^2},\\
 z_H&=\frac{A_N}{c^2(1+s)L},&
 z&=\frac{B_N}{c^2(1+s)L},\\
 F_0(V)&=\frac{2s^2P_0}{c(1+s)L(cT_N+10D_N)},&
 F_1(V)&=\frac{2s^2P_1}{(1+s)LB_N(cT_N+10D_N)}.
\end{align*}
At these moment values all denominators are strictly positive. Indeed $c,s,L,B_N>0$ and $D_N>0$ by~\eqref{eq:mh-coarse-TQ}. The polynomial lemma on the larger rectangle only bounds $P_0,P_1$. Lemma~\ref{lem:mh-coarse-arithmetic} now gives $F_0(V),F_1(V)>0$. Equations~\eqref{eq:mh-F-endpoints} and~\eqref{eq:mh-source-Chebyshev} prove $I>0$ and $D>0$.

Combining the integral bound for $39/20\le\lambda<(3+\sqrt5)/2$, the kernel bound~\eqref{eq:mh-middle-conclusion}, and the case $(\lambda+1)h\ge1$ proves $v'(\lambda)<0$ throughout $707/500\le\lambda<(3+\sqrt5)/2$, with $\lambda<m$. Section~\ref{sec:high} gives the Dini derivative bound for $(3+\sqrt5)/2\le\lambda<m$.


\subsection{Completion of the proof and the finite computations}
\label{subsec:completion}
\begin{proof}[Proof of Theorem~\ref{thm:computer}]
Lemma~\ref{lem:full-response} gives $v\in C^1((1,m))$ and
$v'(\lambda)=-D_\lambda/Z$, with $Z>0$. Fix $1<\lambda<m$.
If $\alpha=1$, Lemma~\ref{lem:calibration} gives $D_\lambda>0$.
If $0<\alpha<1$, the moments of $U$ satisfy
\eqref{eq:hard-calibration}. Proposition~\ref{prop:entire-low-region}
covers $1<\lambda<2/\sqrt3$,
Section~\ref{sec:near-region} covers $2/\sqrt3\le\lambda\le707/500$,
and Section~\ref{sec:middle-high} covers
$707/500\le\lambda<(3+\sqrt5)/2$.
Each proof evaluates $D_\lambda$ under the escape law $\pi_\lambda$.
Together they give $D_\lambda>0$, including at the shared endpoints
$2/\sqrt3$, $707/500$ and $39/20$.

For $(3+\sqrt5)/2\le\lambda<m$, Theorem~\ref{thm:high} gives
$\Dplus v(\lambda)<0$. Differentiability, already proved by
Lemma~\ref{lem:full-response}, identifies this Dini derivative with
$v'(\lambda)$. Thus $v'(\lambda)<0$ at every point of $(1,m)$.
For any $1<\lambda_1<\lambda_2<m$, the mean value theorem on
$[\lambda_1,\lambda_2]$ now gives $v(\lambda_2)<v(\lambda_1)$.
Lemma~\ref{lem:continuity-one} gives continuity at $\lambda=1$.
Together with Theorem~\ref{thm:low}, this proves strict decrease
throughout $(0,m)$.
\end{proof}

The accompanying computations~\cite{Supplement} are organized by the
three regional proofs, with a reading guide that maps each finite
arithmetic lemma to its executable verification and exact data.
The programs reconstruct the rational expressions stated here, check
their polynomial identities and coefficient inequalities using exact
integer or rational arithmetic, and provide the resulting coefficient
arrays. Rational interval decompositions are recorded explicitly when
needed. The programs do not infer signs from floating-point samples.

The preceding proofs show that the denominators are positive on the stated domains and that the coefficient signs imply the required inequalities there. They also bound the infinite tails, verify that the parameters of the escape law lie in those domains, and deduce $D_\lambda>0$ from the resulting bounds.
The supplementary historical source records document the provenance of
the arguments. Their file-identity and dependency checks are separate
from the exact arithmetic verification used for this manuscript.



\bibliographystyle{alpha}
\bibliography{references}
\end{document}
